\documentclass[11pt, a4paper]{article}

\usepackage{estilo-mate}
\tema{Teoría de la Medida}

\hypersetup{pdftitle={Apuntes de Teoría de la Medida}}

\title{
    \vspace{-2cm}
    \textbf{\Huge Apuntes de Teoría de la Medida}\\
    \vspace{0.5cm}
}
\author{Lucas Allara}
\date{}

\begin{document}

\maketitle
\newpage

\tableofcontents
\newpage

\section{Motivación: de Riemann a Lebesgue}

Sea $f: [a, b] \to \mathbb{R}$ acotada. La integral de Riemann se construye dividiendo el dominio
$[a, b]$ en una cantidad finita de subintervalos $[t_{k-1}, t_k]$ y aproximando el área bajo la
curva
con sumas
$$ \sum_{k=1}^n f(\xi_k) (t_k - t_{k-1}), \qquad \xi_k \in [t_{k-1}, t_k]. $$
Cuando estas sumas convergen al refinar la partición, se dice que $f$ es Riemann-integrable.

Esta construcción tiene dos limitaciones:
\begin{itemize}
    \item Funciones con muchas discontinuidades no son Riemann-integrables, aunque
    intuitivamente sí deberían tener área. El ejemplo clásico es
    $\mathbf{1}_{\mathbb{Q} \cap [0,1]}$.
    \item Los teoremas de intercambio entre límite e integral requieren hipótesis fuertes.
\end{itemize}

Un intento previo (Jordan, 1892) fue medir un conjunto acotado $E$ aproximándolo desde afuera
por cubrimientos \emph{finitos} de intervalos:
$$ \bar c(E) = \inf\left\{ \sum_{k=1}^n \ell(I_k) : E \subseteq I_1 \cup \cdots \cup I_n \right\}, $$
donde $\ell$ denota la longitud. Este \emph{contenido de Jordan} falla sobre conjuntos densos:
todo cubrimiento finito de $\mathbb{Q} \cap [0,1]$ por intervalos abiertos es un abierto que
contiene un denso de $[0,1]$, y su clausura (que solo agrega finitos extremos) contiene a
$[0,1]$; la suma de longitudes es entonces $\ge 1$, de modo que
$\bar c(\mathbb{Q} \cap [0,1]) = 1$ y, por el mismo argumento,
$\bar c(\mathbb{Q}^c \cap [0,1]) = 1$: dos conjuntos disjuntos cuya unión es $[0,1]$ tendrían
``medidas'' que suman $2$. La corrección decisiva fue admitir cubrimientos \emph{numerables}:
con ellos $\mathbb{Q} \cap [0,1]$ se cubre por intervalos de longitud total arbitrariamente
pequeña (un intervalo de longitud $\epsilon/2^n$ alrededor del $n$-ésimo racional), y su medida
resulta $0$.

\textbf{Idea de Lebesgue (1904).} En vez de partir el dominio, se parte la imagen. Se toma una
partición $a_0 < a_1 < \cdots < a_n$ del rango de $f$, y para cada $k$ se considera el conjunto
$$ A_k = \{ t \in [a, b] : a_{k-1} < f(t) \le a_k \}. $$
La integral se aproxima por
$$ \sum_{k=1}^n a_k \cdot m(A_k), $$
donde $m(A_k)$ es una noción de \emph{medida} del conjunto $A_k$.

El objetivo es construir, partiendo de una clase razonable de subconjuntos llamados
\emph{medibles}, una función $m: \mathcal{A} \to [0, \infty]$ con propiedades como:
\begin{itemize}
    \item $A \subseteq B \implies m(A) \le m(B)$.
    \item $m(A \cup B) = m(A) + m(B)$ cuando $A \cap B = \emptyset$.
    \item En $\mathbb{R}$, $m([a, b]) = b - a$ y $m$ es invariante por traslaciones.
\end{itemize}
Se verá que \emph{no} todo subconjunto de $\mathbb{R}$ admite una tal medida; hay subconjuntos
 que no son medibles.

\newpage

\section{\texorpdfstring{$\sigma$}{σ}-álgebras}

\begin{definicion}{$\sigma$-álgebra}{sigma_algebra}
Sea $X$ un conjunto. Una familia $\mathcal{A} \subseteq \mathcal{P}(X)$ es una $\sigma$-álgebra
sobre $X$ si:
\begin{enumerate}
    \item $\emptyset \in \mathcal{A}$.
    \item Cerrada por complementos: $A \in \mathcal{A} \implies A^c \in \mathcal{A}$.
    \item Cerrada por uniones numerables: $(A_n) \subset \mathcal{A} \implies
    \bigcup_{n=1}^\infty A_n \in \mathcal{A}$.
\end{enumerate}
A los elementos de $\mathcal{A}$ se los llama \emph{conjuntos medibles} (respecto a $\mathcal{A}$).
\end{definicion}

\obs Por De Morgan, una $\sigma$-álgebra también es cerrada por intersecciones numerables: si
$(A_n) \subset \mathcal{A}$, entonces
$$ \bigcap_n A_n = \left( \bigcup_n A_n^c \right)^c \in \mathcal{A}. $$
Además, $X = \emptyset^c \in \mathcal{A}$.

\subsection{Ejemplos}

\begin{ejemplo}{$\sigma$-álgebra trivial}{}
$\mathcal{A} = \{\emptyset, X\}$ es la $\sigma$-álgebra más pequeña posible sobre $X$. Cumple
los tres axiomas por inspección.
\end{ejemplo}

\begin{ejemplo}{$\sigma$-álgebra discreta}{}
$\mathcal{A} = \mathcal{P}(X)$ es la $\sigma$-álgebra más grande, donde todo subconjunto es
medible. Sobre conjuntos pequeños esto es lo natural.
\end{ejemplo}

\begin{ejemplo}{Conjuntos cofinitos: un álgebra que no es $\sigma$-álgebra}{}
Sea $X$ infinito y $\mathcal{A} = \{A \subseteq X : A \text{ finito o } A^c \text{ finito}\}$.
Esta familia es un \emph{álgebra} (cerrada por uniones e intersecciones finitas y por
complementos) pero \emph{no} es $\sigma$-álgebra.
\end{ejemplo}

\begin{ejemplo}{$\sigma$-álgebra de las sucesiones de monedas}{ej_monedas}
Sea $X = \{0, 1\}^{\mathbb{N}}$ el espacio de las sucesiones infinitas de ceros y unos (los
resultados de tirar una moneda infinitas veces, con $1$ = cara). Para cada $n$ se define el evento
$$ E_n = \{ x \in X : x_n = 1 \} \qquad (\text{``la } n\text{-ésima tirada es cara''}), $$
y se toma $\mathcal{A} = \sigma(E_1, E_2, \dots)$. Esta $\sigma$-álgebra contiene todos los
eventos descriptos por una cantidad numerable de tiradas: por ejemplo $\bigcap_n E_n$
(``siempre cara''), $E_1 \cap E_2^c$ (``cara y después ceca''), o los límites de la
subsección siguiente. Este espacio reaparecerá al estudiar medidas y los lemas de
Borel-Cantelli.
\end{ejemplo}

\begin{proposicion}{Intersección de $\sigma$-álgebras}{interseccion_sigmas}
Si $(\mathcal{A}_i)_{i \in I}$ es una familia arbitraria de
$\sigma$-álgebras sobre un conjunto $X$, entonces la intersección
$$ \mathcal{B} = \bigcap_{i \in I} \mathcal{A}_i $$
es nuevamente una $\sigma$-álgebra sobre $X$.
\end{proposicion}

\begin{proof}\leavevmode\\
Se verifican los tres axiomas que definen una $\sigma$-álgebra.

\textit{Contiene al vacío.} Para cada $i \in I$, $\mathcal{A}_i$ es $\sigma$-álgebra y por lo
tanto $\emptyset \in \mathcal{A}_i$. Como esto vale para todo $i$, $\emptyset \in \mathcal{B}$.

\textit{Cerrada por complementos.} Sea $A \in \mathcal{B}$. Para cada $i \in I$, $A \in
\mathcal{A}_i$, y como $\mathcal{A}_i$ es cerrada por complementos, $A^c \in \mathcal{A}_i$. Por lo
tanto $A^c$ pertenece a todas las $\mathcal{A}_i$, esto es, $A^c \in \mathcal{B}$.

\textit{Cerrada por uniones numerables.} Sea $(A_n)_{n \in \mathbb{N}} \subseteq \mathcal{B}$ una
sucesión. Fíjese $i \in I$. Para cada $n$, $A_n \in \mathcal{B} \subseteq \mathcal{A}_i$, de
modo que $(A_n) \subseteq \mathcal{A}_i$. Como $\mathcal{A}_i$ es cerrada por uniones numerables,
$\bigcup_{n=1}^{\infty} A_n \in \mathcal{A}_i$. Esto vale para todo $i \in I$, lo cual implica
$\bigcup_n A_n \in \mathcal{B}$.
\end{proof}

\obs Esta proposición es lo que permite definir la \emph{$\sigma$-álgebra generada} por una
familia: dado un conjunto $\mathcal{F} \subseteq \mathcal{P}(X)$ existe siempre al menos una
$\sigma$-álgebra que contiene a $\mathcal{F}$ (por ejemplo $\mathcal{P}(X)$), y la intersección
de todas ellas es la $\sigma$-álgebra más pequeña con esa propiedad.

\subsection{\texorpdfstring{$\sigma$}{σ}-álgebra generada}

\begin{definicion}{$\sigma$-álgebra generada}{sigma_generada}
Dada una familia $\mathcal{F} \subseteq \mathcal{P}(X)$, la $\sigma$-álgebra generada por
$\mathcal{F}$ es
$$ \sigma(\mathcal{F}) = \bigcap \{ \mathcal{A} : \mathcal{A} \text{ es }\sigma\text{-álgebra y }
\mathcal{F} \subseteq \mathcal{A} \}. $$
Es la $\sigma$-álgebra más pequeña que contiene a $\mathcal{F}$.
\end{definicion}

\subsection{Estructura en conjuntos numerables}

Sobre un conjunto numerable, toda $\sigma$-álgebra proviene de una partición, cuyos elementos
se llaman \emph{átomos}: los menores conjuntos medibles no vacíos que contienen un punto dado.

\begin{teorema}{$\sigma$-álgebras en conjuntos numerables}{atomos_sigma}
Sean $X$ numerable y $\Sigma$ una $\sigma$-álgebra sobre $X$. Para cada $x \in X$ se define
$$ A_x = \bigcap_{\substack{B \in \Sigma \\ x \in B}} B. $$
La familia $\mathcal{A} = \{A_x : x \in X\}$ es una partición de $X$ por elementos de $\Sigma$,
y $\Sigma = \sigma(\mathcal{A})$.
\end{teorema}

\begin{proof}
Cada $B$ en la intersección contiene a $x$, por lo cual $x \in A_x$. En particular los $A_x$
cubren $X$.

Para ver que $A_x \in \Sigma$ se reformula la intersección en términos del complemento, que es
a lo sumo numerable. Para cada $y \in A_x^c$, como $y$ no pertenece a la intersección,
existe algún $B_y \in \Sigma$ con $x \in B_y$ pero $y \notin B_y$. Se define
$$ A_x' = \bigcap_{y \in A_x^c} B_y, $$
que es intersección a lo sumo numerable de elementos de $\Sigma$ y por lo tanto pertenece a
$\Sigma$. Cada $B_y$ figura entre los conjuntos que definen $A_x$, lo cual da $A_x \subseteq B_y$
para todo $y \in A_x^c$, y entonces $A_x \subseteq A_x'$. Recíprocamente, si $y \in A_x^c$
entonces $y \notin B_y$ y por lo tanto $y \notin A_x'$, lo que muestra $A_x^c \subseteq (A_x')^c$,
es decir $A_x' \subseteq A_x$. Se concluye $A_x = A_x' \in \Sigma$.

Véase ahora una propiedad de minimalidad: si $B \in \Sigma$ es no vacío y $B \subseteq A_x$,
entonces $B = A_x$. Supóngase primero que $x \notin B$. El complemento $B^c$ pertenece a $\Sigma$
y contiene a $x$, por lo cual figura entre los conjuntos que definen $A_x$; entonces $A_x \subseteq
B^c$, equivalente a $A_x \cap B = \emptyset$. Pero $B \subseteq A_x$ no vacío fuerza $A_x \cap B
= B \ne \emptyset$, contradicción. Así $x \in B$, y como $B$ es medible que contiene a $x$,
$A_x \subseteq B$. Junto con la hipótesis $B \subseteq A_x$, queda $B = A_x$.

De aquí sale la propiedad de partición. Si $A_x \cap A_y$ es no vacío, esa intersección
está en $\Sigma$, no es vacía y está contenida en $A_x$, por lo cual la minimalidad obliga
$A_x \cap A_y = A_x$. En consecuencia $A_x \subseteq A_y$, e intercambiando $x$ con $y$ se
obtiene $A_y \subseteq A_x$, esto es $A_x = A_y$. Combinado con que los $A_x$ cubren $X$, queda
demostrado que $\mathcal{A}$ es una partición.

Resta ver $\Sigma = \sigma(\mathcal{A})$. La inclusión $\sigma(\mathcal{A}) \subseteq \Sigma$
es consecuencia de que cada $A_x \in \Sigma$. Para la otra inclusión, dado $B \in \Sigma$
se afirma $B = \bigcup_{x \in B} A_x$. La inclusión $\subseteq$ sigue de $x \in A_x$ para cada $x
\in B$. Para la inclusión $\supseteq$, si $x \in B$ entonces $B$ es uno de los conjuntos que
definen $A_x$, lo que da $A_x \subseteq B$. La unión es numerable porque $B \subseteq X$ lo es,
y cada $A_x \in \sigma(\mathcal{A})$, por lo cual $B \in \sigma(\mathcal{A})$.
\end{proof}

\obs Las $\sigma$-álgebras sobre conjuntos numerables son entonces siempre \emph{discretas}: se
generan por una partición.

\begin{definicion}{$\sigma$-álgebra de Borel}{borel}
La $\sigma$-álgebra de Borel en $\mathbb{R}$ es la generada por los abiertos:
$$ \mathcal{B}(\mathbb{R}) = \sigma(\{U \subseteq \mathbb{R} : U \text{ abierto}\}). $$
Sus elementos se llaman \emph{borelianos}.
\end{definicion}

\subsection{Borel a partir de intervalos}

Aunque la $\sigma$-álgebra de Borel se define a partir de \emph{todos} los abiertos, en
$\mathbb{R}$ basta una familia mucho menor: los intervalos abiertos $(a, b)$. Se prueba algo más:
todo abierto de $\mathbb{R}$ se descompone como unión a lo sumo numerable y \emph{disjunta} de
intervalos abiertos. Estos intervalos son las componentes conexas del abierto.

\begin{proposicion}{Estructura de abiertos en $\mathbb{R}$}{abiertos_R}
Todo abierto $A \subseteq \mathbb{R}$ se escribe como unión a lo sumo numerable de intervalos
abiertos disjuntos. En particular,
$$ \mathcal{B}(\mathbb{R}) = \sigma(\{(a, b) : a < b\}). $$
\end{proposicion}

\begin{proof}
La estrategia es asignarle a cada punto $a \in A$ el mayor intervalo abierto contenido en $A$ que
lo contiene. Esos intervalos cubren $A$ y, por su maximalidad, son disjuntos dos a dos. Como cada
uno contiene un racional, la cantidad de intervalos distintos resulta a lo sumo numerable.

Para cada $a \in A$ se define
$$ x(a) = \sup\{x \in \mathbb{R} \setminus A : x < a\}, \qquad y(a) = \inf\{y \in \mathbb{R}
\setminus A : y > a\}, $$
con la convención $\sup \emptyset = -\infty$, $\inf \emptyset = +\infty$, y se define
$I_a = (x(a), y(a))$.

Como $A$ es abierto y $a \in A$, existe $\delta > 0$ con
$(a - \delta, a + \delta) \subseteq A$. Como ningún punto del entorno $(a - \delta, a)$ está en
$\mathbb{R} \setminus A$, todo punto que figura en el conjunto cuyo supremo es $x(a)$ satisface
$x \le a - \delta$. Por lo tanto $x(a) \le a - \delta < a$, e igualmente $y(a) > a$, de modo que
$a \in I_a$.

Para la inclusión $I_a \subseteq A$, sea $t \in I_a$ con $t \ne a$, y supóngase por absurdo
$t \notin A$. Si $t < a$, entonces $t$ pertenece al conjunto $\{x \in \mathbb{R} \setminus A :
x < a\}$ cuyo supremo es $x(a)$, lo cual fuerza $t \le x(a)$ y contradice $t > x(a)$. El caso
$t > a$ es simétrico respecto de $y(a)$.

Como consecuencia, $A = \bigcup_{a \in A} I_a$: la inclusión $\supseteq$ es $I_a \subseteq A$, y
la $\subseteq$ es $a \in I_a$.

Supóngase $I_a \cap I_b \ne \emptyset$ y sea $z \in I_a \cap I_b$. Por definición de
$I_a = (x(a), y(a))$ y $I_b = (x(b), y(b))$, se cumplen las desigualdades
$$ x(a) < z < y(a) \qquad \text{y} \qquad x(b) < z < y(b). $$

Se prueba primero $y(a) = y(b)$ descartando la posibilidad $y(a) < y(b)$ (el caso $y(a) > y(b)$
es simétrico, intercambiando los papeles de $a$ y $b$). En efecto, supóngase $y(a) < y(b)$. Por
definición de $y(a)$ como ínfimo de $\{y \in \mathbb{R} \setminus A : y > a\}$, vale
$y(a) \notin A$. Por las desigualdades $z < y(a) < y(b)$, se tiene $y(a) \in (z, y(b))$. Pero
$(z, y(b)) \subseteq I_b \subseteq A$, de donde $y(a) \in A$, contradiciendo $y(a) \notin A$.
Se concluye $y(a) = y(b)$.

Un argumento análogo del otro lado prueba $x(a) = x(b)$: si fuera $x(a) > x(b)$, por definición
de $x(a)$ como supremo de $\{x \in \mathbb{R} \setminus A : x < a\}$, vale $x(a) \notin A$. Por las
desigualdades $x(b) < x(a) < z$, se tiene $x(a) \in (x(b), z)$. Pero $(x(b), z) \subseteq
I_b \subseteq A$, de donde $x(a) \in A$, contradiciendo $x(a) \notin A$. El caso $x(a) < x(b)$ es
simétrico. Se concluye $x(a) = x(b)$.

Con $x(a) = x(b)$ y $y(a) = y(b)$ se sigue $I_a = I_b$.

Cada $I_a$ es un abierto no vacío de $\mathbb{R}$, de modo que por
densidad de $\mathbb{Q}$ contiene al menos un racional. Eligiendo un racional por intervalo
distinto se obtiene una aplicación inyectiva de la familia de intervalos distintos en
$\mathbb{Q}$ (es inyectiva porque dos intervalos distintos son disjuntos). En consecuencia, la
familia se inyecta en $\mathbb{Q}$ y es a lo sumo numerable.

Sea $\mathcal{I} = \{(a, b)
: a < b\}$. Para una inclusión, cada $(a, b)$ es abierto, lo cual da $\mathcal{I} \subseteq
\mathcal{B}(\mathbb{R})$ y, por minimalidad, $\sigma(\mathcal{I}) \subseteq
\mathcal{B}(\mathbb{R})$.
Para la otra inclusión, todo abierto se escribe como unión numerable de elementos de
$\mathcal{I}$ por lo demostrado, y por lo tanto pertenece a $\sigma(\mathcal{I})$. Así,
$\sigma(\mathcal{I})$ es una $\sigma$-álgebra que contiene a los abiertos, y por minimalidad de
$\mathcal{B}(\mathbb{R})$, $\mathcal{B}(\mathbb{R}) \subseteq \sigma(\mathcal{I})$.
\end{proof}

\begin{corolario}{Conjuntos en $\mathcal{B}(\mathbb{R})$}{borel_contiene}
$\mathcal{B}(\mathbb{R})$ contiene a los abiertos, los cerrados, los puntos, los intervalos
semicerrados y las semirrectas.
\end{corolario}
 
\subsection{Operaciones con sucesiones}

\begin{definicion}{Límite superior e inferior de conjuntos}{limsup_conjuntos}
Dada $(A_n) \subset \mathcal{A}$:
\begin{align*}
\varlimsup_n A_n &= \bigcap_{m=1}^\infty \bigcup_{k \ge m} A_k = \{x : x \in A_n \text{ para
infinitos } n\}, \\
\varliminf_n A_n &= \bigcup_{m=1}^\infty \bigcap_{k \ge m} A_k = \{x : x \in A_n \text{ a partir de
cierto } n\}.
\end{align*}
\end{definicion}

\obs Si $\mathcal{A}$ es $\sigma$-álgebra y $(A_n) \subset \mathcal{A}$, entonces
$\varlimsup A_n, \varliminf A_n \in \mathcal{A}$.

\begin{proposicion}{Propiedades del límite superior e inferior}{props_limsup_liminf}
Sea $(A_n)$ una sucesión de subconjuntos de $X$. Entonces:
\begin{enumerate}
    \item $\varliminf_n A_n \subseteq \varlimsup_n A_n$.
    \item $\left(\varliminf_n A_n\right)^c = \varlimsup_n A_n^c$ y
    $\left(\varlimsup_n A_n\right)^c = \varliminf_n A_n^c$.
    \item Si $A \subseteq A_n \subseteq B$ para todo $n$, entonces
    $A \subseteq \varliminf_n A_n \subseteq \varlimsup_n A_n \subseteq B$.
    \item Si $(A_n)$ es creciente, $\varliminf_n A_n = \varlimsup_n A_n = \bigcup_n A_n$.
    \item Si $(A_n)$ es decreciente, $\varliminf_n A_n = \varlimsup_n A_n = \bigcap_n A_n$.
\end{enumerate}
\end{proposicion}

\begin{proof}\leavevmode\\
(i) Si $x \in \varliminf_n A_n = \bigcup_m \bigcap_{k \ge m} A_k$, existe $m_0$ tal que $x \in
A_k$ para todo $k \ge m_0$; es decir, $x$ está en todos los $A_k$ salvo finitos. En particular,
$x$ está en infinitos $A_k$, por lo cual para todo $m$ existe $k \ge m$ con $x \in A_k$. Esto
significa $x \in \bigcap_m \bigcup_{k \ge m} A_k = \varlimsup_n A_n$.

(ii) Por las leyes de De Morgan aplicadas dos veces,
$$ \left(\varliminf_n A_n\right)^c = \left(\bigcup_m \bigcap_{k \ge m} A_k\right)^c =
\bigcap_m \left(\bigcap_{k \ge m} A_k\right)^c = \bigcap_m \bigcup_{k \ge m} A_k^c =
\varlimsup_n A_n^c. $$
La otra igualdad se prueba análogamente, intercambiando $\bigcup$ y $\bigcap$.

(iii) Si $x \in A$, entonces $x \in A_n$ para todo $n$, en particular $x \in \bigcap_{k \ge m}
A_k$ para todo $m$. Luego $x \in \bigcup_m \bigcap_{k \ge m} A_k = \varliminf_n A_n$, lo cual
da $A \subseteq \varliminf_n A_n$. Combinado con (i), $A \subseteq \varliminf_n A_n \subseteq
\varlimsup_n A_n$.

Para la cota superior, $A_n \subseteq B$ para todo $n$ equivale a $B^c \subseteq A_n^c$ para todo
$n$. Aplicando la inclusión $B^c \subseteq \varliminf_n A_n^c$ (recién demostrada con $B^c$ en
lugar de $A$ y $A_n^c$ en lugar de $A_n$) y luego (ii),
$$ B^c \subseteq \varliminf_n A_n^c = \left(\varlimsup_n A_n\right)^c, $$
de donde, tomando complemento, $\varlimsup_n A_n \subseteq B$.

(iv) Si $(A_n)$ es creciente, $\bigcap_{k \ge m} A_k = A_m$ para todo $m$, de modo que
$$ \varliminf_n A_n = \bigcup_m \bigcap_{k \ge m} A_k = \bigcup_m A_m. $$
Por otro lado, $\bigcup_{k \ge m} A_k = \bigcup_k A_k$ para todo $m$ (la unión no cambia al
comenzar más adelante en una sucesión creciente), de modo que
$$ \varlimsup_n A_n = \bigcap_m \bigcup_{k \ge m} A_k = \bigcup_n A_n. $$

(v) Si $(A_n)$ es decreciente, $(A_n^c)$ es creciente. Aplicando (iv) a $(A_n^c)$ y luego (ii),
$$ \left(\varlimsup_n A_n\right)^c = \varliminf_n A_n^c = \bigcup_n A_n^c =
\left(\bigcap_n A_n\right)^c, $$
y análogamente para $\varliminf$. Tomando complementos, $\varliminf_n A_n = \varlimsup_n A_n
= \bigcap_n A_n$.
\end{proof}

\begin{ejemplo}{Cálculo de límites de conjuntos}{ej_limites_conjuntos}
\begin{enumerate}
    \item Sea $A_n = [0, 2]$ si $n$ es impar y $A_n = [1, 4]$ si $n$ es par. Un punto está en
    infinitos $A_n$ si y solo si pertenece a $[0,2] \cup [1,4] = [0,4]$, y está en todos salvo
    finitos si y solo si pertenece a $[0,2] \cap [1,4] = [1,2]$. Así
    $$ \varliminf_n A_n = [1, 2], \qquad \varlimsup_n A_n = [0, 4], $$
    y la sucesión no converge como conjuntos.
    \item Sea $A_n = [1/n, 1)$. La sucesión es creciente, y por la Proposición
    \ref{prop:props_limsup_liminf},
    $$ \varliminf_n A_n = \varlimsup_n A_n = \bigcup_n [1/n, 1) = (0, 1). $$
    La sucesión converge al conjunto $(0,1)$; obsérvese que el límite no contiene al $0$
    aunque cada $A_n$ esté ``pegado'' a él.
\end{enumerate}
\end{ejemplo}

\obs En el espacio $X = \{0,1\}^{\mathbb{N}}$ de tiradas de moneda con los eventos
$E_n = \{x : x_n = 1\}$, los límites tienen lectura probabilística directa:
$$ \varlimsup_n E_n = \{\text{sale cara infinitas veces}\}, \qquad
\varliminf_n E_n = \{\text{a partir de algún momento, siempre cara}\}. $$
La inclusión $\varliminf \subseteq \varlimsup$ es, con esta lectura, transparente: si a partir
de un momento siempre sale cara, en particular sale cara infinitas veces.

\newpage

\section{Conjuntos no medibles: el ejemplo de Vitali}

Antes de pasar a la definición de medida y a la construcción de Lebesgue, conviene entender
por qué no puede esperarse que \emph{todo} subconjunto de $\mathbb{R}$ admita una medida con las
propiedades buscadas (invariancia por traslaciones y $\sigma$-aditividad). El ejemplo
clásico que muestra que esto es imposible es la construcción de \emph{Vitali}: un subconjunto
$V \subseteq [0, 1]$ que no puede ser medible respecto de ninguna medida invariante por
traslaciones que extienda al ``largo'' $m([0, 1]) = 1$.

\begin{teorema}{Existencia de conjuntos no medibles}{vitali}
No existe una función $m: \mathcal{P}(\mathbb{R}) \to [0, \infty]$ que cumpla simultáneamente:
\begin{enumerate}
    \item $m$ es $\sigma$-aditiva: si $(A_n)$ es disjunta, $m(\bigcup_n A_n) = \sum_n m(A_n)$.
    \item $m$ es invariante por traslaciones: $m(A + t) = m(A)$ para todo $A \subseteq \mathbb{R}$
    y todo $t \in \mathbb{R}$.
    \item $m([0, 1]) = 1$.
\end{enumerate}
\end{teorema}

\begin{proof}\leavevmode\\
Asumiendo el axioma de elección, se construye un $V \subseteq [0, 1]$ que llevará a una
contradicción bajo (i)--(iii).

\textit{Construcción.} Sobre $[0, 1]$ se define la relación de equivalencia
$$ x \sim y \ \Leftrightarrow \ x - y \in \mathbb{Q}. $$
Cada clase $[x]$ es densa en $[0, 1]$ (porque $\mathbb{Q}$ es denso). Por el axioma de elección,
se elige un representante por cada clase y se forma
$$ V = \{\text{un representante de cada clase}\} \subseteq [0, 1]. $$

\textit{Trasladados.} Enumérese $\mathbb{Q} \cap [-1, 1] = \{r_1, r_2, \dots\}$, que es numerable.
Para cada $k$ se define
$$ V_k = V + r_k = \{v + r_k : v \in V\}. $$

\textit{Los $V_k$ son disjuntos dos a dos.} Si $V_k \cap V_j \ne \emptyset$, existen $v_1, v_2 \in
V$ con $v_1 + r_k = v_2 + r_j$, de donde $v_1 - v_2 = r_j - r_k \in \mathbb{Q}$, es decir
$v_1 \sim v_2$. Como $V$ contiene un único representante por clase, $v_1 = v_2$, y entonces
$r_k = r_j$, es decir $k = j$.

\textit{$[0, 1] \subseteq \bigcup_k V_k \subseteq [-1, 2]$.} La inclusión de la derecha es
consecuencia de que $V \subseteq [0, 1]$ y $r_k \in [-1, 1]$ dan $V_k \subseteq [-1, 2]$. Para la
inclusión de la izquierda, dado $x \in [0, 1]$, su clase $[x]$ tiene un representante $v \in V$,
y $x - v \in \mathbb{Q} \cap [-1, 1]$ (porque $x, v \in [0, 1]$), de modo que $x - v = r_k$ para
algún $k$, es decir $x = v + r_k \in V_k$.

\textit{Contradicción.} Supóngase que $V$ es medible. Por (ii), $m(V_k) = m(V)$ para todo $k$,
y por (i),
$$ m\!\left(\bigcup_k V_k\right) = \sum_{k=1}^\infty m(V_k) = \sum_{k=1}^\infty m(V). $$
Se discute según el valor de $m(V) \in [0, \infty]$:
\begin{itemize}
    \item Si $m(V) = 0$, la suma da $0$, y por monotonía $m([0, 1]) \le m(\bigcup_k V_k) = 0$.
    Pero $m([0, 1]) = 1$ por (iii), contradicción.
    \item Si $m(V) > 0$, $\sum_k m(V) = +\infty$, mientras que por monotonía $m(\bigcup_k V_k)
    \le m([-1, 2]) = 3$ (los últimos por traslación de $[0, 1]$ y aditividad: $m([-1, 2]) =
    m([-1, 0]) + m([0, 1]) + m([1, 2]) = 3$). Contradicción.
\end{itemize}
En cualquier caso se llega a un absurdo. Por lo tanto $V$ no puede ser medible respecto de $m$, y
en particular no existe una función $m: \mathcal{P}(\mathbb{R}) \to [0, \infty]$ que cumpla
(i)--(iii) en toda $\mathcal{P}(\mathbb{R})$.
\end{proof}

\obs No se puede definir medida invariante por traslaciones, $\sigma$-aditiva y
normalizada sobre \emph{toda} la $\sigma$-álgebra $\mathcal{P}(\mathbb{R})$. Lo máximo que se
consigue (con axioma de elección) es una medida sobre una $\sigma$-álgebra estrictamente
más pequeña: los borelianos, o algo más: los Lebesgue-medibles.

\newpage

\section{Medidas}

\begin{definicion}{Medida}{def_medida}
Sea $(X, \mathcal{A})$ un par con $\mathcal{A}$ una $\sigma$-álgebra sobre $X$. Una función
$\mu: \mathcal{A} \to [0, \infty]$ es una \emph{medida} si:
\begin{enumerate}
    \item $\mu(\emptyset) = 0$.
    \item ($\sigma$-aditividad) Si $(A_n) \subset \mathcal{A}$ es una familia disjunta dos a dos,
    $$ \mu\!\left(\bigcup_{n=1}^\infty A_n\right) = \sum_{n=1}^\infty \mu(A_n). $$
\end{enumerate}
A la terna $(X, \mathcal{A}, \mu)$ se la llama \emph{espacio de medida}.

Si $\mu(X) < \infty$, se dice que $\mu$ es \emph{finita}. Si $X = \bigcup_n X_n$ con
$\mu(X_n) < \infty$ para todo $n$, se dice que $\mu$ es \emph{$\sigma$-finita}. Si $\mu(X) = 1$,
$\mu$ es una \emph{probabilidad}.
\end{definicion}

\subsection{Ejemplos}

\begin{proposicion}{Delta de Dirac}{delta_dirac}
Sea $X$ un conjunto no vacío y $x_0 \in X$ un punto fijo. La función $\delta_{x_0}:
\mathcal{P}(X) \to [0, \infty]$ definida por
$$ \delta_{x_0}(A) = \begin{cases} 1, & x_0 \in A, \\ 0, & x_0 \notin A, \end{cases} $$
es una medida (llamada \emph{delta de Dirac} concentrada en $x_0$).
\end{proposicion}

\begin{proof}\leavevmode\\
\textit{$\delta_{x_0}(\emptyset) = 0$.} Como $x_0 \notin \emptyset$, $\delta_{x_0}(\emptyset) = 0$.

\textit{$\sigma$-aditividad.} Sea $(A_n)_{n \in \mathbb{N}} \subset \mathcal{P}(X)$ una familia
disjunta dos a dos, y $A = \bigcup_{n} A_n$. Se separa en dos casos según esté o no $x_0$ en
$A$.

\textit{Caso 1: $x_0 \in A$.} Existe $n_0 \in \mathbb{N}$ con $x_0 \in A_{n_0}$. Como los
$A_n$ son disjuntos dos a dos, ese $n_0$ es único: para todo $n \ne n_0$, $x_0 \notin A_n$.
Por lo tanto
$$ \sum_{n=1}^{\infty} \delta_{x_0}(A_n) = \delta_{x_0}(A_{n_0}) + \sum_{n \ne n_0}
\delta_{x_0}(A_n) = 1 + 0 = 1 = \delta_{x_0}(A). $$

\textit{Caso 2: $x_0 \notin A$.} Entonces $x_0 \notin A_n$ para todo $n$, por lo cual
$\delta_{x_0}(A_n) = 0$ para todo $n$, y
$$ \sum_{n=1}^{\infty} \delta_{x_0}(A_n) = 0 = \delta_{x_0}(A). $$
\end{proof}

\obs $\delta_{x_0}$ es una probabilidad ($\delta_{x_0}(X) = 1$). Concentra toda la masa en el
punto $x_0$.

\begin{proposicion}{Medida de conteo}{medida_conteo}
Sea $X$ un conjunto no vacío. La función $\mu: \mathcal{P}(X) \to [0, \infty]$ definida por
$$ \mu(A) = \#A = \begin{cases} \text{cantidad de elementos de } A, & \text{si } A \text{ finito,}
\\ \infty, & \text{si } A \text{ infinito,} \end{cases} $$
es una medida (llamada \emph{medida de conteo}).
\end{proposicion}

\begin{proof}\leavevmode\\
\textit{$\mu(\emptyset) = 0$.} El conjunto vacío tiene cero elementos.

\textit{$\sigma$-aditividad.} Sea $(A_n)_{n \in \mathbb{N}}$ disjunta dos a dos y
$A = \bigcup_n A_n$. La estructura de la demostración depende de si $A$ es finito o infinito.

\textit{Caso 1: $A$ finito.} Como cada $A_n \subseteq A$, todos los $A_n$ son finitos. Más
aún, dado que los $A_n$ son disjuntos y $A$ es finito, solo finitos de ellos pueden ser no
vacíos; en otro caso $A$ contendría al menos un elemento por cada $A_n$ no vacío,
contradiciendo $\#A < \infty$. Renumerando, supóngase $A_1, \dots, A_N$ no vacíos y $A_n =
\emptyset$ para $n > N$. Por conteo elemental de uniones disjuntas finitas,
$$ \#A = \#(A_1 \sqcup \cdots \sqcup A_N) = \sum_{n=1}^{N} \#A_n = \sum_{n=1}^{\infty} \#A_n, $$
donde la última igualdad usa que $\#A_n = 0$ para $n > N$.

\textit{Caso 2: $A$ infinito.} Entonces $\mu(A) = \infty$. Se distinguen dos subcasos:
\begin{itemize}
    \item Si existe $n_0$ con $A_{n_0}$ infinito, entonces $\#A_{n_0} = \infty$, y
    $$ \sum_{n=1}^{\infty} \#A_n \ge \#A_{n_0} = \infty. $$
    \item Si todos los $A_n$ son finitos, como $A$ es infinito debe haber infinitos $A_n$ no
    vacíos; en otro caso $A$ sería unión finita de finitos, y por lo tanto finito. La serie
    de términos positivos
    $$ \sum_{n=1}^{\infty} \#A_n $$
    tiene entonces infinitos términos $\ge 1$, y diverge a $\infty$.
\end{itemize}
En ambos subcasos $\sum_n \#A_n = \infty = \#A$.
\end{proof}

\obs La medida de conteo es finita si y solo si $X$ es finito, y es $\sigma$-finita si y solo si
$X$ es a lo sumo numerable. Sobre $X = \mathbb{N}$ con la medida de conteo, las integrales
respecto de $\mu$ se reducen a series.

\begin{ejemplo}{Medida de Bernoulli}{ej_medida_bernoulli}
Sobre $X = \{0,1\}^{\mathbb{N}}$ con la $\sigma$-álgebra generada por los eventos
$E_n = \{x : x_n = 1\}$ (Ejemplo \ref{ej:ej_monedas}), existe una única probabilidad $P$ que
modela tiradas independientes de una moneda equilibrada: sobre los eventos que fijan $k$
coordenadas,
$$ P(\{x : x_{i_1} = \epsilon_1, \dots, x_{i_k} = \epsilon_k\}) = \frac{1}{2^k}. $$
Con ella puede calcularse la probabilidad de \emph{alguna vez cara}. El evento es
$\bigcup_n E_n$, que no es unión disjunta; se disjuntiza con
\begin{align*}
A_n = E_n \setminus (E_1 \cup \cdots \cup E_{n-1})
&= \{x : x_1 = \cdots = x_{n-1} = 0, \ x_n = 1\} \\
&= \{\text{la primera cara es la } n\text{-ésima}\},
\end{align*}
que fija $n$ coordenadas, de modo que $P(A_n) = 1/2^n$. Por $\sigma$-aditividad,
$$ P\!\left(\bigcup_n E_n\right) = \sum_{n=1}^\infty P(A_n) = \sum_{n=1}^\infty \frac{1}{2^n}
= 1. $$
Con probabilidad $1$ sale alguna cara.
\end{ejemplo}

\obs De los ejemplos anteriores, solo la delta de Dirac y la medida de conteo quedaron
\emph{construidas}. Para la medida de Bernoulli, y para la medida de Lebesgue en $\mathbb{R}$
(que asigna a cada intervalo su longitud), hasta aquí no se demostró la existencia: ese es el
objetivo de la sección de generación de medidas.

\subsection{Propiedades básicas}

\begin{teorema}{Propiedades de una medida}{prop_medida}
Sea $(X, \mathcal{A}, \mu)$ un espacio de medida. Para $A, B, A_n \in \mathcal{A}$:
\begin{enumerate}
    \item Si $E \subseteq F$, $\mu(E) \le \mu(F)$.
    \item Si $E \subseteq F$ y $\mu(E) < \infty$, $\mu(F \setminus E) = \mu(F) - \mu(E)$.
    \item Si $\mu(A \cap B) < \infty$,
    $$ \mu(A \cup B) = \mu(A) + \mu(B) - \mu(A \cap B). $$
    \item Para cualquier familia (no necesariamente disjunta) $(A_n)$,
    $$ \mu\!\left(\bigcup_n A_n\right) \le \sum_n \mu(A_n). $$
    \item Si $A_n \nearrow A$ (es decir, $A_n \subseteq A_{n+1}$ y $A = \bigcup_n A_n$),
    $$ \mu(A) = \lim_{n \to \infty} \mu(A_n). $$
    \item Si $A_n \searrow A$ y $\mu(A_1) < \infty$,
    $$ \mu(A) = \lim_{n \to \infty} \mu(A_n). $$
    \item Si $\varlimsup_n A_n = \varliminf_n A_n = A$ (es decir, $(A_n)$ converge a $A$ como
    conjuntos) y $\mu\!\left(\bigcup_n A_n\right) < \infty$,
    $$ \mu(A) = \lim_{n \to \infty} \mu(A_n). $$
\end{enumerate}
\end{teorema}

\begin{proof}\leavevmode\\
(i) $F = E \cup (F \setminus E)$ es unión disjunta, de modo que
$\mu(F) = \mu(E) + \mu(F \setminus E)
\ge \mu(E)$.

(ii) Inmediato del paso anterior despejando.

(iii) Se descompone $A \cup B$ en tres conjuntos disjuntos: $A \cup B = (A \setminus B) \sqcup
(B \setminus A) \sqcup (A \cap B)$. Por $\sigma$-aditividad,
$$ \mu(A \cup B) = \mu(A \setminus B) + \mu(B \setminus A) + \mu(A \cap B). $$
Análogamente, $A = (A \setminus B) \sqcup (A \cap B)$ y $B = (B \setminus A) \sqcup
(A \cap B)$, de donde
$$ \mu(A) + \mu(B) = \mu(A \setminus B) + \mu(B \setminus A) + 2 \mu(A \cap B). $$
Restando $\mu(A \cap B)$ (finita por hipótesis, lo que evita indeterminaciones $\infty - \infty$)
se obtiene $\mu(A) + \mu(B) - \mu(A \cap B) = \mu(A \cup B)$.

(iv) Se disjuntiza: $B_1 = A_1$, $B_n = A_n \setminus (A_1 \cup \cdots \cup A_{n-1})$. Los $B_n$ son
disjuntos, $B_n \subseteq A_n$ (así $\mu(B_n) \le \mu(A_n)$), y $\bigcup B_n = \bigcup A_n$. Luego
$$ \mu\!\Big(\!\bigcup A_n\Big) = \sum \mu(B_n) \le \sum \mu(A_n). $$

(v) Si $A_n \nearrow A$, se disjuntiza $B_1 = A_1$, $B_n = A_n \setminus A_{n-1}$.
Entonces $A_n = \bigsqcup_{k=1}^n B_k$ y $A = \bigsqcup_k B_k$, de modo que
$$ \mu(A) = \sum_{k=1}^\infty \mu(B_k) = \lim_n \sum_{k=1}^n \mu(B_k) = \lim_n \mu(A_n). $$

(vi) Si $A_n \searrow A$, se define $C_n = A_1 \setminus A_n$, que crece a $A_1 \setminus A$.
Aplicando (v) y (ii) (usando $\mu(A_1) < \infty$),
$$ \mu(A_1) - \mu(A) = \mu(A_1 \setminus A) = \lim_n \mu(C_n) = \lim_n (\mu(A_1) - \mu(A_n)) =
\mu(A_1) - \lim_n \mu(A_n). $$
Cancelando $\mu(A_1)$, $\mu(A) = \lim \mu(A_n)$.

(vii) Se define
$$ B_m = \bigcap_{k \ge m} A_k, \qquad C_m = \bigcup_{k \ge m} A_k. $$
Por construcción, $B_m \nearrow \varliminf_n A_n = A$ y $C_m \searrow \varlimsup_n A_n = A$.
Además, $B_m \subseteq A_m \subseteq C_m$ por monotonía de la intersección y la unión.

Por (v) aplicada a $(B_m)$,
$$ \lim_m \mu(B_m) = \mu(A). $$
Por (vi) aplicada a $(C_m)$ (que es lícita porque $\mu(C_1) \le \mu(\bigcup_n A_n) < \infty$),
$$ \lim_m \mu(C_m) = \mu(A). $$
Por (i), $\mu(B_m) \le \mu(A_m) \le \mu(C_m)$ para todo $m$. Por el teorema de
intercalación, $\lim_m \mu(A_m) = \mu(A)$.
\end{proof}

\subsection{Lemas de Borel-Cantelli}

Recuérdese la interpretación probabilística de $\varlimsup A_n$: un punto $\omega$ pertenece
a $\varlimsup A_n = \bigcap_m \bigcup_{k \ge m} A_k$ si y solo si está en infinitos $A_n$. Por
eso suele escribirse $\{A_n \text{ infinitas veces}\}$. Los lemas de Borel-Cantelli dan
condiciones bajo las cuales este evento tiene medida cero o uno.

\begin{teorema}{Borel-Cantelli (I)}{borel_cantelli_I}
Sea $(X, \mathcal{A}, \mu)$ un espacio de medida y $(A_n) \subset \mathcal{A}$. Si
$$ \sum_{n=1}^\infty \mu(A_n) < \infty, $$
entonces $\mu(\varlimsup_n A_n) = 0$.
\end{teorema}

\begin{proof}
Por definición, $\varlimsup_n A_n = \bigcap_{m=1}^\infty C_m$, donde $C_m = \bigcup_{k \ge m}
A_k$. Por sub-$\sigma$-aditividad,
$$ \mu(C_m) = \mu\!\left(\bigcup_{k \ge m} A_k\right) \le \sum_{k=m}^\infty \mu(A_k). $$
Como $\sum_n \mu(A_n) < \infty$, su cola $\sum_{k \ge m} \mu(A_k)$ tiende a $0$ cuando $m \to
\infty$. En particular, $\mu(C_1) \le \sum_{k=1}^\infty \mu(A_k) < \infty$, por lo cual la
sucesión decreciente $(C_m)$ satisface las hipótesis de continuidad por sucesiones
decrecientes. Entonces
$$ \mu\!\left(\varlimsup_n A_n\right) = \mu\!\left(\bigcap_m C_m\right) = \lim_{m \to \infty}
\mu(C_m) \le \lim_{m \to \infty} \sum_{k=m}^\infty \mu(A_k) = 0. $$
Por monotonía, $\mu(\varlimsup A_n) = 0$.
\end{proof}

\obs Borel-Cantelli (I) no requiere independencia de los $A_n$ ni que $\mu$ sea una probabilidad:
vale para toda medida.

\begin{definicion}{Independencia de eventos}{independencia_eventos}
Sea $(\Omega, \mathcal{F}, \mathbb{P})$ un espacio de probabilidad.
\begin{enumerate}
    \item Dos eventos $A, B \in \mathcal{F}$ son \emph{independientes} si
    $\mathbb{P}(A \cap B) = \mathbb{P}(A) \mathbb{P}(B)$.
    \item Una familia finita $A_1, \dots, A_n \in \mathcal{F}$ es \emph{independiente} si para todo
    subconjunto de índices $\{i_1, \dots, i_k\} \subseteq \{1, \dots, n\}$ (con $k \ge 1$),
    $$ \mathbb{P}(A_{i_1} \cap \cdots \cap A_{i_k}) = \mathbb{P}(A_{i_1}) \cdots
    \mathbb{P}(A_{i_k}). $$
    \item Una familia $(A_n)_{n \in \mathbb{N}}$ es \emph{independiente} si toda subfamilia finita
    suya lo es.
\end{enumerate}
\end{definicion}

\obs Si $(A_n)$ son independientes, también lo son los complementos $(A_n^c)$ (en cualquier
subfamilia finita). Esto sale por inclusión-exclusión: por ejemplo, $\mathbb{P}(A_1^c \cap
A_2^c) = 1 - \mathbb{P}(A_1) - \mathbb{P}(A_2) + \mathbb{P}(A_1) \mathbb{P}(A_2) = (1 -
\mathbb{P}(A_1))(1 - \mathbb{P}(A_2)) = \mathbb{P}(A_1^c) \mathbb{P}(A_2^c)$; el caso general por
inducción.

\begin{teorema}{Borel-Cantelli (II)}{borel_cantelli_II}
Sea $(\Omega, \mathcal{F}, \mathbb{P})$ un espacio de probabilidad y $(A_n) \subset \mathcal{F}$ una
sucesión \emph{independiente}. Si
$$ \sum_{n=1}^\infty \mathbb{P}(A_n) = \infty, $$
entonces $\mathbb{P}(\varlimsup_n A_n) = 1$.
\end{teorema}

\begin{proof}
Bastará probar $\mathbb{P}((\varlimsup_n A_n)^c) = 0$. Por las propiedades vistas anteriormente,
$$ (\varlimsup_n A_n)^c = \varliminf_n A_n^c. $$
Por sub-$\sigma$-aditividad, alcanza con ver que para todo $m$,
$$ \mathbb{P}\!\left(\bigcap_{k \ge m} A_k^c\right) = 0. $$
Fíjese $m$. Por continuidad por sucesiones decrecientes (la intersección finita $\bigcap_{k =
m}^M A_k^c$ decrece a $\bigcap_{k \ge m} A_k^c$ y tiene probabilidad $\le 1 < \infty$),
$$ \mathbb{P}\!\left(\bigcap_{k \ge m} A_k^c\right) = \lim_{M \to \infty} \mathbb{P}\!\left(
\bigcap_{k = m}^M A_k^c\right). $$
Como $(A_n)$ es independiente, $(A_n^c)$ también lo es. Aplicando la
factorización a la familia finita $A_m^c, \dots, A_M^c$,
$$ \mathbb{P}\!\left(\bigcap_{k = m}^M A_k^c\right) = \prod_{k = m}^M \mathbb{P}(A_k^c) =
\prod_{k = m}^M (1 - \mathbb{P}(A_k)). $$
Se usa la desigualdad $1 - x \le e^{-x}$, válida para todo $x \in \mathbb{R}$:
$$ \prod_{k = m}^M (1 - \mathbb{P}(A_k)) \le \prod_{k = m}^M e^{-\mathbb{P}(A_k)} =
\exp\!\left(-\sum_{k = m}^M \mathbb{P}(A_k)\right). $$
Como $\sum_n \mathbb{P}(A_n) = \infty$, la cola $\sum_{k = m}^M \mathbb{P}(A_k) \to \infty$ cuando
$M \to \infty$ (para cada $m$ fijo). Por lo tanto el exponente tiende a $-\infty$ y
$$ \lim_{M \to \infty} \mathbb{P}\!\left(\bigcap_{k = m}^M A_k^c\right) \le \lim_{M \to \infty}
\exp\!\left(-\sum_{k = m}^M \mathbb{P}(A_k)\right) = 0, $$
es decir, $\mathbb{P}(\bigcap_{k \ge m} A_k^c) = 0$. Como esto vale para todo $m$,
$$ \mathbb{P}\!\left((\varlimsup A_n)^c\right) = \mathbb{P}\!\left(\bigcup_m \bigcap_{k \ge m}
A_k^c\right) \le \sum_m \mathbb{P}\!\left(\bigcap_{k \ge m} A_k^c\right) = 0. $$
Se concluye $\mathbb{P}(\varlimsup A_n) = 1$.
\end{proof}

\begin{ejemplo}{Monedas cada vez más sesgadas}{ej_monedas_sesgadas}
Se tiran monedas independientes, donde la $n$-ésima sale cara con probabilidad $p_n$, y sea
$E_n$ = ``la $n$-ésima sale cara''.
\begin{enumerate}
    \item Si $p_n = 1/n$: como $\sum 1/n = \infty$ y los eventos son independientes,
    Borel-Cantelli (II) da que con probabilidad $1$ salen \emph{infinitas} caras, aunque
    $p_n \to 0$.
    \item Si $p_n = (99/100)^n$: la serie $\sum (99/100)^n$ es geométrica convergente, y
    Borel-Cantelli (I) da que con probabilidad $1$ salen solo \emph{finitas} caras, aunque las
    primeras monedas salgan cara casi siempre.
\end{enumerate}
\end{ejemplo}

\begin{ejemplo}{El mono y Borges}{ej_mono}
Un mono teclea al azar, con teclas equiprobables e independientes, en un teclado de $60$
símbolos. Fíjese una obra de $N$ caracteres (digamos, un cuento de Borges). Se afirma: con
probabilidad $1$, el mono la escribe, y no una sino \emph{infinitas} veces.

Considérense los bloques disjuntos de $N$ teclas consecutivas (el bloque $k$-ésimo ocupa las
posiciones $(k-1)N + 1, \dots, kN$) y sea $S_k$ el evento ``el bloque $k$-ésimo coincide
carácter a carácter con la obra''. Los $S_k$ dependen de teclas disjuntas, de modo que son
independientes, y cada uno tiene probabilidad $p = (1/60)^N$: minúscula, pero positiva y la
misma para todo $k$. Entonces
$$ \sum_{k=1}^\infty \mathbb{P}(S_k) = \sum_{k=1}^\infty \frac{1}{60^N} = \infty, $$
y Borel-Cantelli (II) da $\mathbb{P}(\varlimsup_k S_k) = 1$: infinitos bloques reproducen la
obra completa.
\end{ejemplo}

\obs El ejemplo ilustra el uso típico de Borel-Cantelli (II): para concluir ``infinitas
veces'' no hace falta que los eventos sean probables; alcanza con que sean independientes y que
la serie de sus probabilidades diverja, cosa que una probabilidad constante positiva garantiza.

\obs Combinando ambas direcciones: si $(A_n)$ son \emph{independientes}, entonces
$$ \mathbb{P}(\varlimsup A_n) = \begin{cases} 0, & \text{si } \sum \mathbb{P}(A_n) < \infty, \\
1, & \text{si } \sum \mathbb{P}(A_n) = \infty. \end{cases} $$
Es decir, bajo independencia el evento $\{A_n \text{ infinitas veces}\}$ es siempre de
probabilidad $0$ o $1$. Esta dicotomía es un caso particular de un fenómeno mucho más
general: la ley $0$-$1$ de Kolmogorov.

\subsection{Ley 0-1 de Kolmogorov}

\begin{definicion}{$\sigma$-álgebra de cola}{sigma_cola}
Sea $(\Omega, \mathcal{F}, \mathbb{P})$ un espacio de probabilidad y $(A_n)_{n \in \mathbb{N}}
\subset \mathcal{F}$. La \emph{$\sigma$-álgebra de cola} de la sucesión $(A_n)$ es
$$ \mathcal{T} = \bigcap_{n=1}^\infty \sigma(A_n, A_{n+1}, A_{n+2}, \dots), $$
donde $\sigma(A_n, A_{n+1}, \dots)$ es la $\sigma$-álgebra generada por la sucesión a partir
del índice $n$. Sus elementos se llaman \emph{eventos de cola}: son los que no dependen de
ningún número finito de los $A_k$.
\end{definicion}

\obs Ejemplos de eventos de cola: $\varlimsup_n A_n$ y $\varliminf_n A_n$ (ambos están en
$\sigma(A_n, A_{n+1}, \dots)$ para todo $n$, por su misma definición). El evento ``la sucesión
$\mathbf{1}_{A_n}$ converge'' también es de cola. En cambio, $A_1$ por sí solo no lo es.

\begin{teorema}{Ley 0-1 de Kolmogorov}{kolmogorov_01}
Sea $(\Omega, \mathcal{F}, \mathbb{P})$ un espacio de probabilidad y $(A_n)_{n \in \mathbb{N}}
\subset \mathcal{F}$ una sucesión \emph{independiente}. Si $B \in \mathcal{T}$ es un evento de
cola, entonces
$$ \mathbb{P}(B) \in \{0, 1\}. $$
\end{teorema}

\begin{proof}
La estrategia es probar que $B$ es independiente de sí mismo, lo cual fuerza $\mathbb{P}(B) =
\mathbb{P}(B \cap B) = \mathbb{P}(B)^2$ y por lo tanto $\mathbb{P}(B) \in \{0, 1\}$.

Se usa el siguiente criterio de independencia de $\sigma$-álgebras generadas, consecuencia del
teorema $\pi$-$\lambda$ de Dynkin, cuya demostración se omite: \emph{si dos familias de eventos
$\mathcal{C}_1, \mathcal{C}_2$ son cerradas por intersecciones finitas y todo par $C_1 \in
\mathcal{C}_1$, $C_2 \in \mathcal{C}_2$ es independiente, entonces $\sigma(\mathcal{C}_1)$ y
$\sigma(\mathcal{C}_2)$ son independientes} (todo par de elementos, uno de cada una, lo es).

\textit{$\mathcal{T}$ es independiente de $\sigma(A_1, \dots, A_n)$, para cada $n$.} Sean
$\mathcal{C}_1$ la familia de las intersecciones finitas de conjuntos $A_i$ con $i \le n$, y
$\mathcal{C}_2$ la de las intersecciones finitas de conjuntos $A_j$ con $j > n$. Ambas son
cerradas por intersecciones finitas, y todo par de elementos de una y otra es independiente
porque la sucesión $(A_n)$ lo es (los índices involucrados son todos distintos). Por el
criterio, $\sigma(A_1, \dots, A_n)$ y $\sigma(A_{n+1}, A_{n+2}, \dots)$ son independientes; como
$\mathcal{T} \subseteq \sigma(A_{n+1}, A_{n+2}, \dots)$, en particular $\mathcal{T}$ es
independiente de $\sigma(A_1, \dots, A_n)$.

\textit{$\mathcal{T}$ es independiente de $\sigma(A_1, A_2, \dots)$.} Sea $\mathcal{C} =
\bigcup_n \sigma(A_1, \dots, A_n)$. Esta familia es cerrada por intersecciones finitas (dos
elementos suyos conviven en algún $\sigma(A_1, \dots, A_n)$ con $n$ suficientemente grande),
todo elemento suyo es independiente de todo elemento de $\mathcal{T}$ por el paso anterior, y
$\sigma(\mathcal{C}) = \sigma(A_1, A_2, \dots)$. Aplicando el criterio otra vez (con
$\mathcal{C}$ y $\mathcal{T}$, que es cerrada por intersecciones por ser $\sigma$-álgebra), se
obtiene la independencia de $\sigma(A_1, A_2, \dots)$ y $\mathcal{T}$.

\textit{Conclusión.} Sea $B \in \mathcal{T}$. También $B \in \sigma(A_1, A_2, \dots)$, pues
$\mathcal{T} \subseteq \sigma(A_1, A_2, \dots)$. La independencia recién probada, aplicada al
par $(B, B)$, da
$$ \mathbb{P}(B) = \mathbb{P}(B \cap B) = \mathbb{P}(B) \cdot \mathbb{P}(B), $$
de donde $\mathbb{P}(B)(1 - \mathbb{P}(B)) = 0$ y $\mathbb{P}(B) \in \{0, 1\}$.
\end{proof}

\obs La ley no dice \emph{cuál} de los dos valores toma $\mathbb{P}(B)$, y decidirlo puede ser
difícil. Para el evento de cola $\varlimsup_n A_n$, los lemas de Borel-Cantelli son el
criterio efectivo.

\newpage

\section{Generación de medidas}

\subsection{Medida exterior}

\begin{definicion}{Medida exterior}{medida_exterior}
Una función $\mu^*: \mathcal{P}(X) \to [0, \infty]$ es una \emph{medida exterior} si:
\begin{enumerate}
    \item $\mu^*(\emptyset) = 0$.
    \item Monotonía: $A \subseteq B \implies \mu^*(A) \le \mu^*(B)$.
    \item Sub-$\sigma$-aditividad: $\mu^*(\bigcup_n A_n) \le \sum_n \mu^*(A_n)$ para cualquier
    $(A_n) \subseteq \mathcal{P}(X)$.
\end{enumerate}
\end{definicion}

\obs \emph{Construcción típica de medida exterior.} Dada una función $\ell:
\mathcal{C} \to [0, \infty]$ definida sobre una familia $\mathcal{C} \subseteq \mathcal{P}(X)$
con $\emptyset \in \mathcal{C}$ y $\ell(\emptyset) = 0$, se define
$$ \mu^*(E) = \inf\left\{ \sum_{n=1}^\infty \ell(I_n) : (I_n) \subset \mathcal{C}, \ E \subseteq
\bigcup_n I_n \right\}, $$
con $\mu^*(E) = +\infty$ si $E$ no admite cubrimiento. Entonces $\mu^*$ es una medida exterior
.

\subsection{Criterio de Carathéodory}

\begin{definicion}{Conjuntos $\mu^*$-medibles}{caratheodory_mide}
$A \subseteq X$ es \emph{$\mu^*$-medible} (en el sentido de Carathéodory) si para todo
$E \subseteq X$:
$$ \mu^*(E) = \mu^*(E \cap A) + \mu^*(E \cap A^c). $$
\end{definicion}

\begin{teorema}{Carathéodory}{caratheodory}
Sea $\mu^*$ una medida exterior sobre $X$, y sea
$$ \mathcal{A}^* = \{ A \subseteq X : A \text{ es } \mu^*\text{-medible} \}. $$
Entonces $\mathcal{A}^*$ es una $\sigma$-álgebra y $\mu = \mu^*|_{\mathcal{A}^*}$ es una medida.
\end{teorema}

\begin{proof}\leavevmode\\
Para todo $A \subseteq X$ y todo $E \subseteq X$, la sub-aditividad da $\mu^*(E) = \mu^*((E \cap A) \cup (E \cap A^c)) \le \mu^*(E \cap A) + \mu^*(E \cap A^c)$. Por lo tanto, para ver que un conjunto está en $\mathcal{A}^*$ basta probar la desigualdad opuesta.

\textit{Paso 1: $\mathcal{A}^*$ es un álgebra.} El vacío está en $\mathcal{A}^*$, ya que $\mu^*(E \cap \emptyset) + \mu^*(E \cap X) = 0 + \mu^*(E) = \mu^*(E)$ para todo $E$. El cierre por complementos se sigue de que la condición de Carathéodory es simétrica en $A$ y $A^c$. Para el cierre por uniones finitas, sean $A, B \in \mathcal{A}^*$ y $E \subseteq X$. Aplicando la medibilidad de $A$ a $E$, y luego la de $B$ a $E \cap A^c$,
\begin{align*}
\mu^*(E) &= \mu^*(E \cap A) + \mu^*(E \cap A^c) \\
&= \mu^*(E \cap A) + \mu^*(E \cap A^c \cap B) + \mu^*(E \cap A^c \cap B^c).
\end{align*}
Como $E \cap (A \cup B) = (E \cap A) \sqcup (E \cap A^c \cap B)$, la sub-aditividad da $\mu^*(E \cap (A \cup B)) \le \mu^*(E \cap A) + \mu^*(E \cap A^c \cap B)$. Sustituyendo en la igualdad anterior y usando $E \cap (A \cup B)^c = E \cap A^c \cap B^c$,
$$ \mu^*(E) \ge \mu^*(E \cap (A \cup B)) + \mu^*(E \cap (A \cup B)^c), $$
de modo que $A \cup B \in \mathcal{A}^*$.

\textit{Paso 2: $\mathcal{A}^*$ es una $\sigma$-álgebra.} Sea $(A_n)_{n \in \mathbb{N}} \subseteq \mathcal{A}^*$ una sucesión disjunta y $A = \bigcup_n A_n$. Se escribe $S_n = \bigcup_{k=1}^n A_k$. Se prueba por inducción en $n$ que, para todo $E$,
$$ \mu^*(E \cap S_n) = \sum_{k=1}^n \mu^*(E \cap A_k). $$
El caso $n = 1$ vale porque $S_1 = A_1$. Supuesta la fórmula para $n$, obsérvese que $S_n \in \mathcal{A}^*$ por ser unión finita y se aplica la condición de Carathéodory para $A_{n+1}$ al conjunto $E \cap S_{n+1}$:
$$ \mu^*(E \cap S_{n+1}) = \mu^*(E \cap S_{n+1} \cap A_{n+1}) + \mu^*(E \cap S_{n+1} \cap A_{n+1}^c). $$
Como los $A_k$ son disjuntos, $S_{n+1} \cap A_{n+1} = A_{n+1}$ y $S_{n+1} \cap A_{n+1}^c = S_n$, de modo que
$$ \mu^*(E \cap S_{n+1}) = \mu^*(E \cap A_{n+1}) + \mu^*(E \cap S_n). $$
Aplicando la hipótesis inductiva a $\mu^*(E \cap S_n)$,
$$ \mu^*(E \cap S_{n+1}) = \mu^*(E \cap A_{n+1}) + \sum_{k=1}^n \mu^*(E \cap A_k) = \sum_{k=1}^{n+1} \mu^*(E \cap A_k), $$
que es la fórmula para $n+1$.

Ahora, por la medibilidad de $S_n$,
$$ \mu^*(E) = \mu^*(E \cap S_n) + \mu^*(E \cap S_n^c). $$
En el primer término se usa la fórmula recién probada; en el segundo, como $S_n \subseteq A$ da $S_n^c \supseteq A^c$, la monotonía asegura $\mu^*(E \cap S_n^c) \ge \mu^*(E \cap A^c)$. Por lo tanto
$$ \mu^*(E) \ge \sum_{k=1}^n \mu^*(E \cap A_k) + \mu^*(E \cap A^c), $$
y como esto vale para todo $n$, tomando $n \to \infty$,
$$ \mu^*(E) \ge \sum_{k=1}^\infty \mu^*(E \cap A_k) + \mu^*(E \cap A^c). $$
Por último, $E \cap A = \bigcup_k (E \cap A_k)$, de modo que la sub-$\sigma$-aditividad da $\sum_{k=1}^\infty \mu^*(E \cap A_k) \ge \mu^*(E \cap A)$, y entonces
$$ \mu^*(E) \ge \mu^*(E \cap A) + \mu^*(E \cap A^c), $$
es decir, $A \in \mathcal{A}^*$. Para una unión numerable cualquiera $(B_n) \subseteq \mathcal{A}^*$, se disjuntiza: $A_1 = B_1$ y $A_n = B_n \setminus (B_1 \cup \cdots \cup B_{n-1})$ para $n \ge 2$. Como $\mathcal{A}^*$ es un álgebra, cada $A_n \in \mathcal{A}^*$, los $A_n$ son disjuntos y $\bigcup_n B_n = \bigcup_n A_n \in \mathcal{A}^*$ por lo anterior. Luego $\mathcal{A}^*$ es una $\sigma$-álgebra.

\textit{Paso 3: $\mu = \mu^*|_{\mathcal{A}^*}$ es una medida.} Sea $(A_n) \subseteq \mathcal{A}^*$ disjunta y $A = \bigcup_n A_n$. Se evalúa en $E = A$ la desigualdad obtenida en el Paso 2. Como $A \cap A_k = A_k$ y $A \cap A^c = \emptyset$,
$$ \mu^*(A) \ge \sum_{k=1}^\infty \mu^*(A \cap A_k) + \mu^*(A \cap A^c) = \sum_{k=1}^\infty \mu^*(A_k). $$
La desigualdad opuesta $\mu^*(A) \le \sum_k \mu^*(A_k)$ es la sub-$\sigma$-aditividad, de modo que $\mu(A) = \sum_n \mu(A_n)$; es decir, $\mu$ es $\sigma$-aditiva sobre sucesiones disjuntas de $\mathcal{A}^*$. Como además $\mu(\emptyset) = \mu^*(\emptyset) = 0$, $\mu$ es una medida.
\end{proof}

\subsection{Medida de Lebesgue en \texorpdfstring{$\mathbb{R}$}{R}}

\begin{definicion}{Medida exterior de Lebesgue}{lebesgue_ext}
Para $E \subseteq \mathbb{R}$, se define
$$ \lambda^*(E) = \inf\!\left\{ \sum_{n=1}^\infty (b_n - a_n) : E \subseteq \bigcup_n (a_n, b_n]
\right\}, $$
\end{definicion}

\begin{teorema}{Medida de Lebesgue}{lebesgue}
$\lambda^*$ es una medida exterior. La $\sigma$-álgebra de Carathéodory asociada, anotada
$\mathcal{L}$, contiene a $\mathcal{B}(\mathbb{R})$, y la restricción $\lambda =
\lambda^*|_{\mathcal{L}}$ es una medida (llamada \emph{medida de Lebesgue}) que cumple:
\begin{enumerate}
    \item $\lambda(a, b]) = b - a$.
    \item $\mathcal{B}(\mathbb{R}) \subset \mathcal{L} \subsetneq \mathcal{P}(\mathbb{R})$.
    \item $\lambda$ es invariante por traslaciones: $\lambda(A + t) = \lambda(A)$ para todo $t \in
    \mathbb{R}$ y $A \in \mathcal{L}$.
    \item $(\mathbb{R}, \mathcal{L}, \lambda)$ es completo: si $\lambda(N) = 0$ y $A \subseteq N$,
    entonces $A \in \mathcal{L}$ (y por monotonía $\lambda(A) = 0$).
\end{enumerate}
\end{teorema}

\begin{proof}\leavevmode\\
(i) La desigualdad $\lambda^*((a, b]) \le b - a$ es
se obtiene tomando como cubrimiento el propio $(a, b]$. Se prueba ahora la otra desigualdad. Sea
$\{(a_k, b_k]\}_k$ un cubrimiento numerable de $(a, b]$. Dado $\epsilon > 0$, se agranda cada
intervalo a uno \emph{abierto} $(a_k - \epsilon_k, b_k + \epsilon_k)$ eligiendo $\epsilon_k > 0$
de modo que $\sum_k \epsilon_k < \epsilon$. Así, $\{(a_k - \epsilon_k, b_k + \epsilon_k)\}$ es un
cubrimiento abierto del compacto $[a, b]$.

Por compacidad, hay un subcubrimiento finito $\{(a_{k_1} - \epsilon_{k_1}, b_{k_1} +
\epsilon_{k_1}), \dots, (a_{k_N} - \epsilon_{k_N}, b_{k_N} + \epsilon_{k_N})\}$ de $[a, b]$. La suma de las longitudes
de los intervalos del subcubrimiento es al menos $b - a$:
$$ b - a \le \sum_{j=1}^N (b_{k_j} + \epsilon_{k_j} - (a_{k_j} - \epsilon_{k_j})) = \sum_{j=1}^N
(b_{k_j} - a_{k_j}) + 2 \sum_{j=1}^N \epsilon_{k_j} \le \sum_k (b_k - a_k) + 2\epsilon. $$
Tomando $\epsilon \to 0$ y luego ínfimo sobre cubrimientos, $b - a \le \lambda^*((a, b])$.

(ii) Basta probar que cada semirrecta
$(c, \infty)$ es $\lambda^*$-medible, porque las semirrectas generan $\mathcal{B}(\mathbb{R})$.
Fíjese $c \in \mathbb{R}$ y $E \subseteq \mathbb{R}$. Dado un cubrimiento $E \subseteq \bigcup_k
(a_k, b_k]$, se parte cada intervalo en su pieza a la derecha de $c$ y su pieza a la izquierda:
$$ (a_k, b_k] = ((a_k, b_k] \cap (c, \infty)) \sqcup ((a_k, b_k] \cap (-\infty, c]). $$
Cada pieza es un intervalo semiabierto (eventualmente vacío) cuya longitud suma exactamente
$b_k - a_k$. Aplicando $\lambda^*$ a las dos colecciones de piezas, que cubren respectivamente
$E \cap (c, \infty)$ y $E \cap (-\infty, c]$,
$$ \lambda^*(E \cap (c, \infty)) + \lambda^*(E \cap (-\infty, c]) \le \sum_k (b_k - a_k). $$
Tomando ínfimo sobre cubrimientos en la derecha,
$$ \lambda^*(E \cap (c, \infty)) + \lambda^*(E \cap (c, \infty)^c) \le \lambda^*(E). $$
La desigualdad opuesta es sub-aditividad. Luego $(c, \infty)$ cumple Carathéodory y pertenece a
$\mathcal{L}$.

(iii) Si $\{(a_k, b_k]\}$ cubre $E$, entonces
$\{(a_k + t, b_k + t]\}$ cubre $E + t$ con la misma suma de longitudes. Esto da $\lambda^*(E + t)
\le \lambda^*(E)$, y por simetría (aplicando lo mismo a $-t$) la igualdad. La invariancia se
traslada a la $\lambda^*$-medibilidad: si $A \in \mathcal{L}$, la condición de Carathéodory
para $A + t$ sale aplicando la de $A$ al conjunto $E - t$.

(iv) Si $\lambda(N) = 0$ y $A \subseteq N$, por monotonía $\lambda^*(A)
= 0$. Para cualquier $E$, $\lambda^*(E \cap A) \le \lambda^*(A) = 0$, por lo cual la condición de
Carathéodory para $A$ se reduce a $\lambda^*(E) = \lambda^*(E \cap A^c)$, que vale por
sub-aditividad y monotonía. Luego $A \in \mathcal{L}$.

(v) Para $\mathcal{L} \subsetneq \mathcal{P}(\mathbb{R})$, los conjuntos
de Vitali no pueden ser Lebesgue-medibles porque $\lambda$ cumple las
tres propiedades de la conclusión del teorema.
\end{proof}

\obs Existen conjuntos Lebesgue-nulos con cardinal del continuo (Cantor), así como conjuntos
de medida positiva y nunca densos (Cantor gordo).

\obs Todo conjunto numerable tiene medida de Lebesgue cero: para cada punto,
$\lambda(\{x\}) \le \lambda((x - \epsilon, x]) = \epsilon$ para todo $\epsilon > 0$, de donde
$\lambda(\{x\}) = 0$, y por $\sigma$-aditividad
$$ \lambda(\mathbb{Q}) = \sum_{q \in \mathbb{Q}} \lambda(\{q\}) = 0. $$
Así, la función $\mathbf{1}_{\mathbb{Q} \cap [0,1]}$ de la motivación, que no era
Riemann-integrable, es nula salvo en un conjunto de medida cero: su integral de Lebesgue
será $0$.

\subsection{Lebesgue-Stieltjes}

\begin{definicion}{Medida exterior de Lebesgue-Stieltjes}{lebesgue_stieltjes_ext}
Sea $F: \mathbb{R} \to \mathbb{R}$ no decreciente y continua a derecha. Para $E \subseteq
\mathbb{R}$,
$$ \lambda_F^*(E) = \inf\!\left\{ \sum_{n=1}^\infty (F(b_n) - F(a_n)) : E \subseteq \bigcup_n
(a_n, b_n] \right\}. $$
\end{definicion}

\begin{teorema}{Medida de Lebesgue-Stieltjes}{lebesgue_stieltjes}
Sea $F$ como arriba. Entonces $\lambda_F^*$ es una medida exterior, y la $\sigma$-álgebra de
Carathéodory $\mathcal{L}_F$ contiene a $\mathcal{B}(\mathbb{R})$. La restricción $\lambda_F =
\lambda_F^*|_{\mathcal{L}_F}$ es la \emph{medida de Lebesgue-Stieltjes} de $F$, y cumple
$$ \lambda_F((a, b]) = F(b) - F(a) \qquad \forall a < b. $$
\end{teorema}

\begin{proof}
La identidad $\lambda_F((a, b]) = F(b) - F(a)$ usa, en el paso compacto, que
$F$ es continua a derecha: para mostrar
$F(b) - F(a) \le \sum_k (F(b_k) - F(a_k))$ cuando $\{(a_k, b_k]\}$ cubre $(a, b]$, se ensanchan los
intervalos a $(a_k, b_k + \delta_k]$ con $F(b_k + \delta_k) - F(b_k) < \epsilon/2^k$
 y se compactifica $[a + \delta, b]$ como en Lebesgue.
La $\lambda_F^*$-medibilidad de las semirrectas y la completitud salen exactamente igual.
\end{proof}

\obs Cuando $F$ es la
función de distribución de una variable aleatoria $X$, $\lambda_F$ es la
distribución de $X$, anotada $P_X$, y vale $P_X(B) = P(X \in B)$ para todo boreliano $B$.

\newpage

\section{Funciones medibles}

\begin{definicion}{Función medible}{medible}
Sea $(X, \mathcal{A}, \mu)$ un espacio de medida. Una función $f: X \to \mathbb{R}$ es
\emph{medible} si para todo $a \in \mathbb{R}$,
$$ f^{-1}((a, \infty)) = \{ x \in X : f(x) > a \} \in \mathcal{A}. $$
\end{definicion}

\obs En probabilidad, una función medible se llama \emph{variable aleatoria}.

\obs La definición apunta directamente a la integral de Lebesgue: para aproximar $\int f$
partiendo la \emph{imagen} en valores $a_0 < a_1 < \cdots < a_n$, hay que medir los conjuntos
$$ A_k = \{a_{k-1} < f \le a_k\} = \{f > a_{k-1}\} \setminus \{f > a_k\}. $$
Pedir $\{f > a\} \in \mathcal{A}$ para todo $a$ es exactamente lo que garantiza que esos
conjuntos tengan medida.

\begin{proposicion}{Caracterizaciones equivalentes de medibilidad}{caract_medible}
Sea $(X, \mathcal{A})$ un espacio medible y $f: X \to \mathbb{R}$. Las siguientes condiciones son
equivalentes:
\begin{enumerate}
    \item $\{f > a\} \in \mathcal{A}$ para todo $a \in \mathbb{R}$.
    \item $\{f \ge a\} \in \mathcal{A}$ para todo $a \in \mathbb{R}$.
    \item $\{f < a\} \in \mathcal{A}$ para todo $a \in \mathbb{R}$.
    \item $\{f \le a\} \in \mathcal{A}$ para todo $a \in \mathbb{R}$.
\end{enumerate}
\end{proposicion}

\begin{proof}\leavevmode\\
Se prueban (i) $\Leftrightarrow$ (iv), (ii) $\Leftrightarrow$ (iii) y (i) $\Leftrightarrow$ (ii).

\textit{(i) $\Leftrightarrow$ (iv).} $\{f \le a\} = \{f > a\}^c$, y $\mathcal{A}$ es cerrada por
complementos.

\textit{(ii) $\Leftrightarrow$ (iii).} $\{f < a\} = \{f \ge a\}^c$.

\textit{(i) $\Rightarrow$ (ii).} $\{f \ge a\} = \bigcap_{n=1}^\infty \{f > a - 1/n\}$,
intersección
numerable de elementos de $\mathcal{A}$.

\textit{(ii) $\Rightarrow$ (i).} $\{f > a\} = \bigcup_{n=1}^\infty \{f \ge a + 1/n\}$, unión
numerable de elementos de $\mathcal{A}$.

Con esto, (i)--(iv) son equivalentes.
\end{proof}

\subsection{Ejemplos y operaciones}

\begin{ejemplo}{Funciones constantes}{}
Si $f \equiv c$, entonces $\{f > a\}$ vale $X$ si $c > a$ y $\emptyset$ si $c \le a$. Como
$X, \emptyset \in \mathcal{A}$, $f$ es medible.
\end{ejemplo}

\begin{ejemplo}{Funciones características}{}
Sea $A \in \mathcal{A}$. La función característica $\mathbf{1}_A$ es medible:
$$ \{\mathbf{1}_A > a\} = \begin{cases} X, & a < 0, \\ A, & 0 \le a < 1, \\ \emptyset,
& a \ge 1, \end{cases} $$
y los tres conjuntos están en $\mathcal{A}$.
\end{ejemplo}

\begin{teorema}{Las continuas son Borel-medibles}{continuas_borel}
Toda función continua $f: \mathbb{R} \to \mathbb{R}$ es Borel-medible.
\end{teorema}

\begin{proof}
Hay que ver que $f^{-1}((a, \infty)) \in \mathcal{B}(\mathbb{R})$ para todo $a \in \mathbb{R}$.
Como $(a, \infty)$ es abierto y $f$ continua, $f^{-1}((a, \infty))$ es abierto en $\mathbb{R}$, y
todo abierto pertenece a $\mathcal{B}(\mathbb{R})$.

Más en general, la preimagen por una continua de cualquier abierto es abierta, lo cual da
$f^{-1}(\mathcal{B}(\mathbb{R})) \subseteq \mathcal{B}(\mathbb{R})$: el conjunto $\{B \subseteq
\mathbb{R} : f^{-1}(B) \in \mathcal{B}(\mathbb{R})\}$ es una $\sigma$-álgebra que contiene a los
abiertos, por lo cual contiene a $\mathcal{B}(\mathbb{R})$.
\end{proof}

\begin{teorema}{Las monótonas son Borel-medibles}{monotonas_borel}
Toda función monótona $f: \mathbb{R} \to \mathbb{R}$ es Borel-medible.
\end{teorema}

\begin{proof}
Supóngase $f$ no decreciente (el caso no creciente sale aplicando lo siguiente a $-f$). Fijado
$a \in \mathbb{R}$, considérese el conjunto
$E_a = f^{-1}((a, \infty)) = \{x \in \mathbb{R} : f(x) > a\}$.

Si $x \in E_a$ e $y > x$, la monotonía da $f(y) \ge f(x) > a$, de modo que $y \in E_a$. Por lo
tanto $E_a$ es una semirrecta hacia la derecha: si se denota $c = \inf E_a$ (con la convención
$\inf \emptyset = +\infty$, $\inf$ de no acotado inferiormente $= -\infty$), entonces
$$ E_a = (c, \infty) \quad \text{o bien} \quad E_a = [c, \infty), $$
según el ínfimo se alcance o no. En cualquiera de los dos casos, $E_a \in
\mathcal{B}(\mathbb{R})$ (los semirrectas abiertas y cerradas son borelianas).
\end{proof}

\begin{ejemplo}{Funciones simples}{}
Una función $\varphi: X \to \mathbb{R}$ es \emph{simple} si es combinación lineal finita de
características:
$$ \varphi = \sum_{i=1}^n c_i \, \mathbf{1}_{A_i}, \qquad c_i \in \mathbb{R}, \ A_i \in \mathcal{A}.
$$
Equivalentemente, $\varphi$ toma sólo finitos valores y la preimagen de cada uno está en
$\mathcal{A}$. Las simples son medibles porque cada $\mathbf{1}_{A_i}$ lo es y las operaciones de
suma y producto por escalar preservan la medibilidad (Teorema \ref{teo:operaciones_medibles}).
\end{ejemplo}

\begin{ejemplo}{Representación canónica de una simple}{ej_repr_canonica}
La escritura de una simple como combinación de características no es única. Con intervalos de
$\mathbb{R}$, considérese
$$ \varphi = 2 \cdot \mathbf{1}_{[0,2]} + 3 \cdot \mathbf{1}_{[1,5]}. $$
Los dos intervalos se superponen en $[1,2]$, donde $\varphi$ vale $2 + 3 = 5$. Separando según
los valores que toma $\varphi$,
$$ \varphi = 2 \cdot \mathbf{1}_{[0,1)} + 5 \cdot \mathbf{1}_{[1,2]} + 3 \cdot \mathbf{1}_{(2,5]},
$$
que es la \emph{representación canónica}: los coeficientes son los valores distintos no nulos
de $\varphi$, y los conjuntos $\{\varphi = 2\} = [0,1)$, $\{\varphi = 5\} = [1,2]$ y
$\{\varphi = 3\} = (2,5]$ son disjuntos. Toda simple admite una única representación de esta
forma, que es la que usa la definición de integral.
\end{ejemplo}

\begin{teorema}{Operaciones con medibles}{operaciones_medibles}
Sean $f, g: X \to \mathbb{R}$ medibles y $c \in \mathbb{R}$. Entonces $cf$, $f + g$, $f \cdot g$,
$\max(f, g)$, $\min(f, g)$, $|f|$, $f^+ := \max(f, 0)$ y $f^- := \max(-f, 0)$ son medibles.
\end{teorema}

\begin{proof}\leavevmode\\
Por la Proposición \ref{prop:caract_medible}, para ver que una función es medible basta
verificar cualquiera de las condiciones $\{f > a\}, \{f \ge a\}, \{f < a\}, \{f \le a\} \in
\mathcal{A}$ para todo $a \in \mathbb{R}$; en cada caso se elige la que más convenga.

\textit{Producto por escalar.}
\begin{itemize}
    \item Si $c = 0$, $cf \equiv 0$, que es constante y por lo tanto medible.
    \item Si $c > 0$, $cf(x) > a \Leftrightarrow f(x) > a/c$, de modo que
    $\{cf > a\} = \{f > a/c\} \in \mathcal{A}$.
    \item Si $c < 0$, $cf(x) > a \Leftrightarrow f(x) < a/c$, de modo que 
    $\{cf > a\} = \{f < a/c\} \in \mathcal{A}$.
\end{itemize}

\textit{Máximo y mínimo.} Como
$$ \max(f, g)(x) > a \ \Leftrightarrow \ f(x) > a \ \text{ o } \ g(x) > a, $$
se tiene $\{\max(f, g) > a\} = \{f > a\} \cup \{g > a\} \in \mathcal{A}$ por ser unión finita
de medibles. Análogamente,
$$ \min(f, g)(x) \ge a \ \Leftrightarrow \ f(x) \ge a \ \text{ y } \ g(x) \ge a, $$
de modo que $\{\min(f, g) \ge a\} = \{f \ge a\} \cap \{g \ge a\} \in \mathcal{A}$.

\textit{Valor absoluto.} Para $a < 0$, $\{|f| > a\} = X \in \mathcal{A}$. Para $a \ge 0$,
$$ |f(x)| > a \ \Leftrightarrow \ f(x) > a \ \text{ o } \ f(x) < -a, $$
de modo que $\{|f| > a\} = \{f > a\} \cup \{f < -a\} \in \mathcal{A}$.

\textit{Partes positiva y negativa $f^+, f^-$.} La función constante $0$ es medible, y por el
caso del máximo aplicado a $f$ y $0$ (o a $-f$ y $0$, usando el caso de producto por escalar con
$c = -1$),
$$ f^+ = \max(f, 0), \qquad f^- = \max(-f, 0), $$
son medibles. Además valen las identidades $f = f^+ - f^-$ y $|f| = f^+ + f^-$.

\textit{Suma.} Véase que para todo $a \in \mathbb{R}$,
$$ \{f + g > a\} = \bigcup_{r \in \mathbb{Q}} \left( \{f > r\} \cap \{g > a - r\} \right). $$
\begin{itemize}
    \item ($\subseteq$) Sea $x$ con $f(x) + g(x) > a$, es decir $f(x) > a - g(x)$. Por densidad de
    $\mathbb{Q}$ en $\mathbb{R}$, existe $r \in \mathbb{Q}$ con
    $$ a - g(x) < r < f(x). $$
    De $r < f(x)$ sale $x \in \{f > r\}$, y de $a - g(x) < r$ sale $g(x) > a - r$, es decir $x \in
    \{g > a - r\}$.

    \item ($\supseteq$) Sea $r \in \mathbb{Q}$ con $f(x) > r$ y $g(x) > a - r$. Sumando, $f(x) +
    g(x) > r + (a - r) = a$, es decir $x \in \{f + g > a\}$.
\end{itemize}
El lado derecho es una unión \emph{numerable} de
intersecciones finitas de conjuntos en $\mathcal{A}$, y por lo tanto pertenece a $\mathcal{A}$.

\textit{Cuadrado.} Para $a < 0$, $\{f^2 > a\} = X \in \mathcal{A}$. Para $a \ge 0$,
$$ f(x)^2 > a \ \Leftrightarrow \ |f(x)| > \sqrt{a} \ \Leftrightarrow \ f(x) > \sqrt{a} \ \text{ o }
\ f(x) < -\sqrt{a}, $$
y por lo tanto $\{f^2 > a\} = \{f > \sqrt{a}\} \cup \{f < -\sqrt{a}\} \in \mathcal{A}$.

\textit{Producto.} Se usa la identidad de polarización:
$$ fg = \tfrac{1}{4} \left[ (f + g)^2 - (f - g)^2 \right]. $$
Por los pasos previos, $f + g$ y $f - g = f + (-1) g$ son medibles.
Sus cuadrados $(f + g)^2$ y $(f - g)^2$ son medibles por el caso anterior. Aplicando suma y
producto por escalar una vez más,
$$ fg = \tfrac{1}{4}(f + g)^2 + \left(-\tfrac{1}{4}\right)(f - g)^2 \in \mathcal{A}, $$
es decir, $fg$ es medible.
\end{proof}

\obs Las partes positiva y negativa satisfacen, además de $f = f^+ - f^-$ y $|f| = f^+ + f^-$,
las fórmulas inversas
$$ f^+ = \tfrac{1}{2}(|f| + f), \qquad f^- = \tfrac{1}{2}(|f| - f), $$
con $f^+, f^- \ge 0$ de soportes disjuntos. La descomposición $f = f^+ - f^-$ es la que reduce
la integral de una función general a la de dos no negativas.

\subsection{Sup, inf, lim sup, lim inf}

\begin{teorema}{Sup e inf de medibles}{sup_inf_medibles}
Si $(f_n)$ es una sucesión de funciones medibles, entonces $\sup_n f_n$, $\inf_n f_n$,
$\varlimsup_n f_n$ y $\varliminf_n f_n$ son medibles.
\end{teorema}

\begin{proof}\leavevmode\\
\textit{Caso del supremo.} Para todo $a \in \mathbb{R}$,
$$ \sup_n f_n(x) > a \ \Leftrightarrow \ \exists \, n \in \mathbb{N} : f_n(x) > a, $$
de donde
$$ \{\sup_n f_n > a\} = \bigcup_{n \in \mathbb{N}} \{f_n > a\}. $$
Cada $\{f_n > a\} \in \mathcal{A}$ por medibilidad de $f_n$, y la unión es numerable, de modo que
$\{\sup_n f_n > a\} \in \mathcal{A}$ para todo $a$.

\textit{Caso del ínfimo.} Análogamente,
$$ \inf_n f_n(x) < a \ \Leftrightarrow \ \exists \, n \in \mathbb{N} : f_n(x) < a, $$
de modo que $\{\inf_n f_n < a\} = \bigcup_n \{f_n < a\} \in \mathcal{A}$.

\textit{Límite superior.} Por definición,
$$ \varlimsup_n f_n = \inf_{m \in \mathbb{N}} \sup_{n \ge m} f_n. $$
Para cada $m$, la sucesión $(f_n)_{n \ge m}$ es medible, de modo que $g_m := \sup_{n \ge m} f_n$
es medible por el caso del sup. La sucesión $(g_m)_m$ es entonces medible, y
$\varlimsup_n f_n = \inf_m g_m$ es medible por el caso del inf.

\textit{Límite inferior.} Análogo, intercambiando $\sup$ con $\inf$: $\varliminf_n f_n =
\sup_m \inf_{n \ge m} f_n$, donde $\inf_{n \ge m} f_n$ es medible por el caso del inf y luego el
sup en $m$ por el caso del sup.
\end{proof}

\obs Las cuatro operaciones $\sup, \inf, \varlimsup, \varliminf$ pueden dar funciones con valores
en $[-\infty, +\infty]$. La definición de medibilidad se extiende sin problemas:
$f: X \to [-\infty,
+\infty]$ es medible si $\{f > a\} \in \mathcal{A}$ para todo $a \in \mathbb{R}$.

\begin{corolario}{Límite puntual de medibles es medible}{lim_medibles}
Sea $(f_n)$ una sucesión de funciones medibles.
\begin{enumerate}
    \item Si $f_n \to f$ puntualmente en todo $X$, $f$ es medible.
    \item Si la serie $g := \sum_{k=1}^\infty f_k$ converge puntualmente, $g$ es medible.
\end{enumerate}
\end{corolario}

\begin{proof}\leavevmode\\
\textit{(i)} Si $f_n \to f$ puntualmente, $f = \varlimsup_n f_n = \varliminf_n f_n$, y por el
teorema anterior ambos son medibles. Luego $f$ es medible.

\textit{(ii)} Las sumas parciales $g_n = \sum_{k=1}^n f_k$ son sumas finitas de medibles, y por
el Teorema \ref{teo:operaciones_medibles} son medibles. Como $g = \lim_n g_n$ puntualmente (por
hipótesis), por (i) $g$ es medible.
\end{proof}

\subsection{Aproximación por simples}

\begin{teorema}{Aproximación por simples}{aprox_simples}
Sea $f: X \to \mathbb{R}$ medible y no negativa. Existe una sucesión $(\varphi_n)$ de funciones
simples no negativas tal que $\varphi_n \nearrow f$ puntualmente. Si además $f$ es acotada, la
convergencia es uniforme.
\end{teorema}

\begin{proof}\leavevmode\\
La idea es discretizar los \emph{valores} de $f$. Para $n \in \mathbb{N}$ y $0 \le k \le n 2^n - 1$,
se escribe $I_{n,k} = [k/2^n, (k+1)/2^n)$ y se define
$$ \varphi_n(x) = \begin{cases}
\dfrac{k}{2^n}, & \text{si } f(x) \in I_{n,k}, \\
n, & \text{si } f(x) \ge n.
\end{cases} $$

Los intervalos $I_{n,0}, \dots, I_{n, n2^n - 1}$ junto con $[n, \infty)$ forman una partición de
$[0, \infty)$, de modo que $\varphi_n$ está bien definida y toma a lo sumo $n 2^n + 1$ valores.
Sus conjuntos de nivel son las preimágenes $\{\varphi_n = k/2^n\} = f^{-1}(I_{n,k})$ y
$\{\varphi_n = n\} = f^{-1}([n, \infty))$, medibles por serlo $f$; luego $\varphi_n$ es simple. Por
otro lado, $\varphi_n(x) \le f(x)$ en cada caso ($k/2^n \le f(x)$ cuando $f(x) \in I_{n,k}$, y
$n \le f(x)$ cuando $f(x) \ge n$), de modo que $0 \le \varphi_n \le f$.

Véase que $\varphi_n \le \varphi_{n+1}$ puntualmente. Como $I_{n,k} = I_{n+1, 2k} \sqcup
I_{n+1, 2k+1}$, si $f(x) \in I_{n,k}$ con $k \le n 2^n - 1$ entonces $f(x)$ pertenece a uno de esos
dos intervalos, sobre los cuales $\varphi_{n+1}(x)$ vale $2k/2^{n+1} = k/2^n$ o
$(2k+1)/2^{n+1} > k/2^n$; en ambos $\varphi_{n+1}(x) \ge k/2^n = \varphi_n(x)$. Si en cambio
$f(x) \ge n$, entonces $\varphi_n(x) = n$ y $\varphi_{n+1}(x) \ge n$: vale $n+1$ si $f(x) \ge n+1$,
y si $n \le f(x) < n+1$ existe $k$ con $f(x) \in I_{n+1,k}$ y $k/2^{n+1} \ge n$, de donde
$\varphi_{n+1}(x) = k/2^{n+1} \ge n$. En todos los casos $\varphi_{n+1}(x) \ge \varphi_n(x)$.

Sea ahora $x$ con $f(x) < \infty$ y tómese $n > f(x)$. Entonces $f(x) < n$, de modo que
$f(x) \in I_{n,k}$ para algún $k$ y, por definición de $I_{n,k}$,
$$ 0 \le f(x) - \varphi_n(x) = f(x) - \frac{k}{2^n} < \frac{1}{2^n}. $$
Como $2^{-n} \to 0$, se obtiene $\varphi_n(x) \to f(x)$.

Por último, si $f \le M$, para todo $n > M$ se cumple $f(x) < n$ cualquiera sea $x$, de modo que la
cota $0 \le f(x) - \varphi_n(x) < 2^{-n}$ vale uniformemente en $x$. Como $2^{-n} \to 0$, la
convergencia es uniforme.
\end{proof}

\newpage

\section{Convergencia de funciones medibles}

\begin{definicion}{Convergencia en casi todo punto}{conv_ctp}
$f_n \to f$ \emph{en casi todo punto} (c.t.p.) si existe $N \in \mathcal{A}$ con
$\mu(N) = 0$ tal que $f_n(x) \to f(x)$ para todo $x \notin N$.
\end{definicion}

\begin{definicion}{Convergencia en medida}{conv_medida}
$f_n \to f$ \emph{en medida} si para todo $\epsilon > 0$,
$$ \mu(\{x : |f_n(x) - f(x)| \ge \epsilon\}) \xrightarrow{n \to \infty} 0. $$
\end{definicion}

\begin{definicion}{Convergencia casi uniforme}{conv_casi_unif}
$f_n \to f$ \emph{casi uniformemente} si para todo $\delta > 0$ existe $N \in \mathcal{A}$ con
$\mu(N) < \delta$ tal que $f_n \to f$ uniformemente en $N^c$.
\end{definicion}

\begin{corolario}{Límite c.t.p. de medibles}{lim_ctp_medible}
Si $\mathcal{A}$ es completa respecto de $\mu$, las $f_n$ son medibles y $f_n \to f$ c.t.p.,
entonces $f$ es medible.
\end{corolario}

\begin{proof}
Existe $N \in \mathcal{A}$ con $\mu(N) = 0$ tal que $f_n(x) \to f(x)$ para todo $x \notin N$. Sea
$\tilde f = \varlimsup_n f_n$, medible por el Teorema \ref{teo:sup_inf_medibles}; en $N^c$ la
sucesión converge, de modo que $f = \tilde f$ sobre $N^c$. Fijado $a \in \mathbb{R}$, se descompone
$$ \{f > a\} = \big(\{f > a\} \cap N^c\big) \cup \big(\{f > a\} \cap N\big). $$
El primer conjunto coincide con $\{\tilde f > a\} \cap N^c$, intersección de dos medibles. El
segundo verifica $\{f > a\} \cap N \subseteq N$, y como $\mathcal{A}$ es completa, todo subconjunto
de un conjunto $\mu$-nulo es medible. Por lo tanto $\{f > a\}$ es unión de dos medibles, y $f$ es
medible.
\end{proof}

\obs Sin completitud de $\mathcal{A}$, la convergencia c.t.p. no garantiza medibilidad: si
$\mu(N) = 0$ pero $N$ contiene subconjuntos no medibles, redefinir $f$ arbitrariamente en $N$
puede romper la medibilidad.

\obs En general, en cualquier espacio de medida valen las implicaciones
$$ \text{uniforme} \implies \text{puntual} \implies \text{c.t.p.}, \qquad
   \text{casi uniforme} \implies \text{en medida}, $$
y la casi uniforme también implica c.t.p.
Pero \emph{ninguna} otra implicación vale automáticamente. Ejemplos típicos:
\begin{itemize}
    \item $f_n = \mathbf{1}_{[n, n+1]}$ en $\mathbb{R}$ converge puntualmente
     a 0, pero \emph{no} en medida.
    \item Las funciones: $\mathbf{1}_{[0,1]}, \mathbf{1}_{[0,1/2]},
    \mathbf{1}_{[1/2,1]}, \mathbf{1}_{[0,1/4]}, \dots$ convergen a 0 en medida pero \emph{no}
    puntualmente en ningún punto.
    \item $f_n(x) = x^n$ en $([0, 1], \lambda)$ converge puntualmente a $\mathbf{1}_{\{1\}}$
    pero \emph{no} uniformemente: dado $0 < \epsilon < 1$ y $x < 1$, la condición
    $x^n < \epsilon$ exige $n > \ln \epsilon / \ln x$, cota que explota cuando $x \to 1^-$, de
    modo que ningún $n_0$ sirve para todos los $x$ a la vez (obsérvese además que el límite
    puntual no es continuo aunque cada $f_n$ lo sea). Sí converge \emph{casi uniformemente} a
    $\mathbf{1}_{\{1\}}$: dado $\delta > 0$ basta excluir $N = (1 - \delta, 1)$, y en
    $[0, 1 - \delta]$ la convergencia a $0$ es uniforme porque $x^n \le (1 - \delta)^n \to 0$.
\end{itemize}

\subsection{Teorema de Egoroff}

\begin{teorema}{Egoroff}{egoroff}
Sea $(X, \mathcal{A}, \mu)$ un espacio de medida finita ($\mu(X) < \infty$) y $f_n \to f$ c.t.p.
Entonces:
\begin{enumerate}
    \item $f_n \to f$ casi uniformemente (Definición \ref{def:conv_casi_unif}).
    \item Como consecuencia, $f_n \to f$ en medida.
\end{enumerate}
\end{teorema}

\begin{proof}
Fíjese $\delta > 0$. Para cada $p, m \in \mathbb{N}$, sea
$$ A_m^{(p)} = \{x \in X : |f_m(x) - f(x)| \ge 1/p\}. $$
Véase primero que, para cada $p$,
$$ \mu\!\left(\bigcup_{k \ge m} A_k^{(p)}\right) \xrightarrow{m \to \infty} 0. $$
En efecto, si $x \in \varlimsup_m A_m^{(p)} = \bigcap_m \bigcup_{k \ge m} A_k^{(p)}$, entonces para
cada $m$ existe $k \ge m$ con $|f_k(x) - f(x)| \ge 1/p$, de donde $f_k(x) \not\to f(x)$. Por lo
tanto $\varlimsup_m A_m^{(p)} \subseteq \{x \in X : f_n(x) \not\to f(x)\}$, que es $\mu$-nulo por
hipótesis. Como además $\big(\bigcup_{k \ge m} A_k^{(p)}\big)_m$ decrece a
$\varlimsup_m A_m^{(p)}$ y su primer término tiene medida $\le \mu(X) < \infty$, la continuidad de
la medida da el límite anunciado.

Se prueba (i). Por lo recién visto, para cada $p \in
\mathbb{N}$ se puede elegir $m_p$ tal que $\mu\big(\bigcup_{k \ge m_p} A_k^{(p)}\big) < \delta/2^p$, y
se define
$$ N = \bigcup_{p=1}^\infty \bigcup_{k \ge m_p} A_k^{(p)}, \qquad \text{con } \mu(N) \le
\sum_{p=1}^\infty \frac{\delta}{2^p} = \delta. $$
Si $x \in N^c$, entonces $|f_k(x) - f(x)| < 1/p$ para todo $p \in \mathbb{N}$ y todo $k \ge m_p$.
Así, dado $\epsilon > 0$ y eligiendo $p$ con $1/p < \epsilon$, resulta $|f_k(x) - f(x)| < \epsilon$
para todo $k \ge m_p$ y todo $x \in N^c$. Como $m_p$ no depende de $x$, la convergencia es uniforme
en $N^c$.

Por último, (ii) se deduce de (i): dados $\epsilon > 0$ y $\eta > 0$, por la convergencia casi
uniforme existen $N$ con $\mu(N) < \eta$ y $m_0 \in \mathbb{N}$ tales que $|f_m - f| < \epsilon$ en
$N^c$ para todo $m \ge m_0$. Entonces $\{|f_m - f| \ge \epsilon\} \subseteq N$ para $m \ge m_0$, de
donde $\mu(\{|f_m - f| \ge \epsilon\}) \le \mu(N) < \eta$. Como $\eta$ es arbitrario,
$\mu(\{|f_m - f| \ge \epsilon\}) \to 0$; es decir, $f_m \to f$ en medida.
\end{proof}

\obs El $m_p$ que define a $N$ depende de $p$ pero \emph{no} de $x$.

\newpage

\section{Integración}

Se construye la integral en tres pasos: simples, no negativas, generales.

\subsection{Paso 1: funciones simples positivas}

\begin{definicion}{Integral de una función simple positiva}{int_simple}
Sea $\varphi = \sum_{i=1}^n a_i \mathbf{1}_{A_i}$ simple con $a_i \ge 0$ y $A_i$ disjuntos.
Se define
$$ \int_X \varphi \, d\mu = \sum_{i=1}^n a_i \, \mu(A_i), $$
con la convención $0 \cdot \infty = 0$.
\end{definicion}

\obs La definición no depende de la representación (basta verificarlo para refinamientos).

\begin{proposicion}{Linealidad y monotonía en simples}{lin_mon_simple}
Si $\varphi, \psi$ son simples positivas y $a, b \ge 0$:
\begin{enumerate}
    \item $\int (a \varphi + b \psi) d\mu = a \int \varphi \, d\mu + b \int \psi \, d\mu$.
    \item Si $\varphi \le \psi$, $\int \varphi \, d\mu \le \int \psi \, d\mu$.
\end{enumerate}
Además, para $A \in \mathcal{A}$, $\lambda(A) := \int_A \varphi \, d\mu := \int_X \varphi
\mathbf{1}_A \, d\mu$ es una medida sobre $(X, \mathcal{A})$.
\end{proposicion}

\begin{proof}\leavevmode\\
\textit{(i)} Refinando ambas representaciones a una partición común, puede suponerse que
$\varphi = \sum_i \alpha_i \mathbf{1}_{C_i}$ y $\psi = \sum_i \beta_i \mathbf{1}_{C_i}$ usan la
misma partición $\{C_i\}$ disjunta. Entonces $a \varphi + b \psi = \sum_i (a \alpha_i + b
\beta_i) \mathbf{1}_{C_i}$, y por definición,
$$ \int (a \varphi + b \psi) d\mu = \sum_i (a \alpha_i + b \beta_i) \mu(C_i) = a \sum_i \alpha_i
\mu(C_i) + b \sum_i \beta_i \mu(C_i) = a \int \varphi \, d\mu + b \int \psi \, d\mu. $$

\textit{(ii)} Con la misma partición común, $\varphi \le \psi$ implica $\alpha_i \le \beta_i$
en cada $C_i$ no vacío. Sumando con $\mu(C_i) \ge 0$ se preserva la desigualdad:
$\int \varphi \, d\mu \le \int \psi \, d\mu$.

\textit{$\lambda$ es medida.} $\lambda(\emptyset) = 0$ es directo. Para la $\sigma$-aditividad,
sea $(A_n)$ disjunta y $A = \bigsqcup_n A_n$. Con $\varphi = \sum_i a_i \mathbf{1}_{C_i}$
en forma canónica,
$$ \lambda(A) = \sum_i a_i \mu(C_i \cap A) = \sum_i a_i \sum_n \mu(C_i \cap A_n) =
\sum_n \sum_i a_i \mu(C_i \cap A_n) = \sum_n \lambda(A_n), $$
donde el intercambio de sumas es lícito al ser todos los términos no negativos (teorema de
Tonelli para series).
\end{proof}

\subsection{Paso 2: funciones medibles no negativas}

\begin{definicion}{Integral de una función no negativa}{int_no_neg}
Sea $f: X \to [0, \infty]$ medible. Se define
$$ \int_X f \, d\mu = \sup\left\{ \int_X \varphi \, d\mu : \varphi \text{ simple positiva,}
\ 0 \le \varphi \le f \right\}. $$
\end{definicion}

\begin{proposicion}{Propiedades básicas en no negativas}{props_int_pos}
Para $f, g \ge 0$ medibles:
\begin{enumerate}
    \item Si $f \le g$, $\int f \, d\mu \le \int g \, d\mu$.
    \item Si $A \subseteq B$ medibles, $\int_A f \, d\mu \le \int_B f \, d\mu$.
    \item Si $A \in \mathcal{A}$, $\nu(A) := \int_A f \, d\mu$ define una medida sobre
    $(X, \mathcal{A})$.
\end{enumerate}
\end{proposicion}

\begin{proof}\leavevmode\\
\textit{(i)} Toda simple positiva $\varphi$ con $0 \le \varphi \le f$ satisface también
$0 \le \varphi \le g$, de modo que el supremo que define $\int f$ se toma sobre una subfamilia del que
define $\int g$. Por lo tanto $\int f \le \int g$.

\textit{(ii)} $f \mathbf{1}_A \le f \mathbf{1}_B$ por ser $\mathbf{1}_A \le \mathbf{1}_B$. Por (i),
$\int_A f = \int f \mathbf{1}_A \le \int f \mathbf{1}_B = \int_B f$.

\textit{(iii)} Es claro que $\nu(\emptyset) = 0$. Para $\sigma$-aditividad, sea $(A_n)$ disjunta y
$A = \bigsqcup_n A_n$.

\emph{Cota inferior $\nu(A) \ge \sum_n \nu(A_n)$.} Por (ii), para todo $N$ finito,
$\sum_{n=1}^N \nu(A_n) = \nu\!\left(\bigsqcup_{n=1}^N A_n\right) \le \nu(A)$, donde la primera
igualdad sale por aditividad finita: para una simple $\varphi \le f$, $\int \varphi
\mathbf{1}_{\bigsqcup_{n=1}^N A_n} = \sum_{n=1}^N \int \varphi \mathbf{1}_{A_n}$ por
$\sigma$-aditividad de $\mu$ aplicada al medida que define $\varphi$ (Proposición
\ref{prop:lin_mon_simple}). Tomando $N \to \infty$, $\sum_{n=1}^\infty \nu(A_n) \le \nu(A)$.

\emph{Cota superior $\nu(A) \le \sum_n \nu(A_n)$.} Sea $\varphi$ simple positiva con
$0 \le \varphi \le f \mathbf{1}_A$. Como $\varphi$ se anula fuera de $A$, $\varphi = \sum_n \varphi
\mathbf{1}_{A_n}$ (suma puntual donde a lo sumo un término es no nulo por disjuntez). Por
$\sigma$-aditividad de la medida $E \mapsto \int_E \varphi \, d\mu$ (de nuevo Proposición
\ref{prop:lin_mon_simple}),
$$ \int_A \varphi \, d\mu = \sum_n \int_{A_n} \varphi \, d\mu \le \sum_n \nu(A_n), $$
porque cada $\int_{A_n} \varphi = \int \varphi \mathbf{1}_{A_n} \le \int f \mathbf{1}_{A_n} =
\nu(A_n)$. Tomando supremo sobre $\varphi$, $\nu(A) \le \sum_n \nu(A_n)$.

Combinando ambas desigualdades, $\nu(A) = \sum_n \nu(A_n)$, es decir, $\nu$ es $\sigma$-aditiva.
\end{proof}

\subsection{Paso 3: funciones generales}

\begin{definicion}{Función integrable}{integrable}
Sea $f: X \to \mathbb{R}$ medible. Se descompone $f = f^+ - f^-$, donde $f^\pm \ge 0$. Se dice que
$f$ es \emph{integrable} si $\int f^+ d\mu < \infty$ y $\int f^- d\mu < \infty$, y se define
$$ \int_X f \, d\mu = \int_X f^+ d\mu - \int_X f^- d\mu. $$
\end{definicion}

\obs Si solo una de las dos partes es finita, la integral se define con valor en $\mathbb{R} \cup
\{\pm\infty\}$. Si ambas son infinitas, la integral no está definida.

\begin{proposicion}{Caracterización por $|f|$}{caract_integrable}
$f$ es integrable si y solo si $|f|$ es integrable, y en tal caso
$$ \left| \int_X f \, d\mu \right| \le \int_X |f| \, d\mu \qquad (\text{desigualdad triangular}). $$
Si $|f| \le g$ con $g$ integrable, entonces $f$ es integrable.
\end{proposicion}

\begin{proof}
$|f| = f^+ + f^-$, de modo que $|f|$ es integrable ssi $f^+$ y $f^-$ lo son. Para la triangular,
$|\int f| = |\int f^+ - \int f^-| \le \int f^+ + \int f^- = \int |f|$. La última afirmación
sale por monotonía.
\end{proof}

\begin{proposicion}{Linealidad y monotonía}{lin_int}
Si $f, g$ son integrables y $a, b \in \mathbb{R}$:
\begin{enumerate}
    \item $af + bg$ es integrable y $\int (af + bg) d\mu = a \int f \, d\mu + b \int g \, d\mu$.
    \item Si $f \le g$, $\int f \, d\mu \le \int g \, d\mu$.
    \item $f = 0$ c.t.p. $\implies \int f \, d\mu = 0$ (y recíproca para $f \ge 0$:
    $\int f \, d\mu = 0 \implies f = 0$ c.t.p.).
\end{enumerate}
\end{proposicion}

\begin{proof}[Demostración de (iii)]
Se prueba la recíproca: si $f \ge 0$ y $\int f \, d\mu = 0$, entonces $f = 0$ c.t.p. Para cada $m$,
el conjunto $A_m = \{f \ge 1/m\}$ cumple $f \ge \tfrac{1}{m} \mathbf{1}_{A_m}$, de donde, integrando,
$$ \frac{1}{m} \mu(A_m) \le \int_X f \, d\mu = 0. $$
Luego $\mu(A_m) = 0$ para todo $m$, y como $\{f > 0\} = \bigcup_m A_m$, se concluye
$\mu(\{f > 0\}) = 0$.
\end{proof}

\obs Para funciones con signo, $\int_X f \, d\mu = 0$ no implica $f = 0$ c.t.p.\ (tómese una
$f$ impar en $(\mathbb{R}, \lambda)$, integrable). La versión correcta pide todas las
restricciones: si
$$ \int_A f \, d\mu = 0 \qquad \text{para todo } A \in \mathcal{A}, $$
entonces $f = 0$ c.t.p. En efecto, con $A = \{f \ge 0\}$ resulta $f \cdot \mathbf{1}_A = f^+$,
de donde $\int_X f^+ \, d\mu = 0$ y $f^+ = 0$ c.t.p.\ por el resultado anterior; con
$A = \{f \le 0\}$ sale análogamente $f^- = 0$ c.t.p. Esta observación es la que da la
unicidad en el teorema de Radon-Nikodym y en la esperanza condicional.

\subsection{Desigualdades clásicas}

\begin{teorema}{Desigualdad de Markov}{markov}
Sea $f \ge 0$ medible. Para todo $a > 0$,
$$ \mu(\{f \ge a\}) \le \frac{1}{a} \int_X f \, d\mu. $$
\end{teorema}

\begin{proof}
$f \ge a \cdot \mathbf{1}_{\{f \ge a\}}$, e integrando $\int f \, d\mu \ge a \mu(\{f \ge a\})$.
\end{proof}

\begin{corolario}{Desigualdad de Chebyshev}{chebyshev}
Si $X \in L^p$ con $1 \le p < \infty$, para todo $a > 0$,
$$ \mu(\{|X| \ge a\}) \le \frac{1}{a^p} \int |X|^p \, d\mu = \frac{\|X\|_p^p}{a^p}. $$
\end{corolario}

\begin{proof}
Aplicar Markov a $f = |X|^p$ con umbral $a^p$.
\end{proof}

\obs En probabilidad: $P(|X - E(X)| \ge a) \le \mathrm{Var}(X) / a^2$ (Chebyshev con $p = 2$
aplicado a $X - E(X)$). Es la base de la ley débil de los grandes números.

\begin{teorema}{Desigualdad de Jensen (integral)}{jensen_int}
Sea $\mu$ una probabilidad ($\mu(X) = 1$), $f \in L^1(\mu)$, y $\varphi: \mathbb{R} \to \mathbb{R}$
convexa con $\varphi \circ f$ integrable. Entonces
$$ \varphi\!\left(\int_X f \, d\mu\right) \le \int_X \varphi \circ f \, d\mu. $$
\end{teorema}

\begin{proof}
Sea $c = \int_X f \, d\mu$. Como $\varphi$ es convexa, en cada punto admite una recta soporte:
existe $m \in \mathbb{R}$ tal que
$$ \varphi(y) \ge \varphi(c) + m (y - c) \qquad \forall y \in \mathbb{R}. $$
(El $m$ es cualquier valor entre las derivadas laterales $\varphi_-'(c)$ y $\varphi_+'(c)$, que
existen por convexidad.) Reemplazando $y = f(x)$ e integrando respecto de $\mu$,
$$ \int_X \varphi \circ f \, d\mu \ge \int_X \varphi(c) \, d\mu + m \int_X (f - c) \, d\mu
= \varphi(c) \mu(X) + m \left(\int f \, d\mu - c \mu(X)\right). $$
Como $\mu(X) = 1$ y $\int f \, d\mu = c$, el segundo término es cero y queda
$\int \varphi \circ f \, d\mu \ge \varphi(c) = \varphi(\int f \, d\mu)$.
\end{proof}

\obs Casos particulares útiles: $\varphi(t) = t^p$ para $p \ge 1$ da $|\int f|^p \le \int |f|^p$ (si
$\mu$ es probabilidad); $\varphi(t) = e^t$ da $\exp(\int f) \le \int e^f$; $\varphi(t) = -\ln t$
(con $t > 0$) da $\ln \int f \ge \int \ln f$.

\newpage

\section{Teoremas de convergencia}

¿Cuándo vale $\lim_n \int f_n = \int \lim_n f_n$? No siempre: sea $f_n = n \cdot
\mathbf{1}_{(0, 1/n)}$ en $(\mathbb{R}, \lambda)$. Entonces $f_n \to 0$ puntualmente (fijado
$x > 0$, para $n > 1/x$ vale $f_n(x) = 0$; en $x \le 0$ todas se anulan), pero
$$ \int_{\mathbb{R}} f_n \, d\lambda = n \cdot \frac{1}{n} = 1 \qquad \text{para todo } n, $$
de modo que $\lim_n \int f_n = 1 \ne 0 = \int \lim_n f_n$: la masa se escapa por un pico cada
vez más alto y angosto. El mismo fenómeno en versión de series: con $g_1 = f_1$ y
$g_n = f_n - f_{n-1}$ se tiene $\sum_{n=1}^m g_n = f_m$, y
$$ \sum_{n=1}^\infty \int_{\mathbb{R}} g_n \, d\lambda \ne \int_{\mathbb{R}}
\sum_{n=1}^\infty g_n \, d\lambda. $$
Los teoremas de esta sección dan condiciones (monotonía, dominación) bajo las cuales el
intercambio sí es lícito.

\begin{teorema}{Convergencia Monótona (Beppo Levi)}{mct}
Sea $(f_n)$ una sucesión creciente de funciones medibles positivas con $f_n \nearrow f$
puntualmente. Entonces
$$ \int_X f \, d\mu = \lim_{n \to \infty} \int_X f_n \, d\mu. $$
\end{teorema}

\begin{proof}
$f = \lim f_n$ es medible. Como $f_n \le f$ para todo $n$,
$$ \lim_n \int f_n \, d\mu \le \int f \, d\mu. $$
Para la otra desigualdad, sea $\varphi$ simple positiva con $0 \le \varphi \le f$ y $0 < c < 1$.
Se define $A_n = \{x : f_n(x) \ge c \varphi(x)\}$. Como $f_n \nearrow f \ge \varphi > c\varphi$
puntualmente en $\{\varphi > 0\}$, los $A_n$ crecen a $X$ (o, más precisamente, su unión cubre
$\{\varphi > 0\}$, y fuera de allí no hay nada que probar). Entonces
$$ \int f_n \, d\mu \ge \int_{A_n} f_n \, d\mu \ge c \int_{A_n} \varphi \, d\mu. $$
Como $\nu(A) := \int_A \varphi \, d\mu$ es una medida, por continuidad
$\int_{A_n} \varphi \, d\mu \to \int_X \varphi \, d\mu$. Pasando al límite,
$\lim_n \int f_n \, d\mu \ge c \int \varphi \, d\mu$. Tomando $c \to 1$ y luego sup sobre $\varphi$,
queda $\lim_n \int f_n \, d\mu \ge \int f \, d\mu$.
\end{proof}

\begin{teorema}{Lema de Fatou}{fatou}
Si $(f_n)$ es una sucesión de funciones medibles positivas,
$$ \int_X \varliminf_n f_n \, d\mu \le \varliminf_n \int_X f_n \, d\mu. $$
\end{teorema}

\begin{proof}
Se define $g_m = \inf_{n \ge m} f_n$. La sucesión $(g_m)$ es creciente y $g_m \nearrow
\varliminf_n f_n$. Por MCT,
$$ \int_X \varliminf_n f_n \, d\mu = \lim_m \int_X g_m \, d\mu. $$
Como $g_m \le f_n$ para todo $n \ge m$, $\int g_m \le \int f_n$, de modo que
$\int g_m \le \inf_{n \ge m} \int f_n$. Pasando al límite en $m$,
$\lim_m \int g_m \le \varliminf_n \int f_n$.
\end{proof}

\obs En general la desigualdad es estricta. Ejemplo: $f_n = \mathbf{1}_{[n, n+1]}$ en $\mathbb{R}$
cumple $\varliminf f_n = 0$ pero $\int f_n = 1$ para todo $n$.

\begin{teorema}{Convergencia Dominada (Lebesgue)}{dct}
Sea $(f_n)$ una sucesión de funciones medibles con $f_n \to f$ c.t.p., y supóngase que existe
una función integrable $g$ con $|f_n| \le g$ c.t.p. para todo $n$. Entonces $f$ es integrable y
$$ \int_X f \, d\mu = \lim_{n \to \infty} \int_X f_n \, d\mu. $$
Además, $\int |f_n - f| d\mu \to 0$.
\end{teorema}

\begin{proof}
Pasando a un conjunto de medida total, puede suponerse convergencia y dominación en todo punto.
Como $|f| \le g$, $f$ es integrable. Se aplica Fatou a las dos sucesiones positivas
$g + f_n \ge 0$ y $g - f_n \ge 0$:

\textit{Para $g + f_n$:}
$$ \int g + \int f \le \varliminf \int (g + f_n) = \int g + \varliminf \int f_n, $$
de modo que $\int f \le \varliminf \int f_n$.

\textit{Para $g - f_n$:}
$$ \int g - \int f \le \varliminf \int (g - f_n) = \int g - \varlimsup \int f_n, $$
de modo que $\varlimsup \int f_n \le \int f$.

Combinando, $\varlimsup \int f_n \le \int f \le \varliminf \int f_n$, lo que da que el límite
existe y vale $\int f$. La afirmación $\int|f_n - f| d\mu \to 0$ sale aplicando DCT a
$|f_n - f|$ con dominante $2g$.
\end{proof}

\begin{corolario}{Convergencia acotada}{conv_acotada}
Sea $\mu(X) < \infty$, $f_n \to f$ c.t.p.\ y $|f_n| \le M$ c.t.p.\ para todo $n$, con $M$
constante. Entonces $f$ es integrable y $\int_X f \, d\mu = \lim_n \int_X f_n \, d\mu$.
\end{corolario}

\begin{proof}
La función constante $g \equiv M$ es integrable, porque $\int_X M \, d\mu = M \mu(X) < \infty$,
y domina a todas las $f_n$. Se aplica convergencia dominada.
\end{proof}

\begin{ejemplo}{Un límite calculado con convergencia dominada}{ej_tcd}
Calcúlese
$$ L = \lim_{n \to \infty} \int_0^\pi (\operatorname{sen} x)^n \sqrt{x} \, dx. $$
El integrando converge puntualmente a $0$ en $[0, \pi] \setminus \{\pi/2\}$ (allí
$0 \le \operatorname{sen} x < 1$, de modo que $(\operatorname{sen} x)^n \to 0$), mientras que en
$x = \pi/2$ vale $\sqrt{\pi/2}$ para todo $n$: el límite puntual es $0$ c.t.p. La convergencia
no es uniforme (cerca de $\pi/2$ el integrando se mantiene próximo a $\sqrt{\pi/2}$), de modo
que el criterio clásico de intercambio no aplica. Pero
$$ |(\operatorname{sen} x)^n \sqrt{x}| \le \sqrt{x} =: g(x), \qquad \int_0^\pi \sqrt{x} \, dx =
\tfrac{2}{3} \pi^{3/2} < \infty, $$
y por convergencia dominada
$$ L = \int_0^\pi \lim_n (\operatorname{sen} x)^n \sqrt{x} \, dx = \int_0^\pi 0 \, dx = 0. $$
\end{ejemplo}

\subsection{Implicaciones entre tipos de convergencia}

Se tiene cuatro tipos de convergencia para sucesiones de funciones medibles: \emph{uniforme},
\emph{puntual} (o c.t.p.), \emph{en medida} y \emph{en $L^p$}. Véanse las implicaciones
elementales y dos resultados clásicos finos: el teorema de Riesz (subsucesión c.t.p.\ a partir
de convergencia en medida) y el DCT en medida.

\begin{proposicion}{Convergencia en $L^p$ implica en medida}{Lp_implica_medida}
Si $f_n \to f$ en $L^p$ con $1 \le p < \infty$, entonces $f_n \to f$ en medida.
\end{proposicion}

\begin{proof}
Por Chebyshev aplicada a $|f_n - f|$,
$$ \mu(\{|f_n - f| \ge \epsilon\}) \le \frac{\|f_n - f\|_p^p}{\epsilon^p} \to 0. $$
\end{proof}

\begin{teorema}{Riesz: subsucesión c.t.p.}{riesz_subsucesion}
Si $f_n \to f$ en medida, existe una subsucesión $(f_{n_k})$ con $f_{n_k} \to f$ casi todo punto.
\end{teorema}

\begin{proof}
Se construye una subsucesión $(f_{n_k})$ con
$$ \mu\!\left(\{|f_{n_k} - f| > 2^{-k}\}\right) < 2^{-k}. $$
Esto es posible porque la convergencia en medida da, para cada $k$, un $n_k$ que cumple eso, y
se puede elegir $n_k$ estrictamente crecientes.

Sea $A_k = \{|f_{n_k} - f| > 2^{-k}\}$ y $N = \varlimsup_k A_k$. Por Borel-Cantelli I (porque
$\sum \mu(A_k) < \infty$), $\mu(N) = 0$. Si $x \notin N$, $x$ está en a lo sumo finitos $A_k$,
de modo
que para $k$ grande, $|f_{n_k}(x) - f(x)| \le 2^{-k} \to 0$. Luego $f_{n_k} \to f$ fuera de $N$.
\end{proof}

\obs En general, convergencia en medida \emph{no} implica convergencia c.t.p.\ (las ``olas
caminadoras'' que se mencionaron en la sección de convergencia). Pero \emph{una subsucesión}
sí.

\begin{proposicion}{Convergencia c.t.p.\ implica en medida si $\mu(X) < \infty$}{ctp_implica_medida}
Si $\mu(X) < \infty$ y $f_n \to f$ c.t.p., entonces $f_n \to f$ en medida.
\end{proposicion}

\begin{proof}
Por Egoroff, para todo $\delta > 0$ hay $N$ con $\mu(N) < \delta$ tal que $f_n \to f$ uniformemente
en $N^c$. Dado $\epsilon > 0$, en $N^c$ existe $n_0$ tal que $|f_n - f| < \epsilon$ para
$n \ge n_0$,
o sea $\{|f_n - f| \ge \epsilon\} \subseteq N$. Luego $\mu(\{|f_n - f| \ge \epsilon\}) \le \delta$,
y $\delta$ era arbitrario.
\end{proof}

\begin{teorema}{Convergencia dominada en medida}{dct_en_medida}
Sea $(f_n)$ una sucesión medible con $f_n \to f$ en medida, y supóngase que existe $g \ge 0$
integrable con $|f_n| \le g$ c.t.p.\ para todo $n$. Entonces $f \in L^1$, $\int |f_n - f| d\mu
\to 0$, y en particular $\int f_n d\mu \to \int f \, d\mu$.

Más en general, si $|f_n|^p \le g$ con $g$ integrable y $1 \le p < \infty$, entonces
$f_n \to f$ en $L^p$.
\end{teorema}

\begin{proof}
Por el teorema de Riesz hay una subsucesión $f_{n_k} \to f$ c.t.p., de modo que $|f| \le g$ c.t.p.\
y $f \in L^1$.

Para la convergencia en $L^1$, supóngase por absurdo que $\int |f_n - f| d\mu \not\to 0$. Existe
$\epsilon_0 > 0$ y una subsucesión $(f_{n_j})$ con $\int |f_{n_j} - f| d\mu \ge \epsilon_0$. Como
$f_{n_j} \to f$ en medida, se extrae una sub-subsucesión $(f_{n_{j_l}})$ con $f_{n_{j_l}} \to f$
c.t.p. Por DCT clásico, $\int |f_{n_{j_l}} - f| d\mu \to 0$, lo que contradice la cota inferior
$\epsilon_0$.

La versión $L^p$ es análoga, aplicando DCT a $|f_n - f|^p \le 2^p (|f_n|^p + |f|^p) \le
2^p (g + |f|^p)$.
\end{proof}

\obs En resumen, en medida finita se tiene
$$ \text{uniforme} \Longrightarrow \text{c.t.p.} \Longrightarrow \text{en medida}, \qquad
L^p \Longrightarrow \text{en medida}. $$
Y con dominación integrable, ``en medida'' o ``c.t.p.'' implican ``en $L^p$''. Las demás
implicaciones son falsas en general.

\subsection{Densidad de simples y continuas en \texorpdfstring{$L^p$}{Lp}}

\begin{teorema}{Las simples son densas en $L^p$}{simples_densas}
Sea $1 \le p < \infty$. El conjunto de funciones simples integrables (es decir, simples $\varphi$
con $\|\varphi\|_p < \infty$) es denso en $L^p(\mu)$.
\end{teorema}

\begin{proof}
Sea $f \in L^p$. Se descompone $f = f^+ - f^-$. Para $f^+ \ge 0$, por el teorema de aproximación
por simples hay $\varphi_n \nearrow f^+$ con $\varphi_n$ simples positivas. Como
$0 \le \varphi_n \le f^+$, $\varphi_n \in L^p$. Además $|\varphi_n - f^+|^p \le (f^+)^p \in L^1$,
y converge a 0 puntualmente, de modo que por DCT $\|\varphi_n - f^+\|_p \to 0$. Análogo con $f^-$.
\end{proof}

\begin{teorema}{Las continuas son densas en $L^p(\mathbb{R}, \lambda)$}{continuas_densas}
Sea $1 \le p < \infty$. Las funciones continuas de soporte compacto $C_c(\mathbb{R})$ son densas
en $L^p(\mathbb{R}, \lambda)$.
\end{teorema}

\begin{proof}
Por el teorema anterior basta aproximar funciones simples integrables por continuas. Por
linealidad, basta el caso $f = \mathbf{1}_A$ con $\lambda(A) < \infty$.

\textit{Reducción a intervalos.} Por construcción de la medida de Lebesgue (vía medida
exterior con cubrimientos por intervalos), dado $\epsilon > 0$ existe una unión numerable de
intervalos abiertos $U = \bigcup_n (a_n, b_n)$ con $A \subseteq U$ y $\lambda(U \setminus A) <
\epsilon^p / 2$. Truncando a una unión finita $V = (a_1, b_1) \cup \cdots \cup (a_N, b_N)$,
puede suponerse $\lambda(U \setminus V) < \epsilon^p / 2$. Entonces $\|\mathbf{1}_A -
\mathbf{1}_V\|_p^p \le \lambda(A \triangle V) \le \lambda(U \setminus A) + \lambda(U \setminus V) <
\epsilon^p$.

\textit{Aproximación por continuas en cada intervalo.} Para $\mathbf{1}_{(a, b)}$, se define
$g_\eta: \mathbb{R} \to [0, 1]$ continua con soporte en $(a - \eta, b + \eta)$, que vale $1$ en
$[a, b]$ y baja linealmente a $0$ en las puntas. Entonces $|\mathbf{1}_{(a, b)} - g_\eta|$ se
anula fuera de $[a - \eta, a] \cup [b, b + \eta]$ y está acotado por 1, de modo que
$$ \|\mathbf{1}_{(a, b)} - g_\eta\|_p^p \le 2 \eta \to 0 \quad \text{cuando } \eta \to 0. $$

Combinando, $\|\mathbf{1}_V - \sum_k g_\eta^{(k)}\|_p$ se hace arbitrariamente pequeño, y la suma
finita de continuas de soporte compacto es continua de soporte compacto. Luego para todo
$\epsilon > 0$ hay $g \in C_c(\mathbb{R})$ con $\|\mathbf{1}_A - g\|_p < 2 \epsilon$.
\end{proof}

\obs La densidad permite probar propiedades en $L^p$ verificándolas primero en simples (o
continuas) y luego pasando al límite. Muy útil para extender operadores: si un operador es
lineal y continuo en un denso, se extiende únicamente a todo el espacio.

\newpage

\section{Los espacios \texorpdfstring{$L^p$}{Lp}}

La teoría de los espacios $L^p$ como espacios normados —las desigualdades de Young, Hölder y
Minkowski, la completitud (teorema de Riesz-Fischer), las inclusiones entre los $L^p$ y los
$\ell^p$, la ley del paralelogramo, el espacio dual y el teorema de representación de Riesz—
está desarrollada, con demostraciones, en los apuntes de \emph{Análisis}. Para no duplicarla,
aquí solo se fijan la definición y los enunciados que el resto de estos apuntes usa, y se
demuestra lo que allí no figura: la descomposición ortogonal de un espacio de Hilbert, sobre
la que se apoyan la esperanza condicional y la demostración de Radon-Nikodym.

Recuérdese que, para $1 \le p < \infty$,
$$ L^p(\mu) = \left\{ f : X \to \mathbb{R} \text{ medible} : \int_X |f|^p \, d\mu < \infty
\right\}, \qquad \|f\|_p = \left( \int_X |f|^p \, d\mu \right)^{1/p}, $$
identificando dos funciones que coinciden c.t.p.\ (sin esa identificación, $\|\cdot\|_p$ es
solo una seminorma: $\|f\|_p = 0$ equivale a $f = 0$ c.t.p.). Los hechos que se usarán, todos
demostrados en los apuntes de \emph{Análisis}:

\begin{enumerate}
    \item \emph{Desigualdad de Hölder:} si $\tfrac{1}{p} + \tfrac{1}{q} = 1$ con
    $1 < p, q < \infty$, $f \in L^p(\mu)$ y $g \in L^q(\mu)$, entonces $fg \in L^1(\mu)$ y
    $\|fg\|_1 \le \|f\|_p \, \|g\|_q$. Para $p = q = 2$ es la desigualdad de Cauchy-Schwarz.
    \item \emph{Desigualdad de Minkowski:} $\|f + g\|_p \le \|f\|_p + \|g\|_p$. Con ella,
    $(L^p(\mu), \|\cdot\|_p)$ es un espacio normado.
    \item \emph{Teorema de Riesz-Fischer:} $L^p(\mu)$ es completo, es decir, un espacio de
    Banach. (La demostración corre sobre las herramientas de estos apuntes: convergencia
    monótona, Lema de Fatou y convergencia dominada.)
    \item $L^2(\mu)$ es un \emph{espacio de Hilbert}: su norma proviene del producto interno
    $$ \langle f, g \rangle = \int_X f g \, d\mu, $$
    finito por Cauchy-Schwarz.
    \item \emph{Proyección sobre un convexo cerrado:} si $H$ es un espacio de Hilbert,
    $S \subseteq H$ es convexo, cerrado y no vacío, y $x \in H$, existe un único $s_0 \in S$
    que realiza la distancia: $\|x - s_0\| = \inf\{\|x - s\| : s \in S\}$.
    \item \emph{Representación de Riesz:} toda funcional lineal y continua
    $\varphi: H \to \mathbb{R}$ sobre un espacio de Hilbert se escribe $\varphi(x) =
    \langle x, z \rangle$ para un único $z \in H$. En $L^2(\mu)$: toda funcional lineal
    continua es ``integrar contra una función''. Esta es exactamente la forma en que la
    usará la demostración de Radon-Nikodym.
\end{enumerate}

\obs Si $\mathcal{G} \subseteq \mathcal{F}$ es una sub-$\sigma$-álgebra,
$L^p(\mathcal{G}) \subseteq L^p(\mathcal{F})$ es un subespacio y, al ser él mismo completo
(Riesz-Fischer), es un subespacio \emph{cerrado}. Este hecho lo usará la esperanza
condicional.

\begin{teorema}{Descomposición ortogonal}{descomp_ortogonal}
Sea $S$ un subespacio cerrado de un espacio de Hilbert $H$, y sea
$S^\perp = \{v \in H : \langle v, s \rangle = 0 \text{ para todo } s \in S\}$. Entonces
$H = S \oplus S^\perp$: todo $x \in H$ se escribe de manera única como $x = s_0 + v$ con
$s_0 \in S$ y $v \in S^\perp$.
\end{teorema}

\begin{proof}
Un subespacio es convexo y, por la proyección sobre un convexo cerrado (punto (v) de arriba),
existe el $s_0 \in S$ que realiza la distancia de $x$ a $S$. Se afirma que
$x - s_0 \in S^\perp$. Sean $s \in S$ y $t \in \mathbb{R}$; como $s_0 + t s \in S$,
$$ \|x - s_0\|^2 \le \|x - s_0 - t s\|^2 = \|x - s_0\|^2 - 2 t \langle x - s_0, s \rangle +
t^2 \|s\|^2. $$
La función $\phi(t) = t^2 \|s\|^2 - 2 t \langle x - s_0, s \rangle$ es entonces $\ge 0$ para
todo $t$ y se anula en $t = 0$: su derivada en $0$ debe anularse, es decir
$\langle x - s_0, s \rangle = 0$. Como $s \in S$ era arbitrario, $x - s_0 \in S^\perp$ y
$x = s_0 + (x - s_0) \in S + S^\perp$.

La descomposición es única porque $S \cap S^\perp = \{0\}$: si $w$ pertenece a ambos,
$\langle w, w \rangle = 0$ y $w = 0$.
\end{proof}

El operador $P: H \to S$, $P(x) = s_0$, se llama \emph{proyección ortogonal} sobre $S$.

\newpage

\section{Producto de medidas: Fubini-Tonelli}

Dadas dos medidas, ¿cómo se integra una función en el producto? Los teoremas de Tonelli y
Fubini permiten intercambiar el orden de integración.

\begin{definicion}{Medida producto}{medida_producto}
Sean $(X, \mathcal{A}, \mu)$ y $(Y, \mathcal{B}, \nu)$ espacios de medida $\sigma$-finitos. En el
producto $X \times Y$ se define:
\begin{itemize}
    \item La $\sigma$-álgebra producto $\mathcal{A} \otimes \mathcal{B} = \sigma(\{A \times B :
    A \in \mathcal{A}, B \in \mathcal{B}\})$.
    \item La medida producto $\mu \otimes \nu$ es la única medida sobre $\mathcal{A} \otimes
    \mathcal{B}$ con $(\mu \otimes \nu)(A \times B) = \mu(A) \nu(B)$ para rectángulos medibles.
\end{itemize}
\end{definicion}

\obs En particular, $\lambda \otimes \lambda$ en $\mathbb{R}^2$ extiende a la medida de Lebesgue
$n$-dimensional, que asigna a un rectángulo $[a, b] \times [c, d]$ el área $(b-a)(d-c)$.

\begin{teorema}{Tonelli (funciones positivas)}{tonelli}
Sea $f: X \times Y \to [0, \infty]$ medible respecto de $\mathcal{A} \otimes \mathcal{B}$, con
$\mu$ y $\nu$ $\sigma$-finitas. Entonces:
\begin{enumerate}
    \item Para cada $x \in X$, $y \mapsto f(x, y)$ es $\mathcal{B}$-medible (y simétricamente
    para $y$).
    \item $x \mapsto \int_Y f(x, y) \, d\nu(y)$ es $\mathcal{A}$-medible (y simétrico).
    \item Vale la igualdad
    $$ \int_{X \times Y} f \, d(\mu \otimes \nu) = \int_X \!\left( \int_Y f(x, y) \, d\nu(y)
    \right) d\mu(x) = \int_Y \!\left( \int_X f(x, y) \, d\mu(x) \right) d\nu(y). $$
\end{enumerate}
\end{teorema}

\begin{proof}
\textit{Paso 1: medibilidad de las secciones.} Sea $\mathcal{M} \subseteq \mathcal{A} \otimes
\mathcal{B}$ la clase de conjuntos $E$ tales que la sección $E_x = \{y : (x, y) \in E\}$ es
$\mathcal{B}$-medible para todo $x$. Si $E = A \times B$ rectángulo medible, $E_x = B$ si
$x \in A$, $\emptyset$ si no, en ambos casos medible. $\mathcal{M}$ es una $\sigma$-álgebra
(directo: $(E^c)_x = (E_x)^c$, $(\bigcup E_n)_x = \bigcup (E_n)_x$). Como contiene los
rectángulos, $\mathcal{M} \supseteq \sigma(\text{rectángulos}) = \mathcal{A} \otimes
\mathcal{B}$. Por aproximación por funciones simples, para $f$ medible $f(x, \cdot)$ también lo
es.

\textit{Paso 2: caso $f = \mathbf{1}_E$, $E \in \mathcal{A} \otimes \mathcal{B}$.} Se define
$g_E(x) = \nu(E_x)$. Hay que ver que $g_E$ es $\mathcal{A}$-medible y $\int_X g_E \, d\mu = (\mu
\otimes \nu)(E)$. Para rectángulos $E = A \times B$, $g_E(x) = \nu(B) \mathbf{1}_A(x)$, medible,
y $\int_X g_E \, d\mu = \mu(A) \nu(B) = (\mu \otimes \nu)(E)$.

La clase $\mathcal{D}$ de los $E$ donde vale esta propiedad es una \emph{clase monótona}: cerrada
por uniones crecientes y por intersecciones decrecientes (suponiendo $\sigma$-finitud para
intercambiar limites con integrales vía MCT) y por diferencias de conjuntos comparables. Por el
lema de clases monótonas (o por el lema $\pi$-$\lambda$ de Dynkin), $\mathcal{D}$ contiene a la
$\sigma$-álgebra generada por los rectángulos, que es $\mathcal{A} \otimes \mathcal{B}$.

\textit{Paso 3: caso $f$ simple positiva.} Por linealidad sobre los términos
$c_i \mathbf{1}_{E_i}$.

\textit{Paso 4: caso $f \ge 0$ general.} Se toman $\varphi_n \nearrow f$ simples positivas
(aproximación por simples). Para cada $x$, $\varphi_n(x, \cdot) \nearrow f(x, \cdot)$. Por MCT en
$Y$,
$$ \int_Y \varphi_n(x, y) \, d\nu(y) \nearrow \int_Y f(x, y) \, d\nu(y) =: F(x). $$
$F$ es límite puntual de funciones $\mathcal{A}$-medibles (por Paso 3), de modo que $F$ es
$\mathcal{A}$-medible. Aplicando MCT otra vez en $X$,
$$ \int_X F \, d\mu = \lim_n \int_X \!\int_Y \varphi_n \, d\nu \, d\mu = \lim_n \int_{X \times Y}
\varphi_n \, d(\mu \otimes \nu) = \int_{X \times Y} f \, d(\mu \otimes \nu), $$
usando el caso simple. Simétricamente para el otro orden.
\end{proof}

\begin{teorema}{Fubini (funciones integrables)}{fubini}
Sea $f: X \times Y \to \mathbb{R}$ medible. Si $f$ es integrable respecto de $\mu \otimes \nu$
(o si, por Tonelli, $\int_X \int_Y |f| \, d\nu \, d\mu < \infty$), entonces:
\begin{enumerate}
    \item Para $\mu$-c.t.p.\ $x$, la función $y \mapsto f(x, y)$ está en $L^1(\nu)$.
    \item La función $x \mapsto \int_Y f(x, y) \, d\nu(y)$, definida $\mu$-c.t.p., está en
    $L^1(\mu)$.
    \item Vale la igualdad de Tonelli para $f$, con integrales finitas.
\end{enumerate}
\end{teorema}

\begin{proof}
Se descompone $f = f^+ - f^-$ con $f^\pm \ge 0$ medibles. Por Tonelli aplicado a $|f| = f^+ + f^-$,
$$ \int_X \!\int_Y |f| \, d\nu \, d\mu = \int_{X \times Y} |f| \, d(\mu \otimes \nu) < \infty. $$

Para que el lado izquierdo sea finito, $\int_Y |f(x, y)| \, d\nu(y) < \infty$ para $\mu$-c.t.p.\
$x$, es decir, $f(x, \cdot) \in L^1(\nu)$ $\mu$-c.t.p. En esos puntos $x$ se define
$$ F(x) = \int_Y f(x, y) \, d\nu(y) = \int_Y f^+(x, y) \, d\nu(y) - \int_Y f^-(x, y) \, d\nu(y), $$
que es la diferencia de dos funciones medibles finitas $\mu$-c.t.p., de modo que $F$ es
$\mathcal{A}$-medible y, por Tonelli aplicado a $f^\pm$,
$$ \int_X F \, d\mu = \int_{X \times Y} f^+ \, d(\mu \otimes \nu) - \int_{X \times Y} f^- \,
d(\mu \otimes \nu) = \int_{X \times Y} f \, d(\mu \otimes \nu), $$
y todas las integrales son finitas porque $f \in L^1(\mu \otimes \nu)$.
\end{proof}

\obs En la práctica, el flujo es: \textit{primero} verificar con Tonelli (en $|f|$) que la
integral doble es finita, \textit{después} aplicar Fubini para intercambiar orden con $f$.
Saltarse el chequeo da resultados erróneos (la cuenta clásica
$\int_0^1 \int_0^1 (x^2 - y^2)/(x^2 + y^2)^2 \, dx \, dy = \pi/4$ vs.\ $\int \int = -\pi/4$ ocurre
porque $|f|$ no es integrable).

\newpage

\section{Espacios de medida completos}

\begin{definicion}{Espacio completo}{esp_completo}
$(X, \mathcal{A}, \mu)$ es \emph{completo} si todo subconjunto de un conjunto de medida cero es
medible (en particular, tiene medida cero):
$$ E \in \mathcal{A}, \ \mu(E) = 0, \ F \subseteq E \implies F \in \mathcal{A}. $$
\end{definicion}

\obs Todo espacio de medida puede \emph{completarse}: dado $(X, \mathcal{A}, \mu)$, se define
$\overline{\mathcal{A}} = \{A \cup N : A \in \mathcal{A}, N \subseteq E \text{ con } \mu(E) = 0\}$
y se extiende $\mu$ a $\overline{\mathcal{A}}$. Por ejemplo, $(\mathbb{R}, \mathcal{B}, \lambda)$
no es completo, pero su completación es $(\mathbb{R}, \mathcal{L}, \lambda)$.

\obs En espacios completos, si $f$ es medible y $g = f$ c.t.p., entonces $g$ también es medible.
Esto facilita el manejo de las propiedades ``c.t.p.''.

\newpage

\section{Medidas con signo y el teorema de Radon-Nikodym}

\begin{definicion}{Medida con signo}{medida_signo}
Una función $\nu: \mathcal{A} \to [-\infty, +\infty]$ es una \emph{medida con signo} si:
\begin{enumerate}
    \item $\nu(\emptyset) = 0$.
    \item $\nu$ toma a lo sumo uno de los valores $\pm\infty$.
    \item $\sigma$-aditividad: $\nu(\bigsqcup A_n) = \sum \nu(A_n)$ (con convergencia absoluta
    cuando
    el lado izquierdo es finito).
\end{enumerate}
\end{definicion}

\begin{definicion}{Absoluta continuidad}{abs_cont}
Sean $\mu, \nu$ medidas sobre $(X, \mathcal{A})$. Se dice que $\nu$ es \emph{absolutamente continua}
respecto de $\mu$ (se anota $\nu \ll \mu$) si
$$ \mu(E) = 0 \implies \nu(E) = 0 \qquad \text{para todo } E \in \mathcal{A}. $$
\end{definicion}

\begin{definicion}{Mutuamente singulares}{singular}
$\mu$ y $\nu$ son \emph{mutuamente singulares} (se anota $\mu \perp \nu$) si existen
$A, B \in \mathcal{A}$ disjuntos con $A \cup B = X$ y $\mu(B) = 0 = \nu(A)$.
\end{definicion}

\begin{ejemplo}{Medidas absolutamente continuas via integración}{}
Si $f \ge 0$ es medible y $\mu$ es una medida, $\nu(A) = \int_A f \, d\mu$ es una medida sobre el
mismo espacio medible, y es absolutamente continua respecto de $\mu$: si $\mu(A) = 0$, entonces
$\int_A f \, d\mu = 0$ (la integral sobre un conjunto $\mu$-nulo es cero). El teorema de
Radon-Nikodym, más adelante, dirá que \emph{toda} medida absolutamente continua (bajo
hipótesis razonables) tiene esta forma.
\end{ejemplo}

\subsection{Descomposición de Lebesgue + Radon-Nikodym}

El teorema que se prueba a continuación es la versión combinada: \emph{toda} medida $\nu$ se descompone en una
parte absolutamente continua y una parte singular respecto de $\mu$. Cuando $\nu \ll \mu$, la parte
singular desaparece y queda solo la derivada.

\begin{teorema}{Lebesgue-Radon-Nikodym}{radon_nikodym}
Sean $\mu$ y $\nu$ medidas $\sigma$-finitas sobre $(X, \mathcal{A})$. Entonces:
\begin{enumerate}
    \item Existen medidas únicas $\nu_a, \nu_\perp$ tales que
    $$ \nu = \nu_a + \nu_\perp, \qquad \nu_a \ll \mu, \qquad \nu_\perp \perp \mu. $$
    \item Existe una función medible $h: X \to [0, \infty)$, única $\mu$-c.t.p., tal que
    $$ \nu_a(E) = \int_E h \, d\mu \qquad \forall E \in \mathcal{A}. $$
\end{enumerate}
La función $h$ se llama la \emph{derivada de Radon-Nikodym} de $\nu_a$ respecto de $\mu$, y se
anota $h = \dfrac{d\nu_a}{d\mu}$. En particular, cuando $\nu \ll \mu$, $\nu_\perp = 0$ y
$\nu(E) = \int_E h \, d\mu$.
\end{teorema}

\begin{proof}[Demostración (von Neumann)]
Se hace el caso de medidas finitas; el paso a $\sigma$-finitas es por descomposición en partes
finitas.

\textit{Funcional auxiliar.} Sobre $L^2(\mu + \nu)$ se define
$$ \varphi(f) = \int_X f \, d\nu. $$
Es lineal y continuo: por Cauchy-Schwarz (con $g \equiv 1$),
$$ |\varphi(f)| \le \int_X |f| \, d\nu \le \nu(X)^{1/2} \, \|f\|_{L^2(\nu)} \le \nu(X)^{1/2} \,
\|f\|_{L^2(\mu+\nu)}. $$
Por el teorema de representación de Riesz aplicado a $L^2(\mu + \nu)$, existe $g \in
L^2(\mu + \nu)$ tal que
\begin{equation}\label{eq:rn1}
\int_X f \, d\nu = \int_X f g \, d(\mu + \nu) \qquad \forall f \in L^2(\mu + \nu).
\end{equation}

\textit{$0 \le g \le 1$ casi todo punto.} Tomando $f = \mathbf{1}_A$ en \eqref{eq:rn1},
$$ \nu(A) = \int_A g \, d(\mu + \nu), $$
de donde $0 \le \int_A g \, d(\mu + \nu) \le (\mu + \nu)(A)$. Esto fuerza $0 \le g \le 1$
$(\mu + \nu)$-c.t.p.

\textit{Descomposición.} Sea $N = \{x : g(x) = 1\}$ y $M = N^c = \{x : g(x) < 1\}$. Se define
$$ \nu_\perp(E) = \nu(E \cap N), \qquad \nu_a(E) = \nu(E \cap M). $$
Se prueba que $\nu_\perp \perp \mu$ y $\nu_a \ll \mu$:
\begin{itemize}
    \item $\mu(N) = 0$: tomando $f = \mathbf{1}_N$ en \eqref{eq:rn1},
    $\nu(N) = \int_N g \, d(\mu+\nu)
    = (\mu+\nu)(N) = \mu(N) + \nu(N)$, de donde $\mu(N) = 0$. Así $\nu_\perp$ y $\mu$ son
    mutuamente singulares (separadas por $N$).
    \item Si $\mu(E) = 0$, debe valer $\nu_a(E) = \nu(E \cap M) = 0$. Reorganizando \eqref{eq:rn1}
    (para $f = \mathbf{1}_{E \cap M}$),
    $$ \int_{E \cap M} (1 - g) \, d\nu = \int_{E \cap M} g \, d\mu = 0. $$
    En $M$, $1 - g > 0$, de modo que necesariamente $\nu(E \cap M) = 0$.
\end{itemize}

\textit{Derivada.} Iterando \eqref{eq:rn1} con $f_n = \mathbf{1}_E (1 + g + g^2 + \cdots + g^n)$ y
restando, se obtiene
$$ \int_E (1 - g^{n+1}) \, d\nu = \int_E (g + g^2 + \cdots + g^{n+1}) \, d\mu. $$
En $M$, $g^{n+1} \to 0$ y $g + g^2 + \cdots \to g/(1-g)$. Por convergencia monótona,
$$ \nu_a(E) = \nu(E \cap M) = \int_{E \cap M} \frac{g}{1 - g} \, d\mu = \int_E h \, d\mu, $$
donde $h = \mathbf{1}_M \cdot g/(1 - g)$.

La unicidad sale estándar: si dos descomposiciones coexistieran, su diferencia sería a la vez
absolutamente continua y singular, lo que la obliga a ser cero.
\end{proof}

\subsection{Aplicación: el Teorema Fundamental del Cálculo}

Sea $F: \mathbb{R} \to \mathbb{R}$ no decreciente y continua a derecha, y $\lambda_F$ su medida de
Lebesgue-Stieltjes. Si $\lambda_F \ll \lambda$ (cosa que pasa, por ejemplo, cuando $F$ es
\emph{absolutamente continua} en el sentido clásico), por Radon-Nikodym existe
$f \in L^1_{\mathrm{loc}}(\mathbb{R})$ tal que
$$ F(b) - F(a) = \lambda_F((a, b]) = \int_{(a, b]} f \, d\lambda = \int_a^b f(x) \, dx. $$
La función $f = d\lambda_F / d\lambda$ juega el papel de la derivada de $F$ en sentido
distribucional. Esta es la versión \emph{general} del Teorema Fundamental del Cálculo.

\newpage

\section{Esperanza Condicional}

Sea $(\Omega, \mathcal{F}, P)$ un espacio de probabilidad. Recuérdese la probabilidad condicional
elemental: dados eventos $A, B$ con $P(B) > 0$,
$$ P(A \mid B) = \frac{P(A \cap B)}{P(B)}. $$
Esta es una versión ``básica'' que solo permite condicionar a un evento. Lo que se busca ahora
es condicionar respecto a una \emph{$\sigma$-álgebra entera} $\mathcal{G} \subseteq \mathcal{F}$,
que representa la información disponible.

\subsection{Motivación}

\begin{ejemplo}{Tres tiradas de moneda}{ejemplo_monedas}
Se lanza una moneda 3 veces. $\Omega$ tiene 8 elementos $\{HHH, HHT, \dots, TTT\}$, y $\mathcal{F}
= \mathcal{P}(\Omega)$. Se define $\sigma$-álgebras anidadas que ``van revelando''
información:
\begin{align*}
\widetilde{\mathcal{F}}_0 &= \{\emptyset, \Omega\} && \text{(no se sabe nada)}, \\
\widetilde{\mathcal{F}}_1 &= \sigma(A) = \{\emptyset, \Omega, A, A^c\} && \text{donde } A
= \{\text{el 1er tiro es }H\}, \\
\widetilde{\mathcal{F}}_2 &= \sigma(B_1, B_2, B_3, B_4) && \text{(según los 2 primeros tiros)}, \\
\widetilde{\mathcal{F}}_3 &= \mathcal{P}(\Omega) && \text{(se sabe todo)}.
\end{align*}
Cada $\widetilde{\mathcal{F}}_k$ tiene más información que $\widetilde{\mathcal{F}}_{k-1}$.
\end{ejemplo}

\begin{definicion}{Filtración}{filtracion}
Una \emph{filtración} en $(\Omega, \mathcal{F}, P)$ es una familia creciente de
$\sigma$-álgebras $(\widetilde{\mathcal{F}}_t)_{t \in T} \subseteq \mathcal{F}$:
$$ s \le t \implies \widetilde{\mathcal{F}}_s \subseteq \widetilde{\mathcal{F}}_t. $$
\end{definicion}

\begin{ejemplo}{Esperanza condicional en las tres tiradas}{}
Continuando el ejemplo anterior, sea $X: \Omega \to \mathbb{R}$ con $X(\omega) = $ ``cantidad de
caras en $\omega$''. La esperanza es
$$ E(X) = 0 \cdot \tfrac{1}{8} + 1 \cdot \tfrac{3}{8} + 2 \cdot \tfrac{3}{8} + 3 \cdot \tfrac{1}{8}
= \tfrac{3}{2}. $$
Si ya se vio la 1\textsuperscript{a} tirada (es decir, se conoce $\widetilde{\mathcal{F}}_1$),
se \emph{espera} en promedio:
$$ E(X \mid \widetilde{\mathcal{F}}_1)(\omega) = \begin{cases}
2, & \text{si la 1\textsuperscript{a} fue $H$} \ (= 1 + E(\text{2 tiradas restantes})), \\
1, & \text{si la 1\textsuperscript{a} fue $T$}.
\end{cases} $$
Notar que $E(X \mid \widetilde{\mathcal{F}}_1)$ es una \emph{variable aleatoria}
$\widetilde{\mathcal{F}}_1$-medible (toma valores constantes en $A$ y en $A^c$).
\end{ejemplo}

\obs En el ejemplo puede verificarse la medibilidad a mano:
$\{E(X \mid \widetilde{\mathcal{F}}_1) \ge 2\} = A \in \widetilde{\mathcal{F}}_1$, mientras que
$\{X \ge 2\} = \{HHH, HHT, HTH, THH\} \notin \widetilde{\mathcal{F}}_1$. La esperanza
condicional es $\widetilde{\mathcal{F}}_1$-medible; $X$ no lo es.

\obs (\emph{Definición sobre particiones.}) Cuando $\mathcal{G} = \sigma(G_1, G_2, \dots)$ con
$\{G_i\}$ una partición de $\Omega$ y $P(G_i) > 0$, hay una fórmula explícita. Una variable
$\mathcal{G}$-medible es constante sobre cada átomo $G_i$, y la esperanza condicional resulta
ser el promedio de $X$ sobre el átomo en que cae $\omega$:
$$ E(X \mid \mathcal{G})(\omega) = \frac{1}{P(G_i)} \int_{G_i} X \, dP \qquad \text{si }
\omega \in G_i. $$
Se comprueba la propiedad característica: todo $A \in \mathcal{G}$ es unión de átomos,
$A = \bigcup_{i \in M} G_i$, y
$$ \int_A E(X \mid \mathcal{G}) \, dP = \sum_{i \in M} \left( \frac{1}{P(G_i)} \int_{G_i} X \,
dP \right) P(G_i) = \sum_{i \in M} \int_{G_i} X \, dP = \int_A X \, dP. $$
La definición general (que sigue) extiende esto a $\sigma$-álgebras sin átomos, donde el
promedio local ya no puede escribirse dividiendo por $P(G_i)$.

\subsection{Definición general}

\begin{definicion}{Esperanza condicional}{esp_cond}
Sea $X \in L^1(\Omega, \mathcal{F}, P)$ y sea $\mathcal{G} \subseteq \mathcal{F}$ una
$\sigma$-álgebra. La \emph{esperanza condicional} de $X$ dada $\mathcal{G}$, anotada
$E(X \mid \mathcal{G})$, es la única variable aleatoria $\mathcal{G}$-medible (módulo c.t.p.)
tal que
$$ \int_A E(X \mid \mathcal{G}) \, dP = \int_A X \, dP \qquad \forall A \in \mathcal{G}. $$
\end{definicion}

\begin{teorema}{Existencia (vía Radon-Nikodym)}{exist_esp_cond}
Bajo las hipótesis anteriores, $E(X \mid \mathcal{G})$ existe y es única $P$-c.t.p.
\end{teorema}

\begin{proof}
Supóngase primero $X \ge 0$. Se define $\lambda: \mathcal{G} \to [0, \infty]$ por
$$ \lambda(A) = \int_A X \, dP. $$
$\lambda$ es una medida sobre $(\Omega, \mathcal{G})$. Es absolutamente continua respecto de
$P|_{\mathcal{G}}$: si $P(A) = 0$, entonces $\int_A X \, dP = 0$. Por Radon-Nikodym, existe
$h \in L^1(\mathcal{G})$ con
$$ \int_A X \, dP = \lambda(A) = \int_A h \, dP \qquad \forall A \in \mathcal{G}. $$
Esa $h$ es $E(X \mid \mathcal{G})$. Para $X$ general, se descompone $X = X^+ - X^-$ y se resta.
\end{proof}

\subsection{Propiedades}

\begin{proposicion}{Propiedades de la esperanza condicional}{props_esp_cond}
Sean $X, Y \in L^1$ y $\mathcal{G}, \mathcal{H} \subseteq \mathcal{F}$ $\sigma$-álgebras.
\begin{enumerate}
    \item $E(aX + bY \mid \mathcal{G}) = a E(X \mid \mathcal{G}) + b E(Y \mid \mathcal{G})$.
    \item Si $Y$ es $\mathcal{G}$-medible (y $XY \in L^1$),
    $$ E(XY \mid \mathcal{G}) = Y \cdot E(X \mid \mathcal{G}). $$
    \item Si $X$ es independiente de $\mathcal{G}$ (es decir, $X^{-1}(B) \perp A$ para todo $A \in
    \mathcal{G}$ y boreliano $B$),
    $$ E(X \mid \mathcal{G}) = E(X). $$
    \item Si $\mathcal{H} \subseteq \mathcal{G}$,
    $$ E(E(X \mid \mathcal{G}) \mid \mathcal{H}) = E(X \mid \mathcal{H}). $$
\end{enumerate}
\end{proposicion}

\begin{proof}\leavevmode\\
\textit{(i)} Para todo $A \in \mathcal{G}$,
\begin{align*}
\int_A (aE(X \mid \mathcal{G}) + b E(Y \mid \mathcal{G})) \, dP
&= a \int_A E(X \mid \mathcal{G}) \, dP + b \int_A E(Y \mid \mathcal{G}) \, dP \\
&= a \int_A X \, dP + b \int_A Y \, dP = \int_A (aX + bY) \, dP.
\end{align*}
Por unicidad, $E(aX + bY \mid \mathcal{G}) = a E(X \mid \mathcal{G}) + b E(Y \mid \mathcal{G})$.

\textit{(ii)} Se prueba por aproximación creciente.
\begin{itemize}
    \item Para $Y = \mathbf{1}_C$ con $C \in \mathcal{G}$: para todo $A \in \mathcal{G}$,
    $A \cap C \in \mathcal{G}$, y
    $$ \int_A \mathbf{1}_C \cdot E(X \mid \mathcal{G}) dP = \int_{A \cap C} E(X \mid \mathcal{G}) dP
    = \int_{A \cap C} X dP = \int_A \mathbf{1}_C X dP. $$
    Por unicidad, $E(\mathbf{1}_C X \mid \mathcal{G}) = \mathbf{1}_C E(X \mid \mathcal{G})$.

    \item Para $Y = \sum_i c_i \mathbf{1}_{C_i}$ simple $\mathcal{G}$-medible: por linealidad de
    (i).

    \item Para $Y \ge 0$ $\mathcal{G}$-medible: se toman $Y_n \nearrow Y$ simples positivas. Si
    $X \ge 0$, por convergencia monótona $\int_A Y_n X dP \to \int_A YX dP$ y $\int_A Y_n
    E(X \mid \mathcal{G}) dP \to \int_A Y E(X \mid \mathcal{G}) dP$. Para $X$ general se descompone
    $X = X^+ - X^-$.

    \item Para $Y$ general: $Y = Y^+ - Y^-$ y se aplica el paso anterior a cada parte.
\end{itemize}

\textit{(iii)} Para $A \in \mathcal{G}$, $\mathbf{1}_A$ y $X$ son independientes,
de modo que $E(\mathbf{1}_A X) = P(A) E(X)$:
$$ \int_A X \, dP = E(\mathbf{1}_A X) = E(\mathbf{1}_A) E(X) = P(A) E(X) = \int_A E(X) \, dP. $$
$E(X)$ es constante (en particular $\mathcal{G}$-medible), de modo que por unicidad
$E(X \mid \mathcal{G}) = E(X)$.

\textit{(iv)} Sean $Z = E(X \mid \mathcal{G})$ y $W = E(X \mid \mathcal{H})$. Para
$A \in \mathcal{H} \subseteq \mathcal{G}$,
$$ \int_A Z \, dP = \int_A X \, dP = \int_A W \, dP. $$
Como $W$ es $\mathcal{H}$-medible y la igualdad vale para todo $A \in \mathcal{H}$, por unicidad de
$E(Z \mid \mathcal{H})$ se tiene $E(Z \mid \mathcal{H}) = W$, es decir $E(E(X \mid \mathcal{G})
\mid \mathcal{H}) = E(X \mid \mathcal{H})$.
\end{proof}

\subsection{Interpretación como proyección}

Cuando $X \in L^2$, la esperanza condicional admite una caracterización geométrica:

\begin{teorema}{$E(\cdot \mid \mathcal{G})$ es proyección ortogonal}{proy_L2}
Para $X \in L^2(\mathcal{F})$, $E(X \mid \mathcal{G})$ es la proyección ortogonal de $X$ sobre el
subespacio cerrado $L^2(\mathcal{G}) \subseteq L^2(\mathcal{F})$. Es decir, $E(X \mid \mathcal{G})$
es el único elemento $\mathcal{G}$-medible que minimiza
$$ \|X - Y\|_{L^2}^2 = E((X - Y)^2) \qquad \text{entre } Y \in L^2(\mathcal{G}). $$
\end{teorema}

\begin{proof}
Como $L^2(\mathcal{G})$ es subespacio cerrado de $L^2(\mathcal{F})$ (Hilbert), existe la
proyección $\widehat{X}$ caracterizada por $X - \widehat{X} \perp L^2(\mathcal{G})$, o sea
$$ \int_\Omega (X - \widehat{X}) Y \, dP = 0 \qquad \forall Y \in L^2(\mathcal{G}). $$
Tomando $Y = \mathbf{1}_A$ con $A \in \mathcal{G}$, $\int_A X \, dP = \int_A \widehat{X} \, dP$, y
como $\widehat{X}$ es $\mathcal{G}$-medible, por unicidad $\widehat{X} = E(X \mid \mathcal{G})$.
\end{proof}

\obs Esto explica el nombre ``mejor predictor de $X$ dada la información $\mathcal{G}$'':
$E(X \mid \mathcal{G})$ minimiza el error cuadrático medio entre todas las funciones
$\mathcal{G}$-medibles.

\begin{ejemplo}{Las tres tiradas como espacios anidados}{ej_L2_anidados}
En el ejemplo de las tres tiradas, $L^2(\widetilde{\mathcal{F}}_k)$ es el espacio de las
funciones constantes sobre los átomos de $\widetilde{\mathcal{F}}_k$; identificando cada
variable con el vector de sus $8$ valores $(X(HHH), X(HHT), X(HTH), X(HTT), X(THH), X(THT),
X(TTH), X(TTT))$,
$$ \underbrace{L^2(\widetilde{\mathcal{F}}_0)}_{\dim 1} \subset
\underbrace{L^2(\widetilde{\mathcal{F}}_1)}_{\dim 2} \subset
\underbrace{L^2(\widetilde{\mathcal{F}}_2)}_{\dim 4} \subset
\underbrace{L^2(\widetilde{\mathcal{F}}_3)}_{\dim 8}. $$
La variable $X$ = ``cantidad de caras'' es el vector $(3, 2, 2, 1, 2, 1, 1, 0)$, y las
sucesivas proyecciones ortogonales promedian sobre átomos cada vez más gruesos:
$$ (3,2,2,1,2,1,1,0) \ \mapsto \ \left(\tfrac{5}{2}, \tfrac{5}{2}, \tfrac{3}{2}, \tfrac{3}{2},
\tfrac{3}{2}, \tfrac{3}{2}, \tfrac{1}{2}, \tfrac{1}{2}\right) \ \mapsto \ (2,2,2,2,1,1,1,1)
\ \mapsto \ \left(\tfrac{3}{2}, \dots, \tfrac{3}{2}\right), $$
que son $E(X \mid \widetilde{\mathcal{F}}_2)$, $E(X \mid \widetilde{\mathcal{F}}_1)$ y
$E(X \mid \widetilde{\mathcal{F}}_0) = E(X) = 3/2$. Componer proyecciones
($L^2(\mathcal{F}) \to L^2(\mathcal{G}) \to L^2(\mathcal{H})$ con $\mathcal{H} \subseteq
\mathcal{G}$) es la lectura geométrica de la propiedad de torre.
\end{ejemplo}

\obs Otra caracterización, en términos de covarianzas: para $X \in L^2$,
$E(X \mid \mathcal{G})$ es la única variable $\mathcal{G}$-medible con $E(E(X \mid
\mathcal{G})) = E(X)$ tal que
$$ \mathrm{Cov}(E(X \mid \mathcal{G}), Y) = \mathrm{Cov}(X, Y) \qquad \text{para toda } Y \in
L^2(\mathcal{G}). $$
En efecto, restando,
$$ \mathrm{Cov}(X, Y) - \mathrm{Cov}(E(X \mid \mathcal{G}), Y) = \big\langle X - E(X \mid
\mathcal{G}), \, Y - E(Y) \big\rangle = 0, $$
porque $X - E(X \mid \mathcal{G})$ es ortogonal a $L^2(\mathcal{G})$ e $Y - E(Y) \in
L^2(\mathcal{G})$. La esperanza condicional conserva toda la correlación con lo observable.

\subsection{Desigualdad de Jensen condicional}

\begin{teorema}{Jensen condicional}{jensen_cond}
Sea $X \in L^1$ y $\varphi: \mathbb{R} \to \mathbb{R}$ convexa con $\varphi(X) \in L^1$. Entonces
$$ \varphi(E(X \mid \mathcal{G})) \le E(\varphi(X) \mid \mathcal{G}) \quad P\text{-c.t.p.} $$
\end{teorema}

\begin{proof}
La idea es la misma que en la Jensen integral: por convexidad existe una familia numerable de
rectas soporte $\varphi(y) = \sup_n (a_n y + b_n)$. Para cada $n$ fijo, por linealidad y monotonía
de la esperanza condicional,
$$ E(\varphi(X) \mid \mathcal{G}) \ge E(a_n X + b_n \mid \mathcal{G}) = a_n E(X \mid \mathcal{G})
+ b_n. $$
Tomando supremo sobre $n$, $E(\varphi(X) \mid \mathcal{G}) \ge \varphi(E(X \mid \mathcal{G}))$.
\end{proof}

\obs Casos particulares: $\varphi(t) = |t|^p$ con $p \ge 1$ da
$$ |E(X \mid \mathcal{G})|^p \le E(|X|^p \mid \mathcal{G}), $$
y tomando esperanza, $\|E(X \mid \mathcal{G})\|_p \le \|X\|_p$. Así, el operador esperanza
condicional es \emph{contractivo} en $L^p$.

\newpage

\section{Martingalas}

Las martingalas formalizan la idea de un \emph{juego justo}: la esperanza condicional del valor
futuro, sabiendo todo lo que pasó hasta ahora, coincide con el valor actual.

\subsection{Procesos adaptados y stopping times}

\begin{definicion}{Proceso adaptado}{adaptado}
Sea $(\widetilde{\mathcal{F}}_t)_{t \in T}$ una filtración. Un proceso estocástico
$(X_t)_{t \in T}$ es \emph{adaptado} a la filtración si cada $X_t$ es
$\widetilde{\mathcal{F}}_t$-medible.
\end{definicion}

\obs Si no hay filtración dada, se puede usar la \emph{filtración natural}
$\widetilde{\mathcal{F}}_t = \sigma(X_s : s \le t)$, que es la $\sigma$-álgebra más pequeña que
hace a $(X_s)_{s \le t}$ medibles. Claramente $(X_t)$ es adaptado a su filtración natural.

\begin{definicion}{Stopping time}{stopping}
Sea $(\widetilde{\mathcal{F}}_t)_{t \in T}$ filtración. Un \emph{stopping time} es una variable
aleatoria $\tau: \Omega \to T$ tal que para todo $t$,
$$ \{\tau \le t\} \in \widetilde{\mathcal{F}}_t. $$
En tiempo discreto $T = \mathbb{N}$ equivale a pedir $\{\tau = n\} \in \widetilde{\mathcal{F}}_n$.
\end{definicion}

\obs Intuitivamente, $\tau$ es un instante \emph{decidible con la información hasta ese momento}.
No se requiere conocimiento del futuro.

\begin{ejemplo}{Stopping times comunes}{}
Son stopping times: jugar una cantidad fija de veces ($\tau \equiv N$ constante), ``la primera vez
que aparece el evento $A$'' ($\tau = \inf\{n : \omega_n \in A\}$), o ``la $k$-ésima pasada por
$A$''. \emph{No} lo es ``el instante anterior al máximo'', porque saber que el siguiente paso
será el máximo requiere conocer el futuro. Además, $\max(\tau_1, \tau_2)$ y $\min(\tau_1,
\tau_2)$ son stopping times: $\{\max(\tau_1, \tau_2) \le t\} = \{\tau_1 \le t\} \cap \{\tau_2 \le
t\}$ y $\{\min(\tau_1, \tau_2) \le t\} = \{\tau_1 \le t\} \cup \{\tau_2 \le t\}$ están en
$\widetilde{\mathcal{F}}_t$.
\end{ejemplo}

\begin{teorema}{Lema de Wald}{wald}
Sean $(X_i)$ i.i.d.\ e integrables, $S_n = \sum_{i=1}^n X_i$, y $\tau$ un stopping time
respecto de $\widetilde{\mathcal{F}}_n = \sigma(X_1, \dots, X_n)$ con $E(\tau) < \infty$.
Entonces
$$ E(S_\tau) = E(X_1) \, E(\tau). $$
\end{teorema}

\begin{proof}
Supóngase primero $X_i \ge 0$. Se escribe
$$ S_\tau = \sum_{n=1}^\infty X_n \, \mathbf{1}_{\{\tau \ge n\}}. $$
El evento $\{\tau \ge n\} = \{\tau \le n - 1\}^c \in \widetilde{\mathcal{F}}_{n-1}$ es
independiente de $X_n$, de modo que
$E(X_n \mathbf{1}_{\{\tau \ge n\}}) = E(X_n) \, P(\tau \ge n) = E(X_1) \, P(\tau \ge n)$. Por
convergencia monótona (los sumandos son no negativos),
$$ E(S_\tau) = \sum_{n=1}^\infty E(X_n \mathbf{1}_{\{\tau \ge n\}}) = E(X_1)
\sum_{n=1}^\infty P(\tau \ge n) = E(X_1) \, E(\tau), $$
usando la identidad $E(\tau) = \sum_{n \ge 1} P(\tau \ge n)$, válida para variables a valores
en $\mathbb{N}$ (se ve escribiendo $\tau = \sum_{n \ge 1} \mathbf{1}_{\{\tau \ge n\}}$ e
intercambiando esperanza y suma por convergencia monótona). Para $X_i$ generales se aplica lo
anterior a las partes $X_i^+$ y $X_i^-$ y se resta; ambas esperanzas son finitas por
$E|X_1| < \infty$ y $E(\tau) < \infty$.
\end{proof}

\begin{ejemplo}{Dados compuestos}{ej_wald}
Se tira un dado y sale $N$; luego se tiran $N$ dados y se suma: $S = X_1 + \cdots + X_N$. Aquí
$N$ es independiente de los $X_i$ (un caso particular de stopping time), y
$$ E(S) = E(X_1) \, E(N) = \frac{7}{2} \cdot \frac{7}{2} = \frac{49}{4} = 12.25, $$
sin necesidad de condicionar a mano sobre los $6$ valores de $N$.
\end{ejemplo}

\subsection{Definición de martingala}

\begin{definicion}{Martingala}{martingala}
Un proceso $(X_n)_{n \in \mathbb{N}}$ adaptado a $(\widetilde{\mathcal{F}}_n)$ es una
\emph{martingala} si cada $X_n$ es integrable y
$$ E(X_n \mid \widetilde{\mathcal{F}}_{n-1}) = X_{n-1} \qquad \forall n \ge 1. $$
\end{definicion}

\obs Equivalentemente, $E(X_n \mid \widetilde{\mathcal{F}}_m) = X_m$ para todo $m \le n$ (por
torre).

\obs Como consecuencia, $E(X_n) = E(X_0)$ para todo $n$. Todas las variables de una martingala
tienen la misma esperanza.

\begin{ejemplo}{Constantes y promedios}{}
Si $X \in L^1$ y se define $X_n = E(X \mid \widetilde{\mathcal{F}}_n)$, entonces $(X_n)$ es
martingala respecto de $(\widetilde{\mathcal{F}}_n)$. Por torre,
$E(X_n \mid \widetilde{\mathcal{F}}_{n-1}) = E(E(X \mid \widetilde{\mathcal{F}}_n) \mid
\widetilde{\mathcal{F}}_{n-1}) = E(X \mid \widetilde{\mathcal{F}}_{n-1}) = X_{n-1}$.
\end{ejemplo}

\begin{ejemplo}{Random walk simétrico}{}
Sean $(X_n)_{n \ge 1}$ i.i.d.\ con $E(X_n) = 0$ y $S_n = \sum_{k=1}^n X_k$ (con $S_0 = 0$).
Respecto de la filtración natural,
$$ E(S_n \mid \widetilde{\mathcal{F}}_{n-1}) = E(S_{n-1} + X_n \mid \widetilde{\mathcal{F}}_{n-1})
= S_{n-1} + E(X_n) = S_{n-1}, $$
usando que $S_{n-1}$ es $\widetilde{\mathcal{F}}_{n-1}$-medible y $X_n$ es independiente. Así
$(S_n)$ es martingala.
\end{ejemplo}

\begin{ejemplo}{$S_n^2 - n$}{}
Con $X_k$ Rademacher ($\pm 1$ con probabilidad $1/2$), $(S_n^2 - n)$ es martingala. Cuenta:
\begin{align*}
E(S_n^2 \mid \widetilde{\mathcal{F}}_{n-1}) &= E((S_{n-1} + X_n)^2 \mid
\widetilde{\mathcal{F}}_{n-1})
= S_{n-1}^2 + 2 S_{n-1} E(X_n) + E(X_n^2) \\
&= S_{n-1}^2 + 0 + 1.
\end{align*}
Restando $n$ a ambos lados, $E(S_n^2 - n \mid \widetilde{\mathcal{F}}_{n-1}) = S_{n-1}^2 - (n-1)$.
\end{ejemplo}

\subsection{Estrategia previsible: transformada de martingala}

\begin{definicion}{Proceso previsible}{previsible}
$(H_n)_{n \ge 1}$ es \emph{previsible} (respecto de $(\widetilde{\mathcal{F}}_n)$) si cada $H_n$ es
$\widetilde{\mathcal{F}}_{n-1}$-medible. Intuitivamente, $H_n$ se decide \emph{antes} del paso $n$.
\end{definicion}

\begin{proposicion}{Transformada de martingala}{transformada}
Si $(X_n)$ es martingala y $(H_n)$ es previsible y acotado, entonces
$$ Y_n = \sum_{k=1}^n H_k (X_k - X_{k-1}) $$
es martingala.
\end{proposicion}

\begin{proof}
$Y_n - Y_{n-1} = H_n (X_n - X_{n-1})$. Como $H_n$ es $\widetilde{\mathcal{F}}_{n-1}$-medible,
$$ E(Y_n - Y_{n-1} \mid \widetilde{\mathcal{F}}_{n-1}) = H_n E(X_n - X_{n-1} \mid
\widetilde{\mathcal{F}}_{n-1}) = H_n \cdot 0 = 0. $$
\end{proof}

\obs La interpretación: en un juego justo, ninguna estrategia previsible (que solo usa el
pasado) puede convertir el juego en favorable. Esto justifica matemáticamente que ``no hay
sistema infalible'' en un juego justo.

\obs Una estrategia famosa: \emph{doble o nada} (martingala clásica, de donde viene el nombre).
$H_n = 2^{n-1}$ si $X_{k-1} = -k$ para todo $k < n$, y $H_n = 0$ si ya ganó. La ganancia
acumulada es 1 cuando finalmente gana, pero la ruina llega antes con probabilidad significativa.

\newpage

\section{Teoremas de Doob}

\subsection{Optional Stopping}

\begin{teorema}{Optional Stopping (versión acotada)}{optional_stopping}
Sea $(X_n)$ martingala y $\tau$ stopping time con $\tau \le M$ acotado. Entonces
$$ E(X_\tau) = E(X_0). $$
\end{teorema}

\begin{proof}
Se define $X_n^\tau = X_{\tau \wedge n}$, el proceso ``detenido'' en $\tau$. Se verifica que
$(X_n^\tau)$ es martingala (transformada por la estrategia previsible
$H_n = \mathbf{1}_{\tau \ge n}$,
que es $\widetilde{\mathcal{F}}_{n-1}$-medible porque $\{\tau \ge n\} = \{\tau \le n-1\}^c$).
Así $E(X_\tau) = E(X_M^\tau) = E(X_0^\tau) = E(X_0)$.
\end{proof}

\begin{teorema}{Optional Stopping (martingala uniformemente integrable)}{optional_stopping_UI}
Si $(X_n)$ es martingala uniformemente integrable (existe $X_\infty$ con $X_n \to X_\infty$ en
$L^1$) y $\tau$ es stopping time finito c.t.p., entonces
$$ E(X_\tau) = E(X_0). $$
\end{teorema}

\begin{proof}
Sea $\tau_n = \tau \wedge n$, que es stopping time acotado. Por la versión acotada,
$E(X_{\tau_n}) = E(X_0)$ para todo $n$.

Como $\tau < \infty$ c.t.p., $X_{\tau_n} \to X_\tau$ c.t.p.\ (eventualmente $\tau_n = \tau$). Falta
ver que $\int X_{\tau_n} dP \to \int X_\tau dP$. Se escribe
$$ X_{\tau_n} = X_\tau \mathbf{1}_{\tau \le n} + X_n \mathbf{1}_{\tau > n}. $$
Como $(X_n)$ es uniformemente integrable, en particular $\{X_n \mathbf{1}_{\tau > n}\}$ es UI (sus
valores absolutos están dominados por los de $(X_n)$), y como $P(\tau > n) \to 0$, ese término
tiende a cero en $L^1$. El primer término $X_\tau \mathbf{1}_{\tau \le n}$ converge a $X_\tau$ en
$L^1$ por convergencia dominada (con dominante $|X_\tau| \in L^1$, lo cual sale también de UI:
por Fatou y UI, $E|X_\tau| \le \varliminf E|X_{\tau_n}| < \infty$).

Tomando $n \to \infty$, $E(X_0) = E(X_{\tau_n}) \to E(X_\tau)$.
\end{proof}

\subsection{Desigualdad de Doob}

\begin{teorema}{Desigualdad maximal de Doob}{doob_ineq}
Sea $(X_n)$ martingala con $X_n \ge 0$, y $c > 0$. Entonces
$$ P\left(\max_{k \le n} X_k \ge c\right) \le \frac{E(X_n)}{c}. $$
\end{teorema}

\begin{proof}
Se define el stopping time
$$ \tau = \min\{m : \max_{k \le m} X_k \ge c\} \wedge n. $$
Por construcción $\tau \le n$ y $\{\tau < n\} = \{\max_k X_k \ge c\}$ (si el máximo alcanza $c$
antes de $n$). Por Optional Stopping,
$$ E(X_n) = E(X_\tau) = E(X_\tau \mathbf{1}_{\tau < n}) + E(X_\tau \mathbf{1}_{\tau = n})
\ge c \cdot P(\tau < n), $$
porque en $\{\tau < n\}$ vale $X_\tau \ge c$ y $X_\tau \mathbf{1}_{\tau = n} \ge 0$. Despejando,
$P(\max X_k \ge c) = P(\tau < n) \le E(X_n)/c$.
\end{proof}

\begin{corolario}{Cola gaussiana del supremo del browniano}{cola_browniano}
Sea $(W_t)$ un movimiento browniano. Para todos $a, T > 0$,
$$ P\!\left(\sup_{t \le T} W_t \ge a\right) \le e^{-a^2/(2T)}. $$
\end{corolario}

\begin{proof}
Fíjese $\lambda > 0$. Si $Z \sim \mathcal{N}(0, t)$, completando cuadrados se calcula
$E(e^{\lambda Z}) = e^{\lambda^2 t/2}$, de modo que el proceso
$$ M_t = e^{\lambda W_t - \lambda^2 t / 2} $$
es una martingala positiva: para $s < t$, como $W_t - W_s \sim \mathcal{N}(0, t - s)$ es
independiente del pasado,
$$ E(M_t \mid \mathcal{F}_s) = M_s \, e^{-\lambda^2 (t - s)/2} \, E\big(e^{\lambda (W_t -
W_s)}\big) = M_s \, e^{-\lambda^2 (t-s)/2} \, e^{\lambda^2 (t-s)/2} = M_s. $$
Como $e^x$ es creciente y $e^{-\lambda^2 t/2} \ge e^{-\lambda^2 T/2}$ para $t \le T$, vale la
inclusión de eventos
$$ \left\{ \sup_{t \le T} W_t \ge a \right\} \subseteq \left\{ \sup_{t \le T} M_t \ge
e^{\lambda a - \lambda^2 T/2} \right\}, $$
y la desigualdad maximal de Doob aplicada a la martingala positiva $(M_t)$ (la versión
continua se obtiene de la discreta tomando supremos sobre tiempos diádicos y usando la
continuidad de las trayectorias), con $E(M_T) = E(M_0) = 1$, da
$$ P\!\left(\sup_{t \le T} W_t \ge a\right) \le \frac{E(M_T)}{e^{\lambda a - \lambda^2 T/2}} =
e^{\frac{\lambda^2 T}{2} - \lambda a}. $$
El lado derecho se minimiza en $\lambda = a/T$ (anulando la derivada $\lambda T - a$ del
exponente), y al sustituir queda $e^{a^2/(2T) - a^2/T} = e^{-a^2/(2T)}$.
\end{proof}

\subsection{Convergencia de martingalas}

\begin{teorema}{Convergencia de martingalas (Doob)}{conv_martingalas}
Sea $(X_n)$ martingala con $\sup_n E|X_n| < \infty$ (acotada en $L^1$). Entonces existe
$X \in L^1$ con $X_n \to X$ casi todo punto.
\end{teorema}

\begin{proof}
Técnica de \emph{upcrossings}. Dados $a < b$, se define el número de subidas completas
$U_n^{[a,b]}(\omega)$ como la cantidad de veces que $(X_k(\omega))_{k=0}^n$ cruza desde por debajo
de $a$ hasta por encima de $b$. Más precisamente, sea $\tau_0 = 0$ y, recursivamente,
$$ \tau_{2j+1} = \min\{k > \tau_{2j} : X_k < a\}, \qquad \tau_{2j+2} = \min\{k > \tau_{2j+1} : X_k
> b\}, $$
con la convención $\min \emptyset = \infty$. Entonces $U_n^{[a,b]} = \max\{j : \tau_{2j} \le n\}$.

\textit{Desigualdad de upcrossings de Doob.} Para cualquier martingala $(X_n)$,
$$ (b - a) \, E(U_n^{[a,b]}) \le E((X_n - a)^+). $$
Idea de la prueba: la \emph{estrategia} $H_k = \mathbf{1}_{\tau_{2j+1} < k \le \tau_{2j+2} \text{
para algún } j}$ (comprar después de bajar de $a$, vender después de subir de $b$) es
previsible y acotada. La ganancia acumulada $Y_n = \sum_{k=1}^n H_k(X_k - X_{k-1})$ es
no-negativa porque cada subida completa de $a$ a $b$ aporta al menos $b - a$; más
formalmente, $Y_n \ge (b - a) U_n^{[a,b]} - (X_n - a)^+$ (el término restante absorbe una
posible subida incompleta al final). Como $Y_n$ es martingala (transformada por una previsible),
$E(Y_n) = 0$, y queda $(b - a) E(U_n) \le E((X_n - a)^+)$.

\textit{Conclusión.} Si $\sup_n E|X_n| < \infty$, entonces $E((X_n - a)^+) \le |a| + E|X_n|$
está
acotado uniformemente, de modo que
$$ E(U_\infty^{[a,b]}) = \lim_n E(U_n^{[a,b]}) \le \frac{|a| + \sup_n E|X_n|}{b - a} < \infty. $$
En particular, $U_\infty^{[a,b]} < \infty$ casi todo punto.

Sea $\Lambda_{a,b} = \{\omega : \varliminf X_n(\omega) < a < b < \varlimsup X_n(\omega)\}$.
En $\Lambda_{a,b}$, la sucesión $X_n(\omega)$ cruza infinitas veces el intervalo $[a, b]$, de modo
que $U_\infty^{[a,b]}(\omega) = \infty$. Luego $P(\Lambda_{a,b}) = 0$.

Tomando unión numerable sobre racionales $a < b$,
$$ P\!\left(\bigcup_{a < b \in \mathbb{Q}} \Lambda_{a,b}\right) = 0. $$
Fuera de ese conjunto nulo, $\varliminf X_n = \varlimsup X_n$, es decir el límite
$X = \lim X_n$ existe en $\mathbb{R} \cup \{\pm\infty\}$. Por el Lema de Fatou aplicado a $|X_n|$,
$E|X| \le \varliminf E|X_n| \le \sup E|X_n| < \infty$, de modo que $X \in L^1$ (en particular,
finito c.t.p.).
\end{proof}

\obs La convergencia del teorema es c.t.p., \emph{no} en $L^1$. Contraejemplo: sean $(Z_k)$
i.i.d.\ con $P(Z_k = 0) = P(Z_k = 2) = 1/2$ (así $E(Z_k) = 1$) y $M_n = \prod_{k=1}^n Z_k$.
Entonces $(M_n)$ es martingala respecto de $\widetilde{\mathcal{F}}_n = \sigma(Z_1, \dots,
Z_n)$:
$$ E(M_{n+1} \mid \widetilde{\mathcal{F}}_n) = M_n \, E(Z_{n+1} \mid
\widetilde{\mathcal{F}}_n) = M_n \, E(Z_{n+1}) = M_n, $$
y $E(M_n) = 1$ para todo $n$, con $\sup_n E|M_n| = 1 < \infty$. Por el teorema, $M_n$ converge
c.t.p., y el límite es $0$: en cuanto sale un $Z_k = 0$ (lo que ocurre con probabilidad $1$),
el producto se anula para siempre. Pero entonces
$$ E|M_n - 0| = E(M_n) = 1 \not\to 0: $$
no hay convergencia en $L^1$, y la esperanza se pierde en el límite
($1 = \lim_n E(M_n) \ne E(\lim_n M_n) = 0$).

\begin{teorema}{Convergencia $L^p$, $p > 1$}{conv_Lp_martingalas}
Si $\sup_n E|X_n|^p < \infty$ con $p > 1$, entonces $X_n \to X$ también en $L^p$, y
$X_n = E(X \mid \widetilde{\mathcal{F}}_n)$.
\end{teorema}

\begin{proof}
Como $L^p \subseteq L^1$ en medida finita ($P$ es probabilidad), $\sup E|X_n| < \infty$ y por el
teorema anterior $X_n \to X$ c.t.p.\ con $X \in L^1$.

\textit{Acotación uniforme en $L^p$ via desigualdad de Doob.} Aplicando la desigualdad maximal de
Doob a la martingala positiva $|X_n|^p$ (que es submartingala por Jensen condicional con
$\varphi(t) = |t|^p$), se obtiene la \emph{desigualdad maximal $L^p$}:
$$ E\!\left(\sup_{k \le n} |X_k|^p\right) \le \left(\frac{p}{p - 1}\right)^p \sup_{k \le n}
E|X_k|^p. $$
Tomando $n \to \infty$ y por convergencia monótona, $E(\sup_n |X_n|^p) \le (p/(p-1))^p \sup_n
E|X_n|^p < \infty$. Llámese $M = \sup_n |X_n|$, que cumple $M \in L^p$.

\textit{Convergencia en $L^p$.} Se tiene $|X_n - X|^p \le 2^p (|X_n|^p + |X|^p) \le 2^p (M^p + M^p)
= 2^{p+1} M^p$ (porque $|X| \le M$ por el límite). Como $X_n \to X$ c.t.p.\ y $|X_n - X|^p \le
2^{p+1} M^p \in L^1$, por convergencia dominada $E|X_n - X|^p \to 0$.

\textit{Identidad $X_n = E(X \mid \widetilde{\mathcal{F}}_n)$.} Para $A \in
\widetilde{\mathcal{F}}_n$ y $m \ge n$, por la propiedad de martingala $\int_A X_n dP = \int_A X_m
dP$. Por convergencia en $L^1$, $\int_A X_m dP \to \int_A X dP$. Luego $\int_A X_n dP = \int_A X
dP$, y como $X_n$ es $\widetilde{\mathcal{F}}_n$-medible, por unicidad $X_n = E(X \mid
\widetilde{\mathcal{F}}_n)$.
\end{proof}

\newpage

\section{Movimiento Browniano}

El paso del tiempo discreto al continuo lleva al \emph{movimiento browniano}, límite de los
random walks reescalados. Antes de la definición axiomática se estudia el random walk con
cierto detalle: sus cuentas explícitas son las que motivan las propiedades del proceso
límite.

\subsection{El random walk simple}

Sean $(X_n)_{n \ge 1}$ i.i.d.\ con $P(X_n = 1) = P(X_n = -1) = 1/2$, y $S_n = X_1 + \cdots +
X_n$ (con $S_0 = 0$): una partícula parte del origen y a cada paso se mueve una unidad a la
derecha o a la izquierda según una moneda equilibrada. Las cuentas básicas: $E(X_i) = 0$ y
$\mathrm{Var}(X_i) = E(X_i^2) = 1$, de donde
$$ E(S_n) = 0, \qquad \mathrm{Var}(S_n) = E(S_n^2) = \sum_{i,j} E(X_i X_j) = \sum_{i=1}^n
E(X_i^2) = n, $$
usando que para $i \ne j$, $E(X_i X_j) = E(X_i) E(X_j) = 0$ por independencia. Además, para
$m < n$, $S_n - S_m = \sum_{i = m+1}^n X_i$ es otro random walk, independiente de
$S_1, \dots, S_m$, con $\mathrm{Var}(S_n - S_m) = n - m$.

\begin{ejemplo}{Distribución de $S_n$ y retornos al origen}{ej_retorno_origen}
Tras $n$ pasos, $S_n = a$ exige $k$ pasos $+1$ y $n - k$ pasos $-1$ con $a = k - (n - k) =
2k - n$; en particular $a$ y $n$ tienen la misma paridad, y $k = (n + a)/2$. De los $2^n$
caminos equiprobables, los favorables son las elecciones de dónde colocar los pasos $+1$:
$$ P(S_n = a) = \frac{1}{2^n} \binom{n}{\frac{n+a}{2}}. $$
En particular, la probabilidad de estar en el origen tras $2m$ pasos es $P(S_{2m} = 0) =
\binom{2m}{m} / 2^{2m}$, que con la fórmula de Stirling $n! \approx n^n e^{-n} \sqrt{2 \pi n}$
se estima
$$ P(S_{2m} = 0) = \frac{(2m)!}{(m!)^2 \, 2^{2m}} \approx
\frac{(2m)^{2m} e^{-2m} \sqrt{4 \pi m}}{m^{2m} e^{-m} \cdot m^{m} e^{-m} \cdot (2 \pi m) \cdot
2^{2m}} \cdot \frac{m^{m}}{m^{m}} = \frac{\sqrt{4 \pi m}}{2 \pi m} = \frac{1}{\sqrt{\pi m}}. $$
Decae, pero muy lentamente: si $V = \sum_{m \ge 1} \mathbf{1}_{\{S_{2m} = 0\}}$ es la cantidad
de retornos al origen,
$$ E(V) = \sum_{m=1}^\infty P(S_{2m} = 0) \approx \sum_{m=1}^\infty \frac{1}{\sqrt{\pi m}} =
\infty. $$
(Borel-Cantelli (II) no es aplicable aquí: los eventos $\{S_{2m} = 0\}$ no son
independientes.)
\end{ejemplo}

\obs $E(V) = \infty$ no implica por sí solo $V = \infty$ c.t.p.: una variable puede tener
esperanza infinita y ser finita con probabilidad $1$. Es la \emph{paradoja de San
Petersburgo}: se tira una moneda hasta la primera ceca, y si antes salieron $k$ caras se ganan
$2^k$ pesos; la ganancia esperada es
$$ E(G) = \sum_{k=0}^\infty 2^k \cdot \frac{1}{2^{k+1}} = \sum_{k=0}^\infty \frac{1}{2} =
\infty, $$
pero $G$ es finita con probabilidad $1$. Para el random walk el argumento correcto es otro. Sea
$\rho$ la probabilidad de volver al origen alguna vez. Cada vez que el paseo retorna, el
proceso vuelve a empezar (los pasos posteriores son independientes y con la misma
distribución), de modo que $P(V \ge k) = \rho^k$. Si fuera $\rho < 1$,
$$ E(V) = \sum_{k \ge 1} P(V \ge k) = \sum_{k \ge 1} \rho^k = \frac{\rho}{1 - \rho} < \infty, $$
contradiciendo $E(V) = \infty$. Por lo tanto $\rho = 1$: el random walk vuelve al origen — y,
iterando, lo visita \emph{infinitas veces} — con probabilidad $1$.

\begin{ejemplo}{La tabla de Galton y la campana de Gauss}{ej_galton}
En una tabla de Galton caen bolitas por $n$ niveles de clavos; cada bolita realiza un random
walk de $n$ pasos, y la columna sobre la posición $a$ acumula una fracción de bolitas
proporcional a $\binom{n}{(n+a)/2}$. Midiendo cada columna en unidades de la columna central,
la altura sobre $a$ es
$$ H_a = \binom{n}{\frac{n+a}{2}} \Big/ \binom{n}{\frac{n}{2}} =
\frac{(n/2)! \, (n/2)!}{\left(\frac{n+a}{2}\right)! \left(\frac{n-a}{2}\right)!}. $$
Escríbase $\rho = n/2$ y $c = a/2$; cancelando factoriales y tomando logaritmo,
$$ H_a = \frac{(\rho - c + 1) \cdots \rho}{(\rho + 1) \cdots (\rho + c)}, \quad
\ln H_a = \sum_{k = \rho - c + 1}^{\rho} \ln k - \sum_{k = \rho + 1}^{\rho + c} \ln k
\approx \int_{\rho - c}^{\rho} \ln x \, dx - \int_{\rho}^{\rho + c} \ln x \, dx. $$
Con la primitiva $x \ln x - x$,
$$ \ln H_a \approx 2 \rho \ln \rho - (\rho - c) \ln(\rho - c) - (\rho + c) \ln(\rho + c). $$
El polinomio de Taylor de orden $2$ de $f(x) = x \ln x$ en $\rho$ (con $f'(x) = \ln x + 1$ y
$f''(x) = 1/x$) da $f(\rho \pm c) \approx \rho \ln \rho \pm (\ln \rho + 1) c +
\frac{c^2}{2\rho}$; al sustituir, los términos de orden $0$ y $1$ se cancelan y queda
$$ \ln H_a \approx -\frac{c^2}{\rho} = -\frac{a^2/4}{n/2} = -\frac{a^2}{2n}
= -\frac{1}{2} \left( \frac{a}{\sqrt{n}} \right)^{\!2}, \quad \text{es decir} \quad
H_a \approx e^{-x^2/2} \quad \text{con } x = \frac{a}{\sqrt{n}}. $$
Reescalando el eje horizontal por $\sqrt{n}$, los perfiles se ajustan a la campana de Gauss:
es el teorema de \emph{de Moivre-Laplace},
$$ \frac{S_n}{\sqrt{n}} \xrightarrow{d} \mathcal{N}(0, 1), \qquad
P\!\left( \alpha \le \frac{S_n}{\sqrt{n}} \le \beta \right) \approx \frac{1}{\sqrt{2 \pi}}
\int_\alpha^\beta e^{-x^2/2} \, dx. $$
\end{ejemplo}

\subsection{Del random walk al browniano}

Heurísticamente, el browniano es un random walk con pasos temporales de longitud
$\Delta t = 1/N$ y pasos espaciales de longitud $\epsilon$, en el límite $N \to \infty$: se
define $W_{n/N} = \epsilon S_n$ y se normaliza pidiendo $\mathrm{Var}(W_1) = 1$. Como para
llegar a $t = 1$ se dan $N$ pasos,
$$ 1 = E(W_1^2) = \epsilon^2 E(S_N^2) = \epsilon^2 N \implies \epsilon = \frac{1}{\sqrt{N}}. $$
Con esa elección, para $t = n/N$,
$$ W_t = \frac{S_n}{\sqrt{N}} = \frac{S_n}{\sqrt{n}} \cdot \sqrt{\frac{n}{N}} =
\frac{S_n}{\sqrt{n}} \cdot \sqrt{t}, $$
y como $S_n / \sqrt{n} \sim \mathcal{N}(0, 1)$ para $n$ grande (de Moivre-Laplace), se espera
$W_t \sim \mathcal{N}(0, t)$. De $\mathrm{Var}(S_n - S_m) = n - m$ y de la independencia de los
tramos del random walk se esperan también $W_t - W_s \sim \mathcal{N}(0, t - s)$ e incrementos
independientes. Estas propiedades, junto con la continuidad de las trayectorias, son
exactamente la definición.

\subsection{Definición y propiedades}

\begin{definicion}{Movimiento Browniano}{browniano}
Un \emph{movimiento browniano} (proceso de Wiener) es un proceso estocástico
$(W_t)_{t \ge 0}$ que cumple:
\begin{enumerate}
    \item $W_0 = 0$ casi todo punto.
    \item Para $0 \le s < t$, $W_t - W_s \sim \mathcal{N}(0, t - s)$.
    \item Incrementos independientes: para $0 \le t_0 < t_1 < \cdots < t_n$, las variables
    $W_{t_1} - W_{t_0}, W_{t_2} - W_{t_1}, \dots, W_{t_n} - W_{t_{n-1}}$ son independientes.
    \item $t \mapsto W_t$ es continua casi todo punto.
\end{enumerate}
\end{definicion}

\obs Existencia (Wiener, 1923): tal proceso existe. La construcción usa, por ejemplo,
extensión de Kolmogorov o expansión en serie aleatoria.

\obs Las trayectorias $t \mapsto W_t(\omega)$ son continuas pero, casi seguramente, no son
derivables en ningún punto. La heurística de ambas cosas está en la escala de los
incrementos: como $\mathrm{Var}(W_{t+h} - W_t) = h$, el incremento típico es
$|W_{t+h} - W_t| \approx \sqrt{h}$, que tiende a $0$ (continuidad) pero cuyo cociente
incremental $\sqrt{h}/h = 1/\sqrt{h} \to \infty$ (no derivabilidad). Más aún: una derivada en
$t_0$ podría estimarse con incrementos pasados ($r < t_0$) y determinaría los incrementos
futuros próximos, pero $W_s - W_{t_0}$ es \emph{independiente} de $\{W_r : r \le t_0\}$: el
proceso no puede recordar una pendiente. Esto es lo que hace que el cálculo clásico falle y
obligue a la construcción de la integral de Itô.

\obs Transformaciones que preservan el browniano: si $(W_t)$ es browniano, también lo son
\begin{gather*}
(c \, W_{t/c^2})_{t \ge 0} \quad \text{(autosemejanza)}, \qquad
(W_{t+s} - W_s)_{t \ge 0} \quad \text{(reinicio)}, \\
(t \, W_{1/t})_{t \ge 0} \quad \text{(inversión temporal)}.
\end{gather*}
Con drift y escala, $(\mu t + \sigma W_t)$ se llama \emph{browniano con drift} y
$(C e^{\mu t + \sigma W_t})$ \emph{browniano geométrico}.

\subsection{Variación cuadrática}

\begin{teorema}{Variación cuadrática del browniano}{var_cuad}
Sea $a = t_0 < t_1 < \cdots < t_m = b$ partición de $[a, b]$ con norma
$\delta = \max_k (t_{k+1} - t_k)$. Se define
$$ Q_m = \sum_{k=1}^m (W_{t_k} - W_{t_{k-1}})^2. $$
Entonces $Q_m \xrightarrow{L^2} b - a$ cuando $\delta \to 0$.
\end{teorema}

\begin{proof}
Llámese $\Delta_k = W_{t_k} - W_{t_{k-1}}$ y $\Delta t_k = t_k - t_{k-1}$. Por definición de
browniano, $\Delta_k \sim \mathcal{N}(0, \Delta t_k)$ y los $\Delta_k$ son independientes.

\textit{Esperanza.} $E(\Delta_k^2) = \mathrm{Var}(\Delta_k) = \Delta t_k$, de modo que
$$ E(Q_m) = \sum_{k=1}^m E(\Delta_k^2) = \sum_{k=1}^m \Delta t_k = b - a. $$

\textit{Varianza.} Por independencia de los $\Delta_k^2$,
$$ \mathrm{Var}(Q_m) = \sum_{k=1}^m \mathrm{Var}(\Delta_k^2). $$
Calcúlese $\mathrm{Var}(\Delta_k^2)$: si $Z \sim \mathcal{N}(0, \sigma^2)$, entonces
$Z^2 = \sigma^2
\xi$ con $\xi \sim \chi^2_1$. Como $E(\xi^2) = 3$ (cuarto momento de la normal estándar) y
$E(\xi) = 1$, $\mathrm{Var}(\xi) = 2$. Luego $\mathrm{Var}(Z^2) = \sigma^4 \cdot 2 = 2 \sigma^4$, y
en este caso $\mathrm{Var}(\Delta_k^2) = 2 (\Delta t_k)^2$.

Sustituyendo,
$$ \mathrm{Var}(Q_m) = 2 \sum_{k=1}^m (\Delta t_k)^2 \le 2 \delta \sum_{k=1}^m \Delta t_k =
2 \delta (b - a) \xrightarrow{\delta \to 0} 0. $$

\textit{Conclusión.} $E((Q_m - (b - a))^2) = \mathrm{Var}(Q_m) \to 0$, es decir $Q_m \to b - a$
en $L^2$.
\end{proof}

\obs Heurísticamente, esto se escribe como $(dW)^2 = dt$. También valen, en el límite,
$$ \sum (\Delta t)^2 \to 0, \qquad \sum \Delta t \, \Delta W \to 0, \qquad \sum (\Delta W)^2 \to b
- a, $$
que son las reglas que se usan al construir la integral de Itô y el lema correspondiente.

\newpage

\section{Integral de Itô}

Se querría definir $\int_a^b X_t \, dW_t$ siguiendo Riemann: aproximar por sumas
$\sum X_{\xi_k} (W_{t_{k+1}} - W_{t_k})$ con $\xi_k \in [t_k, t_{k+1}]$. Pero a diferencia del
caso clásico, la elección de $\xi_k$ \emph{sí} importa, porque las trayectorias del browniano
no tienen variación acotada. La elección que da una integral con buenas propiedades martingala
es la del \emph{extremo izquierdo}: $\xi_k = t_k$.

\begin{definicion}{Integral de Itô}{ito_integral}
Sea $(X_t)_{t \in [a, b]}$ un proceso adaptado con $X_\cdot \in L^2([a, b] \times \Omega)$
(suficientemente integrable).
Para una partición $a = t_0 < t_1 < \cdots < t_m = b$ se define la \emph{suma de Riemann por el
extremo izquierdo}
$$ R_m = \sum_{k=1}^m X_{t_{k-1}} (W_{t_k} - W_{t_{k-1}}). $$
La integral de Itô se define como
$$ \int_a^b X \, dW = \lim_{\delta \to 0} R_m, $$
en $L^2(\Omega)$. El límite existe y no depende de la sucesión de particiones.
\end{definicion}

\obs La elección del extremo izquierdo es esencial. Tomando $\xi_k = t_k$ (extremo derecho) o
$\xi_k = (t_{k-1} + t_k)/2$ (punto medio, Stratonovich) se obtienen integrales \emph{distintas}.
La de Itô es la única que conserva la propiedad de ser martingala. La diferencia se
calcula directamente para $\int_a^b W \, dW$: si $B_m$ usa el extremo derecho y $A_m$ el
izquierdo,
$$ B_m - A_m = \sum_k \left( W_{t_k} - W_{t_{k-1}} \right) \left( W_{t_k} - W_{t_{k-1}} \right)
= Q_m \xrightarrow{L^2} b - a \ne 0 $$
por la variación cuadrática: ambos límites existen, pero difieren exactamente en $b - a$.

\subsection{Propiedades de la integral de Itô}

\begin{proposicion}{Propiedades}{props_ito}
Sean $X, Y$ procesos adaptados integrables en sentido de Itô, $\alpha, \beta \in \mathbb{R}$:
\begin{enumerate}
    \item $\int_a^b (\alpha X + \beta Y) dW = \alpha \int_a^b X \, dW + \beta \int_a^b Y \, dW$.
    \item $\displaystyle E\!\left[\int_a^b X \, dW\right] = 0$.
    \item $\displaystyle E\!\left[\left(\int_a^b X \, dW\right)^2 \right] = E\!\left[\int_a^b X^2
    \, dt\right]$.
    \item $\displaystyle E\!\left[\int_a^b X \, dW \cdot \int_a^b Y \, dW\right] =
    E\!\left[\int_a^b XY \, dt\right]$.
\end{enumerate}
\end{proposicion}

\begin{proof}
Se prueban (ii) y (iii) al nivel de las sumas de Riemann por el extremo izquierdo (el caso
general sale pasando al límite en $L^2$). Escríbase $R = \sum_k X_{t_{k-1}} \Delta W_k$ con
$\Delta W_k = W_{t_k} - W_{t_{k-1}}$ y $\Delta t_k = t_k - t_{k-1}$.

\textit{(ii)} Como $X_{t_{k-1}}$ es $\mathcal{F}_{t_{k-1}}$-medible (el proceso es adaptado) y
$\Delta W_k$ es independiente de $\mathcal{F}_{t_{k-1}}$ con media $0$,
$$ E(X_{t_{k-1}} \Delta W_k) = E(X_{t_{k-1}}) \, E(\Delta W_k) = 0, $$
y sumando, $E(R) = 0$.

\textit{(iii)} Desarrollando el cuadrado,
$$ E(R^2) = \sum_k E\big(X_{t_{k-1}}^2 (\Delta W_k)^2\big) + 2 \sum_{j < k}
E\big(X_{t_{j-1}} \Delta W_j \, X_{t_{k-1}} \Delta W_k\big). $$
En los términos cruzados con $j < k$, el factor $X_{t_{j-1}} \Delta W_j X_{t_{k-1}}$ es
$\mathcal{F}_{t_{k-1}}$-medible, y $\Delta W_k$ es independiente de él y centrado: cada
término cruzado se anula. En los diagonales, $(\Delta W_k)^2$ es independiente de
$X_{t_{k-1}}^2$ y $E((\Delta W_k)^2) = \Delta t_k$, de donde
$$ E(R^2) = \sum_k E(X_{t_{k-1}}^2) \, \Delta t_k = E\!\left( \sum_k X_{t_{k-1}}^2 \, \Delta
t_k \right) \xrightarrow{\delta \to 0} E\!\left( \int_a^b X^2 \, dt \right). $$
La propiedad (i) es la linealidad de las sumas de Riemann, y (iv) se deduce de (iii) por
polarización, $XY = \tfrac{1}{4} \left[ (X + Y)^2 - (X - Y)^2 \right]$.
\end{proof}

\obs La isometría de Itô es el cambio formal $(dW)^2 \mapsto dt$ que se mencionó: el cuadrado
de la integral estocástica se traduce en la integral clásica de $X^2$ respecto del tiempo.

\subsection{Lema de Itô}

El cambio de variables en cálculo estocástico no es el clásico: aparece un término
adicional debido a $(dW)^2 = dt$.

\begin{teorema}{Lema de Itô}{lema_ito}
Sea $g \in C^2(\mathbb{R})$. Si $Z_t = g(W_t)$,
$$ dZ_t = g'(W_t) \, dW_t + \tfrac{1}{2} g''(W_t) \, dt. $$
En forma integral:
$$ g(W_b) - g(W_a) = \int_a^b g'(W_s) \, dW_s + \frac{1}{2} \int_a^b g''(W_s) \, ds. $$
\end{teorema}

\begin{proof}
Fíjese una partición $a = t_0 < t_1 < \cdots < t_m = b$ con norma $\delta = \max_k (t_k -
t_{k-1})$. Se escribe $\Delta W_k = W_{t_k} - W_{t_{k-1}}$ y $\Delta t_k = t_k - t_{k-1}$.
Por suma telescópica,
$$ g(W_b) - g(W_a) = \sum_{k=1}^m [g(W_{t_k}) - g(W_{t_{k-1}})]. $$

\textit{Taylor de orden 2.} Como $g \in C^2$, para cada $k$ existe $\xi_k$ entre $W_{t_{k-1}}$ y
$W_{t_k}$ con
$$ g(W_{t_k}) - g(W_{t_{k-1}}) = g'(W_{t_{k-1}}) \Delta W_k + \tfrac{1}{2} g''(\xi_k) (\Delta
W_k)^2. $$
Sumando sobre $k$, la suma se separa en dos términos:
$$ g(W_b) - g(W_a) = \underbrace{\sum_{k=1}^m g'(W_{t_{k-1}}) \Delta W_k}_{I_m} +
\underbrace{\tfrac{1}{2} \sum_{k=1}^m g''(\xi_k) (\Delta W_k)^2}_{J_m}. $$

\textit{Término $I_m$.} Por la construcción de la integral de Itô como límite $L^2$ de
sumas de Riemann por el extremo izquierdo,
$$ I_m \xrightarrow{L^2(\Omega)}{} \int_a^b g'(W_s) \, dW_s. $$

\textit{Término $J_m$.} Se lo separa en dos:
$$ J_m = \tfrac{1}{2} \sum_k g''(W_{t_{k-1}}) (\Delta W_k)^2 + \tfrac{1}{2} \sum_k [g''(\xi_k) -
g''(W_{t_{k-1}})] (\Delta W_k)^2 =: J_m^{(1)} + J_m^{(2)}. $$

Para $J_m^{(1)}$, se usa la variación cuadrática: dado $\epsilon > 0$, particiones finas hacen
que $\sum_k (\Delta W_k)^2 \to b - a$ en $L^2$. Más aún, para una función continua
$h$, vale la convergencia
$$ \sum_k h(W_{t_{k-1}}) (\Delta W_k)^2 \xrightarrow{L^2}{} \int_a^b h(W_s) \, ds. $$
(Este resultado es el análogo del teorema fundamental para la variación cuadrática.)
Aplicando con $h = g''$,
$$ J_m^{(1)} \xrightarrow{L^2}{} \tfrac{1}{2} \int_a^b g''(W_s) \, ds. $$

Para $J_m^{(2)}$: como $g''$ es continua y los $\xi_k$ están entre $W_{t_{k-1}}$ y $W_{t_k}$ con
$|W_{t_k} - W_{t_{k-1}}| \to 0$ por continuidad uniforme del browniano en compactos,
$\max_k |g''(\xi_k) - g''(W_{t_{k-1}})| \to 0$ casi todo punto. Como $\sum (\Delta W_k)^2$ está
acotado en $L^1$ (su esperanza es $b - a$),
$$ |J_m^{(2)}| \le \tfrac{1}{2} \max_k |g''(\xi_k) - g''(W_{t_{k-1}})| \cdot \sum_k (\Delta W_k)^2
\xrightarrow{}{} 0. $$

\textit{Conclusión.} Tomando $\delta \to 0$ en
$g(W_b) - g(W_a) = I_m + J_m^{(1)} + J_m^{(2)}$, queda
$$ g(W_b) - g(W_a) = \int_a^b g'(W_s) \, dW_s + \tfrac{1}{2} \int_a^b g''(W_s) \, ds. $$
\end{proof}

\textbf{Versión con $t$.} Si $g \in C^{1, 2}(\mathbb{R}_{\ge 0} \times \mathbb{R})$ y
$Z_t = g(t, W_t)$,
$$ dZ_t = \frac{\partial g}{\partial t} dt + \frac{\partial g}{\partial x} dW_t + \frac{1}{2}
\frac{\partial^2 g}{\partial x^2} dt. $$

\subsection{Ejemplos}

\begin{ejemplo}{Cuadrado del browniano}{cuadrado_W}
Aplicando Itô a $g(x) = x^2$, $g'(x) = 2x$, $g''(x) = 2$. Resulta
$$ W_b^2 - W_a^2 = \int_a^b 2 W_s \, dW_s + \int_a^b 1 \, ds, $$
de donde $\int_a^b W \, dW = \tfrac{W_b^2 - W_a^2}{2} - \tfrac{b - a}{2}$. Notar el término
$-(b-a)/2$, ausente en el cálculo clásico.
\end{ejemplo}

\begin{ejemplo}{Browniano geométrico}{browniano_geo}
Resuélvase la ecuación diferencial estocástica
$$ dX_t = \mu X_t \, dt + \sigma X_t \, dW_t, \qquad X_0 \text{ dado}. $$
Se toma $Y_t = \ln X_t$ y se aplica Itô con $g(x) = \ln x$, $g'(x) = 1/x$, $g''(x) = -1/x^2$:
\begin{align*}
d(\ln X_t) &= \frac{1}{X_t} dX_t - \frac{1}{2 X_t^2} (dX_t)^2 \\
&= \frac{1}{X_t} (\mu X_t \, dt + \sigma X_t \, dW_t) - \frac{1}{2 X_t^2} \sigma^2 X_t^2 \, dt \\
&= \left(\mu - \tfrac{\sigma^2}{2}\right) dt + \sigma \, dW_t.
\end{align*}
Integrando, $\ln X_t = \ln X_0 + (\mu - \tfrac{\sigma^2}{2}) t + \sigma W_t$, de modo que
$$ X_t = X_0 \exp\!\left(\sigma W_t + \left(\mu - \tfrac{\sigma^2}{2}\right) t\right). $$
Este es el modelo clásico de precio de un activo en finanzas (modelo de Black-Scholes).
\end{ejemplo}

\newpage

\subsection{Procesos estocásticos generales}

\begin{definicion}{Proceso de Itô}{proceso_ito}
Un \emph{proceso de Itô} es un proceso de la forma
$$ X_t = X_0 + \int_0^t f(X_s, s) \, ds + \int_0^t g(X_s, s) \, dW_s, $$
escrito en forma diferencial como
$$ dX_t = f(X_t, t) \, dt + g(X_t, t) \, dW_t. $$
$f$ se llama el \emph{drift} y $g$ el \emph{coeficiente de difusión}.
\end{definicion}

\begin{teorema}{Itô generalizado}{ito_general}
Si $X_t$ es proceso de Itô como arriba y $h \in C^{1,2}(\mathbb{R}_{\ge 0} \times \mathbb{R})$,
$M_t = h(t, X_t)$ cumple
$$ dM_t = \left(\frac{\partial h}{\partial t} + f \frac{\partial h}{\partial x} + \frac{g^2}{2}
\frac{\partial^2 h}{\partial x^2}\right) dt + g \frac{\partial h}{\partial x} \, dW_t. $$
\end{teorema}

\begin{proof}
Misma estrategia que el lema de Itô. Se particiona $[0, T]$ con $0 = t_0 < t_1 < \cdots < t_m =
T$ y se aplica Taylor de orden 2 a $h$:
$$ h(t_k, X_{t_k}) - h(t_{k-1}, X_{t_{k-1}}) = h_t \Delta t_k + h_x \Delta X_k + \tfrac{1}{2}
[h_{tt} (\Delta t_k)^2 + 2 h_{tx} \Delta t_k \Delta X_k + h_{xx} (\Delta X_k)^2] + \cdots, $$
con derivadas evaluadas en $(t_{k-1}, X_{t_{k-1}})$. Como $\Delta X_k = f \Delta t_k + g \Delta W_k
+ o(\sqrt{\Delta t_k})$,
$$ (\Delta X_k)^2 = g^2 (\Delta W_k)^2 + 2 f g \Delta t_k \Delta W_k + f^2 (\Delta t_k)^2 +
\cdots. $$
Las reglas asintóticas de variación cuadrática dan
$$ \sum (\Delta t_k)^2 \to 0, \qquad \sum \Delta t_k \Delta W_k \to 0, \qquad \sum (\Delta W_k)^2
\to \int dt, $$
de modo que solo sobrevive el término $h_{xx} g^2 (\Delta W)^2 / 2 \to \tfrac{g^2}{2} h_{xx} \, dt$.
Sumando los términos de orden 1, $h_t \, dt + h_x (f \, dt + g \, dW) = (h_t + f h_x) dt + g h_x
\, dW$. Sumando con el término cuadrático de $h_{xx}$ que sobrevive, queda la fórmula
enunciada.
\end{proof}

\newpage

\section{Teorema de Girsanov}

¿Qué pasa con la distribución de un proceso si se cambia la medida de probabilidad?
Girsanov dice exactamente cómo se transforma un browniano bajo un cambio de medida
absolutamente continua.

\subsection{Motivación: cambio de media de una normal}

Sea $Z \sim \mathcal{N}(0, 1)$ bajo $P$. Véase que existe una medida $Q$, equivalente a $P$ (es
decir, $P \ll Q \ll P$), tal que $Z \sim \mathcal{N}(\mu, 1)$ bajo $Q$. Para esto, se computa
primero
$$ E_P(e^{\mu Z}) = \frac{1}{\sqrt{2\pi}} \int_{-\infty}^\infty e^{\mu t} e^{-t^2/2} \, dt =
\frac{1}{\sqrt{2\pi}} \int e^{-(t-\mu)^2/2} dt \cdot e^{\mu^2/2} = e^{\mu^2/2}. $$
Por lo tanto $E_P(e^{\mu Z - \mu^2/2}) = 1$.

\textbf{Cambio de medida.} Para $A \in \mathcal{F}$ se define
$$ Q(A) = \int_A e^{\mu Z - \mu^2/2} \, dP. $$
$Q$ es una medida de probabilidad (positiva, $\sigma$-aditiva, $Q(\Omega) = 1$). Bajo $Q$,
se calcula la distribución de $Z$:
\begin{align*}
F_Q(a) = Q(\{Z \le a\}) &= \int_{\{Z \le a\}} e^{\mu Z - \mu^2/2} \, dP \\
&= \frac{1}{\sqrt{2\pi}} \int_{-\infty}^\infty \mathbf{1}_{\{t \le a\}} e^{\mu t - \mu^2/2}
e^{-t^2/2} dt \\
&= \frac{1}{\sqrt{2\pi}} \int_{-\infty}^a e^{-(t - \mu)^2/2} dt,
\end{align*}
o sea $Z \sim \mathcal{N}(\mu, 1)$ bajo $Q$.

\subsection{Versión continua}

Antes del enunciado, obsérvese por qué el candidato a densidad es una martingala. Sea
$$ Y_t = \int_0^t X_s \, dW_s - \frac{1}{2} \int_0^t X_s^2 \, ds, \qquad M_t = e^{Y_t}, $$
es decir $dY_t = X_t \, dW_t - \tfrac{1}{2} X_t^2 \, dt$, con $(dY_t)^2 = X_t^2 \, dt$ por las
reglas de variación cuadrática. Por el lema de Itô con $g(s) = e^s$,
$$ dM_t = e^{Y_t} \, dY_t + \tfrac{1}{2} e^{Y_t} (dY_t)^2 = e^{Y_t} \left( X_t \, dW_t -
\tfrac{1}{2} X_t^2 \, dt + \tfrac{1}{2} X_t^2 \, dt \right) = M_t X_t \, dW_t. $$
El término de drift se cancela exactamente — para eso está el $-\tfrac{1}{2} \int_0^t X_s^2
\, ds$ en la definición — y queda $M_t = 1 + \int_0^t M_s X_s \, dW_s$: una integral de Itô,
que es martingala de media $E(M_t) = 1$ siempre que el integrando sea suficientemente
integrable. La condición de Novikov del enunciado garantiza exactamente eso.

\begin{teorema}{Girsanov}{girsanov}
Sea $(W_t)_{t \in [0, T]}$ movimiento browniano bajo $P$ y sea $(X_t)$ proceso adaptado
suficientemente integrable. Se define
$$ M_t = \exp\!\left( \int_0^t X_s \, dW_s - \frac{1}{2} \int_0^t X_s^2 \, ds \right). $$
Bajo condiciones razonables (\emph{condición de Novikov}: $E[e^{\tfrac{1}{2} \int_0^T X_s^2 ds}]
< \infty$), $M_t$ es martingala y $E(M_T) = 1$. Se define la medida
$$ Q(A) = \int_A M_T \, dP \qquad (A \in \mathcal{F}_T). $$
Entonces $Q$ es equivalente a $P$, y el proceso
$$ \widetilde{W}_t = W_t - \int_0^t X_s \, ds $$
es movimiento browniano bajo $Q$.
\end{teorema}

\obs En el caso particular $X_t \equiv \mu/\sigma$ (constante), tomando $W_t' = \sigma W_t$ con
drift $\mu$, Girsanov dice que $\widetilde{W}_t = W_t - \mu t/\sigma$ es browniano bajo
$Q$ con $dQ/dP = \exp(\tfrac{\mu}{\sigma} W_T - \tfrac{\mu^2}{2\sigma^2} T)$. Esto justifica el
cambio del modelo geométrico de Black-Scholes desde la medida del mundo real a la medida
neutral al riesgo, sobre la que se apoya la valuación de derivados.

\obs Girsanov es la versión continua de Radon-Nikodym aplicada a procesos. La derivada
$dQ/dP = M_T$ se llama \emph{exponencial de Doléans-Dade} de $\int X \, dW$.

\newpage

\section{Aplicación: el modelo de Black-Scholes}

Las herramientas construidas hasta aquí (browniano, Itô, Girsanov) tienen una aplicación
emblemática: la fórmula de Black-Scholes para valuar opciones financieras. Es un caso donde
todo el aparato se pone a trabajar.

\subsection{Modelo del activo}

Considérese un \emph{activo subyacente} con precio $S_t$ que evoluciona según un browniano
geométrico:
$$ dS_t = \mu S_t \, dt + \sigma S_t \, dW_t, \qquad S_0 \text{ dado}, $$
donde $\mu \in \mathbb{R}$ es la \emph{tasa de retorno esperado} y $\sigma > 0$ la
\emph{volatilidad}. Como se vio al resolver la EDE con Itô,
$$ S_t = S_0 \exp\!\left(\sigma W_t + \left(\mu - \tfrac{\sigma^2}{2}\right) t\right). $$
Al lado existe un \emph{activo libre de riesgo} (bono) que crece a tasa constante $r$:
$$ B_t = e^{rt}, \qquad dB_t = r B_t \, dt. $$

\textbf{Derivado.} Un \emph{derivado} es un contrato que paga una función $\Phi(S_T)$ en un
tiempo futuro $T$. Por ejemplo, una \emph{opción call} europea con \emph{strike} $K$ paga
$$ \Phi(S_T) = (S_T - K)^+ = \max(S_T - K, 0). $$
El problema: ¿cuál es el precio justo $C_0$ a pagar hoy por este contrato?

\subsection{Argumento de no arbitraje y medida neutral al riesgo}

\textbf{Estrategia.} Si se puede construir, con $C_0$ dólares hoy, una cartera que en $T$
replique exactamente el pago $\Phi(S_T)$, entonces $C_0$ es el único precio compatible con la
\emph{ausencia de arbitraje} (sin posibilidad de ganancia segura sin inversión).

\textbf{Reducir al activo descontado.} Sea $\widetilde{S}_t = S_t / B_t = e^{-rt} S_t$, el precio
del activo medido en bonos. Por Itô,
$$ d\widetilde{S}_t = (\mu - r) \widetilde{S}_t \, dt + \sigma \widetilde{S}_t \, dW_t. $$
Si se pudiera hacer desaparecer el drift $(\mu - r) \widetilde{S}_t \, dt$, $\widetilde{S}$
sería
martingala. Esto es exactamente lo que da Girsanov: tomando
$$ X_t \equiv \frac{\mu - r}{\sigma}, $$
la medida $Q$ con $dQ/dP = \exp(-\tfrac{\mu - r}{\sigma} W_T - \tfrac{1}{2}(\tfrac{\mu-r}{\sigma})^2
T)$
hace que
$$ \widetilde{W}_t = W_t + \tfrac{\mu - r}{\sigma} t $$
sea browniano bajo $Q$. Reescribiendo,
$$ d\widetilde{S}_t = \sigma \widetilde{S}_t \, d\widetilde{W}_t \qquad \text{(bajo $Q$)}, $$
o sea $\widetilde{S}$ es martingala bajo $Q$.

\begin{definicion}{Medida neutral al riesgo}{Q_neutral}
La \emph{medida neutral al riesgo} es la única medida $Q$ equivalente a $P$ tal que
$\widetilde{S}_t = e^{-rt} S_t$ es martingala bajo $Q$. Bajo $Q$, el precio del activo evoluciona
con drift $r$ (no $\mu$):
$$ dS_t = r S_t \, dt + \sigma S_t \, d\widetilde{W}_t, \qquad S_t = S_0 \exp\!\left(\sigma
\widetilde{W}_t + (r - \tfrac{\sigma^2}{2}) t\right). $$
\end{definicion}

\obs El nombre ``neutral al riesgo'' es porque bajo $Q$ \emph{todo activo} (tanto $S$ como $B$)
crece, en promedio, a la misma tasa $r$: no hay premio por riesgo. La actitud de los inversores
hacia el riesgo, que afecta el drift $\mu$ del mundo real, desaparece en $Q$.

\subsection{Fórmula de precio}

\begin{teorema}{Precio Black-Scholes (versión general)}{bs_general}
Bajo no arbitraje y completitud del mercado, el precio justo en $t = 0$ de un derivado con pago
$\Phi(S_T)$ es
$$ C_0 = e^{-rT} E_Q(\Phi(S_T)). $$
\end{teorema}

\begin{proof}
\textit{Cartera autofinanciada.} Una cartera se describe por procesos $(\alpha_t, \beta_t)$
adaptados que indican cuántas unidades del activo y del bono se tienen en cada momento. Su valor
es
$$ V_t = \alpha_t S_t + \beta_t B_t. $$
La cartera es \emph{autofinanciada} si su variación solo viene de las variaciones del precio:
$$ dV_t = \alpha_t \, dS_t + \beta_t \, dB_t. $$
En términos del valor descontado $\widetilde{V}_t = V_t / B_t$, por Itô aplicado a
$V_t \cdot e^{-rt}$,
$$ d\widetilde{V}_t = \alpha_t \, d\widetilde{S}_t. $$

\textit{Replicación y martingala bajo $Q$.} Supóngase que existe una cartera autofinanciada que
\emph{replica} el pago: $V_T = \Phi(S_T)$. Bajo $Q$, $\widetilde{S}_t$ es martingala (por
Girsanov), y entonces
$$ \widetilde{V}_t - \widetilde{V}_0 = \int_0^t \alpha_s \, d\widetilde{S}_s $$
es una integral de Itô respecto de una martingala, por lo tanto $\widetilde{V}_t$ es martingala
bajo $Q$ (suponiendo la integrabilidad necesaria). Tomando esperanza,
$$ \widetilde{V}_0 = E_Q(\widetilde{V}_T) = E_Q(e^{-rT} \Phi(S_T)). $$
Y $V_0 = \widetilde{V}_0$ porque $B_0 = 1$, lo que da $C_0 = V_0 = e^{-rT} E_Q(\Phi(S_T))$.

\textit{Existencia de la cartera replicante.} En el modelo de Black-Scholes (un activo, un bono),
el mercado es \emph{completo}: por un teorema de representación de martingalas (la martingala
descontada $\widetilde{C}_t = E_Q(e^{-rT} \Phi(S_T) \mid \widetilde{\mathcal{F}}_t)$ se escribe como
integral de Itô respecto de $\widetilde{W}$, lo que da $\alpha_t$), siempre existe tal cartera.
El argumento de \emph{no arbitraje} implica que su precio inicial es único.
\end{proof}

\subsection{Fórmula explícita para el call europeo}

Tómese $\Phi(s) = (s - K)^+$. Bajo $Q$, $S_T = S_0 \exp(\sigma \widetilde{W}_T + (r -
\sigma^2/2) T)$, con $\widetilde{W}_T \sim \mathcal{N}(0, T)$, es decir
$\ln S_T \sim \mathcal{N}(\ln S_0 + (r - \sigma^2/2) T, \ \sigma^2 T)$.

Se definen
$$ d_1 = \frac{\ln(S_0/K) + (r + \sigma^2/2) T}{\sigma \sqrt{T}}, \qquad
d_2 = d_1 - \sigma \sqrt{T}, $$
y $N(\cdot)$ la función de distribución de una normal estándar.

\begin{teorema}{Fórmula de Black-Scholes}{bs_formula}
El precio de una opción call europea con strike $K$ y vencimiento $T$ es
$$ C_0 = S_0 \, N(d_1) - K e^{-rT} \, N(d_2). $$
\end{teorema}

\begin{proof}[Demostración (cálculo directo)]
Se calcula $E_Q((S_T - K)^+)$. Bajo $Q$ se escribe
$$ S_T = S_0 \, e^{\sigma \sqrt{T} Z + (r - \sigma^2/2) T}, \qquad Z \sim \mathcal{N}(0, 1). $$
La condición $S_T > K$ equivale a
$$ Z > \frac{\ln(K/S_0) - (r - \sigma^2/2) T}{\sigma \sqrt{T}} = -d_2. $$
Por lo tanto,
\begin{align*}
E_Q(S_T \mathbf{1}_{S_T > K}) &= S_0 e^{(r - \sigma^2/2) T} \int_{-d_2}^\infty e^{\sigma \sqrt{T} z}
\frac{e^{-z^2/2}}{\sqrt{2\pi}} dz \\
&= S_0 e^{rT} \int_{-d_2}^\infty \frac{e^{-(z - \sigma\sqrt{T})^2/2}}{\sqrt{2\pi}} dz
\qquad \text{(completando cuadrados)} \\
&= S_0 e^{rT} \cdot P(Z' > -d_2 - \sigma\sqrt{T}) \quad \text{con } Z' \sim \mathcal{N}(0,1) \\
&= S_0 e^{rT} \, N(d_2 + \sigma\sqrt{T}) = S_0 e^{rT} \, N(d_1).
\end{align*}
Y $E_Q(K \mathbf{1}_{S_T > K}) = K \cdot P(Z > -d_2) = K \, N(d_2)$. Sustituyendo en
$C_0 = e^{-rT} E_Q((S_T - K)^+) = e^{-rT}[E_Q(S_T \mathbf{1}_{S_T > K}) - K E_Q(\mathbf{1}_{S_T >
K})]$,
$$ C_0 = e^{-rT}[S_0 e^{rT} N(d_1) - K N(d_2)] = S_0 N(d_1) - K e^{-rT} N(d_2). $$
\end{proof}

\subsection{Interpretación financiera}

\begin{itemize}
    \item $N(d_2)$ es la probabilidad bajo $Q$ de que la opción \emph{termine in-the-money}, es
    decir $S_T > K$. El segundo término $K e^{-rT} N(d_2)$ es el valor descontado de pagar $K$
    en ese caso.
    \item $S_0 N(d_1)$ es la \emph{cobertura delta}: la cantidad de subyacente (en valor presente)
    que hay que mantener en cartera para replicar el pago de la call.
    \item La fórmula \emph{no depende} de $\mu$, el drift real del activo. Solo depende de
    $\sigma$, $r$, $S_0$, $K$, $T$. Esto es una sorpresa histórica: el precio justo no requiere
    estimar cómo crece la acción, solo cuánto fluctúa.
\end{itemize}

\obs El modelo tiene supuestos fuertes (volatilidad constante, sin saltos, sin costos de
transacción, posibilidad de continuous trading). En la práctica se observan desviaciones
sistemáticas (volatility smile, jumps) que motivan modelos más sofisticados (Heston, modelos con
saltos, volatilidad estocástica). Pero Black-Scholes sigue siendo la referencia base.

\end{document}
